% =====================================================================
%  Harald Høffding, SPINOZAS LIV OG LÆRE.  Et Bidrag til Tænkningens Historie
%  i det 17de Aarhundrede.  København: Rudolph Klein, 1877 (trykt hos Nielsen &
%  Lydiche).  VIII, 166 pp.
%
%  TRANSCRIPTION -- GENERATED; DO NOT EDIT BY HAND.  Made in two steps, both rerunnable:
%    python3 ebook-method/spinozas-liv-og-laere/draft.py --emit
%      this template + the scan's text layer (Google's OCR, in HathiTrust's PDF of the
%      Harvard copy, Phil 3827.3; ~/bibliotek/Høffding, Harald/1877-spinozas.pdf, see
%      SCANS.tsv), then patches.py (corrections read at the image that a word decision
%      cannot express: word order, notes run together, italics, letterspacing, pencil
%      marks in the Harvard copy read as text);
%    python3 ebook-method/collate/check.py spinozas-liv-og-laere --apply
%      the decisions in ebook-method/spinozas-liv-og-laere/decisions.tsv, taken against an
%      independent tesseract reading of the images (rules.py writes the ones settled by
%      rule, each marked with its rule; the rest were read at the image).
%  No e-book of this book exists.  State (2026-10-03): every disagreement between the
%  two OCRs settled (OPEN 0); italics found by stroke slant and letterspacing by the gaps
%  between letters, both confirmed at the image (short spaced words may be missed); the
%  front matter read at the image.  Misprints that BOTH OCRs normalised can remain: no
%  full image pass was made.  The print's errata (Rettelser, p. 166) are noted at their
%  sites, the text printed as is.
%
%  CONVENTIONS.  "% --- p. N ---" + \opage{N} at the first word of printed page N
%  (a word broken over the turn is split with %).  Footnotes as printed, marked
%  *), **), ***) and numbered per page (\footnote[k]).  Italics \textit; letterspacing
%  \emph.  A note running over a page is set whole at its mark.  Running heads, folios and
%  signature marks are not transcribed.  Print governs; misprints are kept and
%  marked PRINTED AS IS.
% =====================================================================
\documentclass[11pt]{article}
\usepackage[utf8]{inputenc}
\usepackage[T1]{fontenc}
\usepackage{libertinus}
\usepackage{libertinust1math}
\usepackage[danish]{babel}
\usepackage[protrusion=true,expansion=true]{microtype}
\usepackage{setspace}
\usepackage[margin=1.3in]{geometry}
\usepackage{fancyhdr}
\usepackage{hyperref}
\hypersetup{pdftitle={Spinozas Liv og Lære}, pdfauthor={Harald Høffding}, hidelinks}
\pagestyle{fancy}
\fancyhf{}
\fancyhead[LE,RO]{\thepage}
\fancyhead[C]{\textit{Harald Høffding: Spinozas Liv og Lære (1877)}}
\setstretch{1.3}
\usepackage{marginnote}
\usepackage{graphicx}
\DeclareUnicodeCharacter{0254}{\reflectbox{c}}   % ɔ (id est)
\renewcommand{\thefootnote}{%
  \ifcase\value{footnote}\or *)\or **)\or ***)\else \arabic{footnote})\fi}
\newcommand{\tocline}[3]{\noindent\makebox[2.2em][r]{#1}\hspace{0.6em}#2\dotfill\ #3\par}
\newcommand{\opage}[1]{\marginnote{\footnotesize\textit{[#1]}}}

\begin{document}

% --- front matter: title (PDF 8), read at the image.  The library's shelfmarks
% (Phil 3827.3, Walker fund) on the verso are not transcribed.
\begin{center}
{\LARGE SPINOZAS LIV OG LÆRE.}\\[1em]
% (an ornamental rule)
{\small ET BIDRAG TIL TÆNKNINGENS HISTORIE I DET 17DE AARHUNDREDE.}\\[2.5em]
AF\\[1em]
Dr. H. HØFFDING.\\[2em]
% (a vignette: an engraved portrait of Spinoza in a medallion)
\vspace{12em}
KØBENHAVN.\\[0.3em]
FORLAGT AF RUDOLPH KLEIN.\\[0.3em]
1877.
\end{center}
\clearpage
% --- verso (PDF 9) ---
\begin{center}\vspace*{\fill}{\small TRYKT HOS NIELSEN \& LYDICHE.}\end{center}
\clearpage

% ===== FORORD (pp. III--V) =====
% --- p. III ---
\setcounter{footnote}{0}%
\opage{III}\begin{center}\emph{Forord.}\end{center}


Dette lille Skrifts Affattelse skyldes den ærede
Udgivers Opfordring til mig om at indlemme Spinoza
i Rækken af hans „kulturhistoriske Personligheder“.
Jeg har ikke været blind for Vanskelighederne ved en
saadan Opgave. Spinozas kulturhistoriske Betydning er
nøje knyttet til hans Plads i den abstrakte Spekulations
Historie og kan ej ret forstaas uden nærmere Indblik i
hans Filosofi. Jeg har stræbt at give en saa let fattelig
Oversigt over denne som muligt og har anlagt Skriftet
saaledes, at Læseren i de tre første Afsnit stifter et
foreløbigt Bekendtskab med Spinozas vigtigste Ideer, før
der i sidste Afsnit gives en sammenhængende Fremstilling
af hans Lære.

Vi have i vor Litteratur en grundig og skarpsindig
Monografi over Spinoza af afdøde Professor Brøchner.
Ved den klare Belysning, hvori dette Arbejde stiller de
forskellige Elementer i Spinozas Filosofi, indtager det
en høj Plads i den store Spinozalitteratur. Men det
er ikke let tilgængeligt for den, der uden videre gaaende
filosofiske Forkundskaber vil stifte Bekendtskab
med Spinoza. Jeg har tænkt mit lille Skrift som en
% --- p. IV ---
\setcounter{footnote}{0}%
\opage{IV}Slags Indledning til Brøchners, i det jeg har givet
Afkald paa Fuldstændigheden og Gennemførelsen i det
enkelte, saa vel som paa den lærde Dokumentation.
Dog har jeg kunnet benytte, hvad der endnu ej forelaa
fuldstændig, da Brøchner skrev sin Afhandling, det
interessante Skrift „Om Gud og Mennesket“, som viser
os den første Form for Spinozas Filosofi og tydelig
henpeger paa hans historiske Forudsætninger, lige som
ogsaa de Breve fra og til Spinoza, man har fundet, og
som oplyse flere Punkter i hans Biografi. Alt her hen
hørende har van Vloten samlet i sin Bog „Ad Benedicti
de Spinoza opera qvæ supersunt omnia supplementum“
(Amstelodami 1862).

Det Synspunkt, hvorfra jeg betragter og bedømmer
Spinozas Filosofi, afviger fra Brøchners. Min afdøde
Lærer, om hvem jeg saa ofte er bleven mindet under
Udarbejdelsen af dette Skrift, stod som spekulativ Filosof
Spinoza nærmere, end jeg gør. Hans Kritik er
derfor heller ikke rettet mod selve Principet i Spinozas
Spekulation, men mod Inkonsekventserne i dets Gennemførelse.
Jeg maatte gaa videre, ikke blot i Medfør af
mine egne filosofiske Anskuelser, men ogsaa fordi jeg
satte mig den Opgave at angive Betydningen af Spinozas
Teorier, naar man tager Hensyn til Videnskabernes
senere Historie og særlig til vor Tids Forsken.
Men trods denne videre gaaende Kritik er min Beundring
for den store Tænker snarere tiltagen end
aftagen. Spinoza kan siges at have fejret en dobbelt
Opstandelse, efter at han i mere end et Aarhundrede
havde ligget hen „som en død Hund“: første Gang i
Slutningen af det 18de Aarhundrede ved den filosofiske,
% --- p. V ---
\setcounter{footnote}{0}%
\opage{V}poetiske og religiøse Opvækkelse, anden Gang i vore
Dage ved de store naturvidenskabelige Hypoteser og
Filosofiens Genopvaagnen efter den Slappelse, den spekulative
Overmættelse havde fremkaldt. Det er karakteristisk
for hans Stilling i Tænkningens Historie, at to
saa forskellige Epoker kunne mødes i Tilslutning og
Sympati over for ham. I mange Maader er hans Metode
og Tankegang endog rent modsat den, som hersker
i vor Tids Forsken; men man har her et Exempel paa,
hvorledes --- for at bruge et Billede, jeg har anvendt
et Steds i \emph{dette} Skrift --- Arbejdere, der grave paa en
Tunnel hver fra sin Side, mødes i Midten, naar de kun,
hver fra sit Udgangspunkt, gaa frem ad den rette Vej.

\medskip
\noindent\hspace*{3em}Juni 1877.\par
\hfill\emph{Harald Høffding.}\par

\clearpage
% --- p. VII --- INDHOLD (PDF 14), read at the image.  Its section numbers are
% 1.-5., where the text heads the sections I.-V.; it gives section 3 as p. 102,
% where p. 102 ends section II and III begins on p. 103; it calls the fourth part
% "Spinozas Lære", headed "Spinozas Filosofi." on p. 72.  All as printed.
\opage{VII}
\begin{center}{\large\emph{Indhold.}}\end{center}
\medskip
\hfill Side.\par
\tocline{}{Indledning}{1}
\tocline{}{Spinozas Levned}{19}
\tocline{}{Spinoza og hans Korrespondenter}{44}
\tocline{}{Spinozas Lære}{72}
\tocline{}{\quad 1. Grundbegreberne}{72}
\tocline{}{\quad 2. Aand og Materie}{89}
\tocline{}{\quad 3. Erkendelseslære}{102}
\tocline{}{\quad 4. Psykologi og Etik}{120}
\tocline{}{\quad 5. Statslære. Religionsfilosofi}{147}
\clearpage

% ===== TEXT (pp. 1--164), LIST OF TERMS (p. 165), VERSE AND ERRATA (p. 166) =====
% --- p. 1 ---
\setcounter{footnote}{0}%
\opage{1}\begin{center}Indledning.\end{center}


Den store Betydning, Videnskabens Historie har,
bliver i vore Dage mere og mere anerkendt. Ikke
blot gaar der ved Siden af Interessen for selve de
videnskabelige Problemer en stærk historisk Interesse,
et Ønske om at kende disse Problemers tidligere Stilling
og den Maade, hvorpaa de behandledes tidligere;
men med Hensyn til mange Spørgsmaal gælder det, at
de først ret forstaas, naar man ser, hvorledes de historisk
ere komne frem og blevne opfattede. Ofte er den
Vej, ad hvilken en videnskabelig Sandhed efter Haanden
er bleven vunden, den samme, som i forkortet Maalestok
maa tilbagelægges af enhver, der vil sætte sig ind i den.
Men der er sikkert ingen Videnskab, for hvilken Historien
har saa stor Betydning som for Filosofien. Dette
maa for det første udledes deraf, at den er rigere paa
Bestræbelser end paa Resultater. Ved en given Sum af
Resultater kan man til en vis Grad glemme, naar, hvorledes
og af hvem de ere vundne. Den pytagoræiske Læresætning
kan fuldt tilegnes, om man saa aldrig har hørt
noget om Pytagoras. Men hvor Hovedvægten maa
lægges paa den stedse fornyede Stræben, --- hvor det
% --- p. 2 ---
\setcounter{footnote}{0}%
\opage{2}netop gælder at skærpe Problemernes Braad for paa
ny at sætte Tænkningen i Bevægelse, der er den historiske
Orientering af en særegen Betydning. De tidligere
Forsøg blive dels Forbilleder, dels Advarsler for de
senere. Men det ligger dernæst i de filosofiske Ideers
Natur, at de tit komme for tidlig til Verden. De ere
ofte Idealer, der opstilles, vundne ved en dristig Spekulation,
som gaar ud over den reale Erkendelses faktiske
Resultater og forudgriber et Maal, der i Virkeligheden
kun kan naas møjsommelig og efter lange Tider. Da
Filosofiens sidste Opgave er at udvikle en sammenhængende
og omfattende Verdensanskuelse, ere de filosofiske
Ideer rettede mod Enheden og Helheden i Tingene,
men forudsætte derhos de enkelte Tings Natur og
Vexelforhold bekendte. Det gaar i øvrigt paa samme
Maade ogsaa inden for de reale Videnskaber med de
Anskuelser og Hypoteser, der udtrykke Fænomenernes
Sammenhæng og give deres Forklaring. Huyghens opstillede
i det 17de Aarhundrede den saakaldte Bølgeteori
til Forklaring af Lysfænomenerne, og først i vort
Aarhundrede har den vundet almindelig Anerkendelse.
Lamarcks Teori om Arternes Oprindelse er først ved
Darwins Undersøgelser bleven Genstand for større Opmærksomhed
og Tilslutning. Betydningen af de videnskabelige
Ideer, i Filosofien som i andre Videnskaber,
beror paa, hvor vidt de virkelig ere grebne ud af Naturen
og gennemførte med Konsekvents og Klarhed.
Have de Rod i Virkeligheden og indre Konsekvents,
vil deres Sejr komme, tidlig eller sent.

Spinozas Filosofi er i begge Henseender af stor
Interesse. Det er ikke tilfældigt, at Opfordringen til
% --- p. 3 ---
\setcounter{footnote}{0}%
\opage{3}at mindes 200 -Aarsdagen efter hans Død ved Oprejsning
af hans Billedstøtte rundt i Evropa har fundet
Tilslutning ikke blot hos Filosofer af Faget, men ogsaa hos
fremragende Repræsentanter for de reale Videnskaber.
Ti hans System har ikke blot klassisk Betydning ved
sin Renhed og klare Sammenhæng, saa at vi der, som
i et Mønster, kunne lære, hvorledes en videnskabelig
Verdensanskuelse maa bygges. Dette vilde ikke være
nok til at vække en almindelig Interesse. Netop den
systematiske Form er det ved en filosofisk Lære, der
snarest hjemfalder til Forgængelighed. Vi kunne bevæge
os i Spinozismen, som i de andre filosofiske Systemer,
paa lignende Maade som i Oldtidens Bygningsværker,
hvilke vi beundre, og af hvilke vi kunne lære
meget, skønt vi ikke selv nu vilde bygge paa samme
Maade. Men man kan ikke i vore Dage studere Filosofi
og almindelig Naturvidenskab, uden at Tanken
paa mange Maader føres tilbage til Spinoza. Mange
af hans Grundtanker kunne siges at have vundet deres
Bekræftelse ved den nyere Videnskabs Resultater. Af
Systemets forkrænkelige Form udvikler Ideernes Kærne
sig, og det har da sin store Interesse at se, hvorledes,
ved hvilken Metode og gennem hvilken Motivering
Spinoza kom til de Ideer, vi nu, til Dels ad andre
Veje, ere naaede til.

Men Spinozas Ideer vare for ham mere end Genstande
for teoretisk Spekulation. De smeltede sammen
med hans personlige Livsanskuelse, og netop den Alvor,
hvormed han hævdede den videnskabelige Forskens Ret
og Evne til at bestemme Livsanskuelsen, blev afgørende
for hans Skæbne og førte ham til Strid med den kirke%
% --- p. 4 ---
\setcounter{footnote}{0}%
\opage{4}lige Avtoritet, der her kun anerkender Traditionens Ret.
Derved faar ogsaa hans personlige Liv og hans Karakter
en Betydning i filosofisk Henseende. Filosofien er vel
nærmest et teoretisk Anliggende; men den kan og skal
ikke opgive den Fordring, at den erkendte Sandhed
bliver bestemmende for Livet; den skal virke som et
Salt over for Dogmatismens Ferskhed. Det er derfor
dens Stolthed, at den i sin Historie har Skikkelser som
Spinozas at pege hen paa, og vi skulle i det Følgende
se, hvorledes hans Karakter og Biografi har spillet en
ikke ringe Rolle i Kampen for den frie Forskens Berettigelse.


Som allerede antydet maa et filosofisk System bedømmes
fra en dobbelt Side. Man maa undersøge dets
indre Sammenhæng og Konsekvents og dets Overensstemmelse
med Virkeligheden, saaledes som denne er
os tilgængelig ved videnskabelig Erfaring. Intet af
disse to Synspunkter kan undværes. Hvad Tanken
søger i Filosofien, er en videnskabelig Livs- og Verdensopfattelse,
hvor alle enkelte givne Kendsgerninger ere forbundne
til en Helhed og bragte i indre Sammenhæng med
hverandre. Der skal altsaa være Enhed og Sammenhæng
til Stede; ikke blot logiske Modsigelser, men
ogsaa det løsrevne, isolerede og rent empiriske er
ufilosofisk. Men denne Enhed og Sammenhæng maa
ikke købes ved Tilsidesættelse af væsentlige Fakta, ej
heller ved Forvanskning og vilkaarlig Omtydning af
saadanne. Det Princip, der ligger til Grund for et
filosofisk System, maa derfor vælges saaledes, at det i
sin videre Udfoldelse kan rumme og belyse alt, hvad
der staar fast ved videnskabelig Erfaring. Til forskellige
% --- p. 5 ---
\setcounter{footnote}{0}%
\opage{5}Tider har man i Filosofien lagt forskellig Vægt paa
disse to Sider. Den dogmatiske og idealistiske
Retning i Filosofien lægger Hovedvægten paa Principets
Gennemførelse; den kritiske og realistiske lægger derimod
Hovedvægten paa Overensstemmelsen med Erfaringen
og paa selve Principets Begrundelse. Den sidst
nævnte Retning synes mere og mere at faa Overhaand
og er i hvert Tilfælde i vore Dage den herskende.

Filosofiens Historie frembyder for os en Række
Systemer, som hvert lægge sit Princip til Grund.
Hvilket Princip der gribes i det enkelte System, er
naturligvis ikke vilkaarligt. Det vil for en dybere Betragtning\footnote[1]{Se H. Brøchner: Bidrag til Opfattelsen af Filosofiens historiske Udvikling. Kbhvn. 1869.}
vise sig, at et filosofisk Systems Hovedidé
fremgaar af hele Aandsudviklingen, lige som den atter
virker tilbage paa denne. Den enkelte Tænker finder
saa at sige sit Princip i den aandelige Atmosfære, hvori
han har udviklet sig, og udskiller det af denne, ligesom
en Kemiker udskiller et Stof fra dets Forbindelse med
flere andre Stoffer for at betragte det for sig selv.
Mangfoldige Principer, der ere væsentlige Elementer i
den almindelige Bevidstheds Indhold, have saaledes
faaet deres nærmere Behandling og Belysning i forskellige
filosofiske Systemer. Men Filosofien maa nødvendigvis
komme til den Indsigt, at det er af største
Vigtighed at undersøge, hvor vidt selve Principerne
kunne begrundes, og hvorledes det er muligt at naa
et sandt og almengyldigt Princip. Denne Vej slog den
engelske Filosofi strax fra Begyndelsen af ind paa, og
% --- p. 6 ---
\setcounter{footnote}{0}%
\opage{6}Kant, den nyere Tids største Filosof, har med stor
Dybsindighed paavist den som den eneste sikre Vej til
filosofisk Erkendelse. Den Filosofi, der gaar ud fra et
Princip, som ikke selv begrundes, vil altid komme til
at staa Teologien nær; der vil være et aandeligt Slægtskab,
hvor meget vedkommende Filosof end hævder den
rationelle Erkendelses Selvstændighed, og hvor energisk
han end har emanciperet sig fra Avtoritetstroen. Det
er navnlig i Metoden, at dette Slægtskab mellem den
dogmatiske eller spekulative Filosofi og Teologien viser
sig. Begge gaa deduktivt frem, søge at udlede deres
Lære af ét højeste Princip, som for Teologien er givet
ved Avtoritet og Tradition, medens det i den spekulative
Filosofi er en Idé, der med umiddelbar Klarhed
fremstiller sig for Fornuften og synes at være selvindlysende.
Den kritiske Filosofi er saa at sige Filosofi
i anden Potents. Den vil ikke blot filosofere over
et givet Princip, men vil forstaa dets Udspring og dets
Berettigelse. Den fordrer derhos Regnskab ikke blot
for Avtoritetens Princip, men ogsaa for den „selvindlysende“
Idé. Den gaar induktivt og analytisk til
Værks, forlanger at faa fast Grund under Fødderne
og anerkender kun de Principer, der med Nødvendighed
fremgaa af Erfaringskendsgerningernes Analyse.

Spinoza tilhører afgjort den dogmatisk -s pekulative
Retning i Filosofien. Hans Metode er den deduktive.
Hans Grundbegreb, Substantsens Idé, er opgaaet for
ham dels ved Paavirkning af den Folkereligions Lære,
inden for hvilken han var opvoxet, dels ved den forudgaaende
Filosofis Indflydelse. Det stod for ham med
en indre Sandhed, hævet over Nødvendigheden af
% --- p. 7 ---
\setcounter{footnote}{0}%
\opage{7}egentligt Bevis. Substantsen (hollandsk: zelfstandigheid)
er det i og for sig værende, som ikke forudsætter
noget som helst andet. Substantsens Idé ligger til Grund
for de højere Folkereligioners Forestilling om Guddommen
som det absolute Væsen, Ophav til alt, evigt,
uforgængeligt og uden Skranker. Men den religiøse,
af Fantasien beherskede Forestilling kan ikke gennemføre
denne Idé med fuld Konsekvents; den stiller stedse
uvilkaarlig den endelige Tilværelse uden for det absolute
Væsen, ja tillægger den endog Evne til at hindre eller
dog hæmme dets Virken. I en dunkel Fornemmelse
af den Modsigelse, hvori den herved hilder sig, griber
den religiøse Forestilling saa til Skabelsesdogmet som
Gaadens Løsning: vel er der noget uden for det absolute
Væsen, med selvstændig Realitet, men det er blevet til
ved det absolute Væsens egen Vilje og har kun
Realitet ved den. Spinoza optager Substantsens Idé
saaledes, at han fører den ud over den teologiske Dualisme
(Tohedslære). Men Sammenhængen med den
religiøse Forestillingskreds ses ikke blot af, at han i
sit Hovedskrift (Etiken) ligefrem betegner Substantsen
som Gud, men især af hele Tonen og Udtryksmaaden
i et ældre Skrift af ham („Om Gud, Mennesket og
dettes Lyksalighed“), der for 15 Aar siden blev fundet
i Amsterdam. Han benytter her hele Tiden religiøse
Forestillinger, men giver dem en rationalistisk Anvendelse
og Gennemførelse. Hvor nær han her endnu
staar det gamle Testamentes Forestillingskreds, ses af
en Udtalelse som denne: „Den sande Kærlighed udspringer
stedse af den Erkendelse, at dens Genstand
er herlig og god. Kan det da flamme heftigere over for
% --- p. 8 ---
\setcounter{footnote}{0}%
\opage{8}nogen end over for Herren vor Gud? Ti han alene er
herlig og er det fuldkomne Gode“. Spinoza er en
religiøs Tænker. Gudsidéen staar for ham med indre
Klarhed og Gyldighed, og hele hans Tænkning bevæger
sig inden for den. Men ved den videre Udvikling, han
giver Gudsidéen, træder han ud over den positive Religions
Skranker og paadrager sig derved Beskyldningen
for Irreligiøsitet og Ateisme. Han saa klart, at det i
og for sig værende Væsen maa være det eneste Væsen
og indeslutte alt i sig, saa at der ikke kan gives nogen
sand Væren uden for det. Han ophævede altsaa den
dualistiske Modsætning mellem Gud og Verden og forkastede
Skabelsesdogmet. Derved er han den egentlige
Grundlægger af Monismen eller Enhedslæren. Filosofien,
der søger Enheden og Sammenhængen i alt,
maa nødvendigvis være Monisme, og for saa vidt havde
Hegel Ret, naar han sagde: „Spinoza er et saadant
Hovedpunkt i den moderne Filosofi, at man i Sandhed
kan sige: enten har du Spinozismen, eller du har ingen
Filosofi.“

Spinoza er naturligvis ikke den første, der har opstillet
Monismen. Inden for den græske Filosofi har han
Forgængere i Eleaterne, Platon og Aristoteles; men
deres Idéer have næppe haft nogen Indflydelse paa
ham. Anderledes synes det, efter hvad man kan spore
af den nys nævnte Afhandling, at forholde sig med de
jødiske og arabiske Religionsfilosofer for Middelalderen, % PRINTED AS IS: „for“ (Rettelser: S. 8, L. 5 f. n., l. fra)
hos hvem der fremtræder en stærk panteistisk Tendents.
I det 11te Aarhundrede levede der i Saragossa en jødisk
Digter ved Navn Ibn-Gebirol, hvis religiøse Hymner
vandt stor Berømmelse. Han var den egentlige Genføder
% --- p. 9 ---
\setcounter{footnote}{0}%
\opage{9}af den hebraiske Poesi og indtager den første Plads
mellem de jødiske Digtere i Middelalderen. Men den
lærde Kender af Middelalderens jødiske og arabiske
Litteratur, Franskmanden S. Munk, har tillige godtgjort
(se hans „Mélanges de philosophie juive et arabe“),
at Ibn-Gebirol er den samme som den Forfatter, de
kristne Skolastikere polemiserede saa ivrig imod under
det fordrejede Navn Avicebron, og som kan kaldes
Middelalderens klassiske Panteist, lige som Spinoza den
nyere Tids. I hans vigtigste Skrift „Livets Kilde“
\textit{(fons vitæ)} udviklede han --- støttende sig til Aristoteles
og Nyplatonikerne ---
Idéen om én Form og én Materie
som aabenbarende sig i alt, baade i det aandelige
og i det legemlige. I Væsenernes uendelige Rækkefølge
er der intet Spring. Der er én Substants, som
bestaar ved sig selv og bærer alle Forskelle, som antager
alle Former og giver alt sit Væsen og sit Navn.
Ved Siden af disse panteistiske Idéer gaar der vel et
Forsøg paa at fastholde Skabelsesdogmet; men „Livets
Kilde“ er dog et rent filosofisk Skrift, det eneste jødiske
Skrift fra Middelalderen, hvori der ingen Skriftsteder
eller Traditioner citeres. Der er, i Følge S. Munk, en
indre Forbindelse mellem Ibn-Gebirols Poesi og hans
Filosofi: „det er den samme Kreds af Tanker, der
aabenbarer sig under to forskellige Former, og naar
man i hans Digte stedse sér en Aand, der er hensunken
i dybe Grublerier og fremmed for hvad der interesserer
Mængden, saa bære hans filosofiske Værker ofte Spor
af Digterens levende Indbildningskraft og Begejstring.“
En arabisk Filosof fra det følgende Aarhundrede, Ibn
Roschd (Skolastikernes Avèrroës, født i Cordova, død % PRINTED AS IS: „Avèrroës“ (an accent on the first e; possibly a speck)
% --- p. 10 ---
\setcounter{footnote}{0}%
\opage{10}1198), udtaler, som det synes uden Kendskab til Ibn-Gebirols
Skrifter, lignende Ideer. Han polemiserer
ligefrem mod Skabelsen af intet. Den evige Materie
indeholder alle Formerne i sig. Medens baade muhamedanske,
jødiske og kristne Teologer ophæve Naturens
Begreb, hævder Filosofien en stadig Overgang
fra Mulighed til Virkelighed efter den evige Nødvendigheds
Lov. Mennesket tænker ved Delagtighed i den
absolute Fornuft, som er evig, medens ethvert enkelt
tænkende Individ er forgængeligt. Den sande Erkendelse,
som naas ved at blive delagtig i den guddommelige
Fornuft, er den højeste Fuldkommenhed; er den
naaet, spørges der ikke længer om Løn og Straf.
--- En noget yngre samtidig af Ibn Roschd var den jødiske
Teolog og Filosof Moses ben Maimon (Maimonides),
født i Cordova 1138, død i Cairo 1204. Han søgte at
forene Tro og Viden, i det han fortolkede Bibelen paa
rationalistisk Vis, opfattede de religiøse Ceremonier
som Hjælpemidler til Indøvelse i den sande Erkendelse
og hævdede videnskabelig Indsigt i Guds Væsen som
nødvendig Forudsætning for Kærligheden til ham.
Hans store og mærkelige Skrift „De vildfarendes Lærer“
(Moreh Nebuchim), som er oversat paa Fransk af S.
Munk, har haft stor Indflydelse paa mange fremragende
Aander inden for Jødedommen; det var Læsningen af
det, der førte Mænd som Spinoza, Mendelsohn og Maimon
fra Synagogen over til Filosofien. Det er vistnok
gennem det, at Spinoza har lært Ibn Roschd's Ideer
at kende. --- Efter hvad Grätz anfører i sin „Geschichte
der Juden“ (10de Bd.) skal der finde en vis Overensstemmelse
Sted mellem nogle af Spinozas Ideer og
% --- p. 11 ---
\setcounter{footnote}{0}%
\opage{11}Udtryk og den saa kaldte Kabbala, en inden for Jødedommen
opstaaet, fantastisk religionsfilosofisk Mystik.
Grätz mener særlig at kunne paavise en Lighed mellem
et i Aaret 1656 udgivet kabbalistisk Værk af Abraham
de Herrera og flere Sætninger i første Bog af Spinozas
„Etik“.

Da vi vide af Spinozas Biografi, at han tilbragte
sine første Ungdomsaar med at studere jødisk Teologi
og Filosofi, har det sin Interesse at se, hvilke vækkende
Paavirkninger i Retning af hans senere Standpunkt han
her direkte eller indirekte kan have modtaget. Men
disse Paavirkninger have sikkert været rent foreløbige;
selv udtalte han, at det var Descartes's Skrifter, der
havde lært ham Filosofien. Der er en stor Forskel
mellem hine middelalderlige Spekulationer og Spinozas
klare og ædruelige Tænkning, der har gjort Matematiken
til sit Forbillede.

Af den oven for anførte ældre Afhandling ser man
ogsaa, at Spinoza har været paavirket af Renaissancetidens
største Filosof Giordano Bruno. Denne urolige
og gærende Aand viser selv tilbage til Ibn Gebirol og
Ibn Roschd, saa vel som til Nyplatonikerne og Aristoteles.
Men afset fra disse ældre Ideer, han søgte at
genføde, er Giordano Brunos hele Verdensanskuelse
væsentlig bestemt og motiveret ved Kopernikus's Opdagelse
af Jordens Bevægelse omkring Solen. Han var
denne Opdagelses begejstrede Profet rundt om i Evropa
og blev dens Martyr. Ved denne Opdagelse udslettedes
den gamle Modsætning mellem Himmel og Jord, og der
aabnede sig en uendelig Horizont, i hvilken Brunos
Fantasi gerne fordybede sig. Hans Naturbegejstring
% --- p. 12 ---
\setcounter{footnote}{0}%
\opage{12}kan endnu spores i enkelte af Spinozas Udtalelser.
Hvad der adskiller Spinoza fra Bruno, er ikke blot
den strenge, systematiske Tænkning, men navnlig det
store Fremskridt, som paa Naturvidenskabens og Filosofiens
Omraade gjordes i Mellemrummet mellem deres
Optræden. Brunos Naturopfattelse er poetisk og idealistisk
og beror ikke paa en Erkendelse af de bestemte
Naturlove og af Fænomenernes virkelige Sammenhæng.
Men omtrent paa samme Tid, som han led Martyrdøden
paa Inkvisitionens Baal, gjorde Galilæi sine store Opdagelser
og lagde Grunden til den moderne exakte Naturvidenskab
ved at anvende Mekaniken paa Astronomien.
Bevægelsens mekaniske Love, der især træde
tydelig frem i Himmellegemernes Bevægelser, bleve nu
Grundlaget for Naturopfattelsen og Naturforklaringen.
Galilæi er Stifteren af den matematiske Fysik, der
hviler paa det Princip, at alle fysiske Fænomener
kunne forklares som Bevægelser, som man kan maale
enten direkte eller indirekte. Ofte er det kun ved
Experimentets Hjælp, at man bliver i Stand til at
maale; men saa snart dette er lykkedes, er Udgangspunktet
vundet for den matematiske Beregning. Det
Maal, hvortil den af Galilæi grundlagte matematiske
Naturvidenskab stræber, er at udlede de i Erfaringen
givne Kendsgerninger med matematisk Nødvendighed
af saa faa og saa simple Forudsætninger som muligt.

Men den experimentale Forskens Vej er lang; den
opnaar i lange Tider kun spredte Resultater og aabner
kun fjerne Udsigter til Totalanskuelser og omfattende
Synspunkter. For mange Filosofer har der derfor
været noget utilfredsstillende ved denne Metode. Kun
% --- p. 13 ---
\setcounter{footnote}{0}%
\opage{13}den engelske Filosofi sluttede sig strax --- med Baco,
Hobbes og Locke --- til den experimentale Forskens
Princip og Metode. Men den kunde kun gøre det ved
foreløbig at stille de store spekulative Problemer i Bero.
Hvor der var en levende Trang til Drøftelse af de
Spørgsmaal, der ligge paa Grænsen mellem Religionen
og Filosofien, kunde denne Metode ikke endnu gøre
Fyldest. Den engelske Empirisme (Erfaringsfilosofi)
har i lang Tid meget vel kunnet blomstre med de
sædvanlige teologiske Forestillinger som Baggrund.
Anderledes paa Fastlandet, hvor Filosofien paa Grund
af sin spekulative Tendents strax stødte sammen med
Teologien og maatte erklære sig for eller imod. Den
experimentale Forsken fik her kun underordnet Betydning;
den væsentlige Interesse er rettet mod de almindelige
Ideers Udvikling og Gennemførelse. Dette fremtræder
klart hos Spinozas nærmeste Forgænger, Franskmanden
Réné Descartes.

Descartes var utilfredsstillet ved den skolastiske
Filosofi, han havde lært i Jesuitkollegiet, og søgte en
ny, som havde en fastere Grundvold. Han besluttede
da at se bort fra alle sine tidligere Forudsætninger,
fra alt Erkendelsesindhold, han tidligere havde optaget,
at tvivle om alt, ikke for at blive Skeptiker, men for
at lære den menneskelige Erkendelses Væsen at kende
og se, hvor vidt Fornuften kan naa ved sine egne
Kræfter. Da nu selve Tvivlen er en Tænken, er Tænkningen
selv det eneste, om hvis Virkelighed der ikke
kan tvivles, og alt, hvad jeg indser lige saa klart og
tydelig, som at jeg tænker, maa være sandt. Derfor
opstiller Descartes den Sandhedsregel, som gjorde saa
% --- p. 14 ---
\setcounter{footnote}{0}%
\opage{14}stærkt Indtryk paa Spinoza, at klar og tydelig Begriben
er Betingelsen for, at vi kunne antage noget for sandt.
Den førte Brug, han gør af denne Regel, er at vise,
at selve Begrebet Tvivl fører os til Begrebet om et
absolut fuldkomment Væsen. Ti, siger han, hvorledes
kan jeg erkende, at jeg tvivler, og at jeg attraar en
sand Erkendelse, d. v. s. at der fattes mig noget, og
at jeg ikke er ganske fuldkommen, naar jeg ikke har
Ideen om et fuldkomnere Væsen end mit eget, ved
Sammenligning med hvilket jeg kan se min Naturs
Mangler? Gud eller det absolute Væsen, der i sig
indeslutter al Fuldkommenhed, bliver derfor det højeste
Princip for Filosofien, og den sande Erkendelse bestaar
i ved klare og tydelige Slutninger at udlede alt af
Guds Væsen. „Da Gud“, siger Descartes, „alene er
den sande Aarsag til alt, hvad der er eller kan være,
er det aabenbart, at vi ville følge den bedste Maade
at filosofere paa, naar vi af Guds Væsen søge at udlede
Forklaringen af de skabte Ting.“ Saaledes udleder han
de højeste Naturlove af Guds Uforanderlighed, som er
en af hans Fuldkommenheder, og af de højeste Naturlove
eller Bevægelseslovene de specielle Love. Galilæi's
Fremgangsmaade er Descartes for empirisk; han
bebrejder Galilæi, at han ikke gaar tilbage til de sidste
Grunde. Den første af de Naturlove, han udleder af
Guds Uforanderlighed, er netop den af Galilæi experimentalt
paaviste Inertiens Lov: enhver Ting forbliver,
saa vidt det beror paa den, i samme Tilstand og kan
ikke nogen Sinde forandres uden af ydre Aarsager.

Descartes vil ikke, som Galilæi, gaa fra Virkningerne
til Aarsagerne, men omvendt, fra Aarsagerne til
% --- p. 15 ---
\setcounter{footnote}{0}%
\opage{15}Virkningerne, og den højeste Aarsag er umiddelbart
given for ham i Tankens eget Væsen. Hans Metode
er derfor aldeles deduktiv. Derimod faa Erfaringskendsgerningerne
kun den Betydning at lede Tanken i
en vis bestemt Retning. Det Forbillede, Descartes her
er paavirket af, er Matematiken. Den tænkte han allerede
paa, da han opstillede den klare og tydelige Erkendelse
som Særkende for Sandheden. Den matematiske
Viden staar som et stort Exempel paa gennemsigtig
Klarhed, i det alt er ført tilbage til Størrelsesforhold, der
finde deres nøjagtige Bestemmelse. Descartes er enig
med Galilæi i, at Naturvidenskaben skal udvikles matematisk;
men han mener af de almindelige matematiske
Love at kunne forklare alle specielle Fænomener. Han
gør altsaa ingen Forskel mellem Fysik og Matematik.
„Hele min Fysik“, siger han i et af sine Breve, „er
ikke andet end Matematik.“ Alt hvad der ikke havde
matematisk Klarhed, var for ham dunkelt og forvirret.
Materiens Væsen saa han i Udstrækningen. Alle andre
Egenskaber ved Materien kan man nemlig efter hans
Mening tænke borte, kun ikke Udstrækningen i Længde,
Brede og Højde; af den, som det væsentlige, maa derfor
alle andre Egenskaber afledes. Grundegenskaberne
ved Materien ere Delelighed og Bevægelighed, og de
forskellige materielle Fænomener opstaa alle ved Delenes
Bevægelse.

Det laa saa meget des nærmere for Descartes at
gøre det matematiske til Hovedsagen i sin Verdensopfattelse,
som han selv var en Matematiker af første
Rang; han er den første, der har anvendt Algebra paa
Geometrien, og altsaa Stifteren af den analytiske Geo
% --- p. 16 ---
\setcounter{footnote}{0}%
\opage{16}metri og Mekanik. Og det maatte i det hele ligge
nær for den Tænker, der vilde opstille et nyt System
og en ny Metode i Steden for den middelalderlige
Skolastik med dens Ræsonnementsspidsfindigheder og
Ordforklaringer, at gribe den matematiske Metode, som
er al Forskens Ideal og Mønster. Der laa tillige heri
mere end en Form; der laa et helt Program, en profetisk
Bebudelse. Verden staar ikke længer som et Spil
af dunkle Kræfter, men som en uendelig Række Fænomener,
der indbyrdes staa i bestemte Forhold, for hvilke
klare Love gælde. Descartes opstiller altsaa den Fordring,
at alt skal forklares efter nødvendige Love, en
Fordring, hvis fremskridende Opfyldelse den nyere Naturvidenskab
er. Hvad Galilæi paaviste paa et enkelt
Omraade, gjorde Descartes til Princip for hele Verdensopfattelsen.


At Descartes med saa forbavsende Hurtighed og
Lethed foretager Overgangen fra sit Udgangspunkt,
Tvivlen om alt, til Ideen om det absolute Væsen og
gør denne Idé til sin Læres Princip, kan kun forklares
ved den religiøse Tros Indflydelse, og vi have her et
klassisk Exempel paa den oven for fremhævede Ejendommelighed
ved den spekulative Filosofis Principer.
Men hvad her særlig maa fremhæves, fordi Spinoza
netop paa dette Punkt træder i Modsætning til Descartes,
er, at denne optager Ideen om det absolute
Væsen i samme Form, hvori den fremtræder i den
positive Religion. Han var selv troende Katolik, og
saa energisk han end fastholder Naturlovenes Gyldighed
og uindskrænkede Herredømme, saa fører han dem dog
i sidste Instants tilbage til Guds almægtige Vilje og
% --- p. 17 ---
\setcounter{footnote}{0}%
\opage{17}forklarer dem ved Guds Uforanderlighed. \emph{Hvis} Gud
vilde, kunde han ophæve Naturlovene, ja endogsaa de
matematiske Love; 2 og 3 ere kun 5, fordi Gud vil
det. Og paa et andet Punkt bliver den guddommelige
Indgriben endog ligefrem nødvendig for Descartes's System.
Naar Materiens Væsen ene bestaar i Udstrækningen,
hvor skal saa den materielle Bevægelses Aarsag
søges? Da den ikke ligger i Materiens Væsen, som er
rent geometrisk, maa den ligge i Guds Vilje. Gud har,
siger Descartes, ved Skabelsen sat Materiens Dele i en
vis Bevægelse.

Da Descartes endnu fastholder Skabelsesdogmet,
antager han andre Substantser foruden den absolute.
Gud har frembragt Aanden, den tænkende Substants,
og Materien, den udstrakte Substants. Men her staa
vi da netop ved den Modsigelse, som den filosoferende
Tænkning opdager i den religiøse Forestilling, at
vi paa den ene Side have en absolut Substants, og saa
dog paa den anden Side endelige Substantser, frembragte
af den. Descartes bemærker rigtig nok, at han ikke
tager Ordet Substants i samme Betydning om Gud og
om de endelige Væsener; men han viser ikke, hvorledes
han --- naar han éngang har brugt det om Gud --- i
det hele taget kan faa Ret til at bruge det om endelige
Væsener. I denne Inkonsekvents ligger det Problem
skjult, som var det dybeste Motiv i Spinozas Tænkning.
Naar Descartes betegner baade Aand og Materie
som Substantser, kommer han derved ikke blot til at
stille dem i Modsætning til den absolute Substants,
men ogsaa til hinanden indbyrdes. Han faar ikke blot
en Dualisme mellem det uendelige og det endelige
% --- p. 18 ---
\setcounter{footnote}{0}%
\opage{18}(Gud og Verden), men ogsaa mellem Aand og Materie.
Med Hensyn til begge Spørgsmaal er han paavirket af
de teologiske Forestillinger. Han havde klart indset
Forskellen mellem det aandelige (Bevidsthedsfænomenerne)
og det materielle og Umuligheden af at udlede
det ene af det andet. Men denne Forskel blev for
ham til en fuldstændig Modsætning, saa at han og hans
Skole til sidst appellerede til et Mirakel for at forklare
Sjælens Forbindelse med Legemet. Ogsaa her griber
Spinoza ind med klar Konsekvents og fører Problemet
et mægtigt Skridt videre.

Efter at vi nu i Korthed have fremhævet de historiske
Betingelser, under hvilke Spinozas Tænkning udviklede
sig, gaa vi over til en Skildring af hans Liv.

% --- p. 19 ---
\setcounter{footnote}{0}%
\opage{19}\begin{center}Spinozas Levned.\end{center}


Man vil allerede af det foregaaende have set, at
Spinozas Lære stod i en stærk Modsætning til de paa
hans Tid ene herskende religiøse Forestillinger, og naar
det nu erindres, hvor mægtigt det politiske og kirkelige
Avtoritetsherredømme var i det 17de Aarhundrede,
staar det som et Vidunder, at han fik Lov til at leve
saa rolig og uforstyrret, som han gjorde. Der var
sikkert heller ikke noget andet Land end netop Holland,
hvor dette havde været muligt. Holland er Religionsfrihedens
Vugge. I Aaret 1609 havde Hollænderne
endt deres heltemodige Frihedskamp mod det mægtige
Spanien, --- som en Historieskriver har sagt, det eneste
moderne Sidestykke til Grækernes Frihedskampe. Som
Grækerne reddede den vesterlandske Kultur, saaledes
reddede Hollænderne den religiøse og politiske Frihed.
I det Manifest, hvori de 1581 sagde sig løs fra det
spanske Herredømme, beraabte de sig paa „Naturens
Lov“, --- et Varsel om Nordamerikanernes Uafhængighedserklæring
og om Menneskerettighedernes Erklæring
under den franske Revolution. Frihed paa ét Omraade
maa i Længden føre Frihed paa alle andre med sig.
% --- p. 20 ---
\setcounter{footnote}{0}%
\opage{20}Protestantismen var i sin første Tid ikke tolérant, og
den havde ogsaa Vanskelighed ved at være det, da det
gjaldt en Kamp paa Liv og Død med Romerkirken.
Der maatte holdes stræng Mandstugt inden for de
kæmpendes Rækker. Vilhelm den første af Oranien
synes at have opfattet Religionskampen fra et mere
universelt Synspunkt, som Kamp for almindelig religiøs
Frihed. Men hans Efterfølgere lode sig navnlig
lede af deres Higen efter at grunde et Monarki og
sluttede sig derfor til den ortodoxe Gejstlighed, hvis
Indflydelse paa Almuen derved ogsaa kom dem til
gode. Over for Hoffet, Hæren, Gejstligheden og Almuen
stod det republikanske Parti, der især støttede
sig til den velhavende Borgerstand. Det sluttede sig i
religiøs Henseende til de saa kaldte Arminianere, som
bestred den ortodoxe Kalvinisme, hævdede Samvittighedens
Frihed og gav enhver Lov til at fortolke Bibelen,
som han selv vilde. Medens alle andre Kirkesamfund
paa den Tid omgærdede sig med et forsvarligt
Apparat af „symbolske Bøger“, kunde enhver være
Medlem af Arminianernes Menighed, der var fri for
Afgudsdyrkelse, forargeligt Levned og Intolerants og
anerkendte Skriften som Regel for sit Liv, hvorledes
han saa end fortolkede den. Det maa fremhæves til
deres Ære, at de hævdede Tolerancen, før de selv
blev forfulgte. Ti Forfølgelsen udeblev ikke. Den
reformerte Kirkeforsamling i Dortrecht fordømte deres
Lære og udstødte dem af Kirkesamfundet. Moritz af
Oranien støttede de ortodoxe, som holdt sig til ham,
fordi Republikanerne og Arminianerne vilde indskrænke
de gejstliges Indflydelse. Først nogen Tid efter Moritz's
% --- p. 21 ---
\setcounter{footnote}{0}%
\opage{21}Død fik Arminianerne Lov til at holde Gudsdyrkelse.
Af dem udsprang igen flere mindre Sekter, som fulgte
friere Former i Tro og Liv.

Det var intet Under, at de, der i andre Lande
forfulgtes for deres Tro, kastede deres Blik paa Holland
og valgte det til deres Tilflugtssted. Medens
Hollænderne kæmpede for Friheden, ja hele det 17de
Aarhundrede igennem, brændte Inkvisitionens Baal
lystig i Spanien og Portugal. Man forfulgte særlig de
Jøder, der af Tvang havde ladet sig døbe, de saa
kaldte Marannos, af hvilke mange i Smug fastholdt
deres gamle Tro og deres gamle Skikke. Allerede
under den hollandske Uafhængighedskrig vare mange
Marannos flygtede til Holland. 1598 byggedes den
første Synagoge i Amsterdam, og der kom snart to til.
I Begyndelsen bleve de vel mistænkte for at være
spanske Spioner; men da man havde overbevist sig om
det ugrundede heri, og da de tillige medbragte store
Kapitaler og ved deres Handelsaand gav Amsterdam
et Stød fremad, viste Øvrigheden dem stor Imødekommen.
Den jødiske Menighed voxede stærkt. I
Aaret 1639 oprettedes en Højskole for jødisk Sprog
og Teologi, hvor Spinoza gjorde sine første Studier.

Baruch Spinoza er født den 24de November 1632.
Hovedkilden til hans Biografi, Præsten Colerus, som
dog ikke er noget paalideligt Vidne med Hensyn til
Spinozas Barndom og Ungdom, fortæller, at han fødtes
i Amsterdam. Men Grätz, den lærde Forfatter af
„Geschichte der Juden“, har af et Sted i et af Spinozas
Breve sluttet, at han maa være bleven født i Spanien
og have tilbragt sin første Barndom der. Spinoza
% --- p. 22 ---
\setcounter{footnote}{0}%
\opage{22}siger nemlig, at han har kendt en Jøde, som kaldtes
Judas den trofaste, der med stort Heltemod led Martyrdøden
paa Baalet. Grätz oplyser nu, at denne
Judas den trofaste var en spansk Adelsmand, der ved
ivrigt Studium af det gamle Testamente lededes til
Overbevisningen om Jødedommens Sandhed og derfor
blev brændt af Inkvisitionen i Aaret 1644. Hvor kan
nu Spinoza have kendt ham, hvis han ikke som Barn
har levet i Spanien? --- Spinozas Forældre kaldes


portugisiske Jøder; men dette var det almindelige Navn
for alle Jøder fra den pyrenæiske Halvø. Enten før
eller, ifølge Grätz's Hypotese, efter Sønnens Fødsel have
de, som saa mange andre Marannos taget deres Tilflugt
til Holland. Faderen synes at have været en
velhavende Handelsmand. Paa Grund af Drengens
gode Evner blev han bestemt til den lærde Stand og
besøgte den jødiske Højskole, hvor han fordybede sig
i Studiet af det gamle Testamente og Talmud. Ogsaa
de jødiske Religionsfilosofer fra Middelalderen og den
kabbalistiske Mystik synes han at have stiftet Bekendtskab
med; om en af hans Lærere, Manasse ben Israel,
vides det, at han hyldede de kabbalistiske Lærdomme.
Hans skarpe Forstand fandt snart Modsigelser i den
teologiske Lære. Ofte satte han sine Lærere i Forlegenhed
ved sine Indvendinger og Spørgsmaal. En af
hans Biografer, Lægen Lucas, fortæller, at da Rygtet
om hans Tvivl udbredte sig, kom to af hans Meddisciple
til ham og trængte ind paa ham med Spørgsmaal
angaaende forskellige Punkter af den ortodoxe
Lære, som Spinoza forgæves søgte at undgaa ved at
sige, at de jo havde Moses og Profeterne. Senere
% --- p. 23 ---
\setcounter{footnote}{0}%
\opage{23}optraadte de samme to Personer som Vidner mod
ham og beskyldte ham for at have haanet Moses og
Jødefolket. Dette gjorde saa meget des større Opsigt
i den jødiske Menighed, som man troede i den begavede
unge Mand at se en fremtidig Støtte for Synagogen
og Højskolen. Han havde været sin Talmudlærer
Morteiras Yndling, og kun ugerne bestemte denne
sig til strenge Skridt imod ham, tilmed da Spinoza
var meget beskeden og tilbageholden og førte et
stille og roligt Liv. Da Spinoza mærkede, at man
mistænkte ham, trak han sig endnu mere tilbage, og
hans Besøg i Synagogen bleve sjeldnere. Han begyndte
nu at omgaas kristne og følte Trang til en universellere
Dannelse end den bibelsk -t eologiske. Han lærte
Latin, som jo den Gang var det fælles videnskabelige
Tungemaal, og fik siden til Lærer i de klassiske Sprog
en Læge ved Navn van Ende, en bekendt Humanist,
der i sit Hus havde et Slags Undervisningsanstalt for
unge Mennesker. Colerus fortæller, at Spinoza blev
forelsket i hans Datter, som hjalp sin Fader med
Undervisningen, men at hun foretrak en af hans Meddisciple,
der bedre forstod sig paa Galanteriet. Men
denne Fortælling, hvorpaa Auerbach har bygget en
Roman, kan næppe være sand, da det er blevet oplyst,
at van Endes Datter paa den Tid, da Spinoza blev
exkommuniceret, kun var tolv Aar gammel. Van Ende
var meget fri i sine religiøse Anskuelser og lærte, som
Colerus siger, sine Disciple andet end Latin\footnote[1]{Van Ende bosatte sig siden i Paris og blev der hængt som hollandsk Spion og Medvider i Ridderen af Rohan's Sammensværgelse.}. Spinozas
% --- p. 24 ---
\setcounter{footnote}{0}%
\opage{24}Tvivl og Betænkeligheder have her fundet deres Næring.
Snart opgav han ganske Teologien og kastede sig over
naturvidenskabelige og filosofiske Studier. „Descartes's
Skrifter faldt ham i Hænderne; han læste dem med
Begærlighed og har siden ofte erklæret, at han af dem
havde øst sine filosofiske Kundskaber. Han var henrykt
over den Grundsætning af Descartes, at man ikke
skal antage noget for sandt, som ikke er bevist med
gode og solide Grunde“ (Colerus). I Descartes's Lære
fandt han en Verdensanskuelse, der overalt var bygget
paa Tænkning og Fornuft. Bekendtskabet med den
har fuldendt Emancipationen fra den jødiske Teologi.
Hans Trosfæller nærede Mistanke om, at han havde i
Sinde at træde over til Kristendommen, og frygtede,
at han vilde blive en farlig Modstander af den jødiske
Religion. Det maatte, som Grätz bemærker, gøre
Rabbinerne lidet stemte til Tolerance, at samtidig en
Mængde Jøder lede Martyrdøden i Spanien. Spinoza
var desuden ikke det første Exempel paa Frafald fra
den fædrene Tro i Amsterdammermenigheden. En
Snes Aar i Forvejen havde Synagogen gentagne Gange
banlyst Uriel da Costa, en Ætling af Marannos, der
ved Læsning af det gamle Testamente førtes tilbage
til Jødedommen og flygtede med sin Familie til Amsterdam,
hvor han lod sig omskære, men som, skuffet
i sin Tro paa Jødedommens Idealitet, optraadte med
Heftighed mod „Farisæerne“ og ved fortsat Grublen
førtes til deistiske Anskuelser. Da Costas urolige og
uklare Sind førte ham ikke til noget fast Standpunkt.
Fortvivlelse over ved Banlysningen at være udelukket
fra al Omgang med Stammefrænder gjorde, at han
% --- p. 25 ---
\setcounter{footnote}{0}%
\opage{25}underkastede sig en haard og ydmygende Bod for igen
at optages i Synagogen. Men snart efter haanede han
paa ny Jødedommen, blev atter banlyst og dræbte til
sidst sig selv i Fortvivlelse.

Med Spinoza vilde man dog helst undgaa et
aabenbart Brud. Man tilbød ham en ikke ubetydelig
aarlig Pension, naar han vilde love ikke at træde
offentlig op imod Jødedommen og af og til at besøge
Synagogen. Da dette Tilbud ikke modtoges, gjorde
en Fanatiker et Snigmordsforsøg paa Spinoza; men
denne slap med nogle Rifter i sin Kappe. Til sidst
skred man (1656) til den højtidelige Banlysning, der
udelukkede ham fra alle Stammefrænders Omgang.
Lucas bemærker: „Banlysningen har en saadan Vægt
hos Jøderne, at den banlystes bedste Ven ikke tør
vove at gøre ham nogen Tjeneste eller blot at tale til
ham, for ikke at komme til at lide samme Straf. Derfor
ville de, der frygte Ensomhed, hellere lide hvilken
som helst anden Straf end Forbandelsen“. Spinoza
havde tidlig lært at sætte Pris paa Ensomheden, og
han var tillige ved sine Studier og sin hele Aandsudvikling
bleven indviet i en universellere Verden og
havde faaet et større aandeligt Fædreland. Da han erfarede,
at man vilde udstøde ham, skal han have sagt:
„Man tvinger mig her ikke til noget, som jeg ikke
vilde have gjort af mig selv. Da man nu vil det saa,
betræder jeg den Vej, der aabner sig for mig, og jeg
kan trøste mig med, at jeg gaar bort mere uskyldig
end de gamle Hebræere gik fra Ægypten, skønt mit
Underhold ikke er mere sikret end deres“. Men Rabbinerne
lode sig ikke nøje med Banlysningen. De
% --- p. 26 ---
\setcounter{footnote}{0}%
\opage{26}angav for Øvrigheden, at Spinoza havde spottet Gud
og Moses, og forlangte ham fjernet fra Amsterdam.
Maaske have de frygtet, at Spinoza skulde drage flere
med sig; allerede Banformlen omtaler, at han har forført
nogle, og der er ogsaa, som vi snart skulle høre,
andre Spor af, at han allerede paa denne Tid har haft
Disciple. Da den reformerte Gejstlighed sluttede sig
til de jødiske Kolleger, blev Spinoza en Tid forvist fra
Byen „til Opretholdelse af Orden og Subordination“.
Et spansk Forsvarsskrift, han skal have forfattet i
denne Anledning, er ikke blevet bevaret; formodentlig
har han deri udtalt de Anskuelser om Religionsfriheden,
han senere udførlig udviklede i sin „teologisk-politiske
Afhandling“. Han tog Ophold hos en Ven paa Landet
uden for Amsterdam. Det kom ham nu til gode, at
han i sin Tid, efter den Forskrift, Talmud giver alle
lærde, havde lært et Haandværk, nemlig at slibe optiske
Glas. Ved dette Arbejde skaffede han sig, hvad
han behøvede til Livets Ophold.

Kort efter Udstødelsen af Synagogen maa det
ældste Arbejde, vi have fra Spinozas Haand, være forfattet.
Det er Afhandlingen „Om Gud, Mennesket og
dettes Lyksalighed“, oprindelig skrevet paa Latin, som
for 15 Aar siden fandtes i en hollandsk Oversættelse.
Den har stor Interesse ved at vise os Spinozas Anskuelse
i sin Vorden. Han har korrigeret Kartesianismen
ved den Erkendelse, at naar Substantsbegrebet
opfattes konsekvent, kan der kun være én Substants,
hvis Aabenbaringsformer alle Ting ere. Ligeledes
fremtræder klart og bestemt hans ejendommelige deterministiske
Opfattelse; alt, ogsaa den menneskelige
% --- p. 27 ---
\setcounter{footnote}{0}%
\opage{27}Vilje, følger Nødvendighedens Lov. Men han har endnu
den kartesianske Opfattelse af Forholdet mellem Sjæl
og Legeme, som han lader paavirke hinanden indbyrdes.
Endelig anvender han hele Tiden religiøse
Forestillinger og Udtryk og antager Hensigtsaarsager
i Naturen. Dette sidste er af særlig Interesse, da det
afgjort viser hen til andre Forgængere end Descartes,

--- til Middelalderens og Renaissancens Tænkere\footnote[1]{Se ovenfor p. 8---11.}.
Hverken Descartes eller Spinozas senere Lære tilskriver
Hensigtsaarsagerne Gyldighed i Naturopfattelsen; begge
hævde paa ethvert Punkt den strenge mekaniske Forklaringsmaade.


I Fortalen til den hollandske Oversættelse af denne
Afhandling hedder det, at den i sin Tid var nedskreven
af Spinoza til Brug for hans Disciple, der vilde dyrke
Studiet af Etik og Filosofi. --- Afhandlingen selv
ender med disse Ord: „Jeg beder de Venner, for hvem
jeg har skrevet dette, ikke at undre eder over det
nye deri, da I vel vide, at noget ikke hører op at
være sandt, fordi det ikke antages af mange. Og da
Beskaffenheden af den Tidsalder, i hvilken vi leve,
ikke er eder ubekendt, beder jeg eder indstændig at
være meget forsigtige med at gøre andre bekendte med
disse Ting. Jeg vil ikke sige, I at skulle beholde dem
ganske for eder selv, men kun, at naar I nogen Sinde
begynde med at meddele andre dem, maa intet andet
Øjemed bevæge eder end eders Medmenneskers Vel“.

I Aaret 1661 flyttede Spinoza (der efter Udstødelsen
af Synagogen havde ombyttet sit hebraiske Navn
% --- p. 28 ---
\setcounter{footnote}{0}%
\opage{28}Baruch med det enstydige latinske Benedictus) til
Rhynsburg, en lille By ved Leyden. Her havde der
under den arminianske Forfølgelse dannet sig en lille
fri Menighed, som ikke vilde vide af Præster eller
Trosbekendelser, men kun ganske i Almindelighed forlangte
en Tro paa Gud og Kristus; de kaldtes Kollegianter
eller Rhynsburgere og bestod endnu i Begyndelsen
af dette Aarhundrede. Spinoza synes at
have staaet paa en venskabelig Fod med dem, saa vel
som med de saa kaldte Mennoniter (Gendøbere). Det
var i Kollegianternes forrige Vajsenhus, man for en
Del Aar siden fandt flere hidtil ubekendte Breve fra
og til Spinoza. Spinoza gik aldrig over til Kristendommen,
men synes at have sympatiseret navnlig med
de friere, halvt rationalistiske Former, hvori han fandt
den hos Arminianere og Kollegianter.

I sit landlige Opholdssted fortsatte han med Iver
sine Studier, og man ser af hans Breve, at hans System
i alle sine ejendommelige Grundtræk har været
færdigt i Begyndelsen af Opholdet i Rhynsburg. Den
Kreds af yngre Mænd, der allerede havde dannet sig
om ham i Amsterdam, og som udgjorde en Art filosofisk
Menighed, der stræbte at leve sig ind i den nye
Livsanskuelse, vedblev han at staa i Forbindelse med.
Mærkeligt nok har den første Udgiver af Spinozas
Breve i disse strøget enhver Hentydning til denne
Kreds af Tilhængere; vel sagtens fordi det dengang,
strax efter Spinozas Død, ikke var raadeligt at henlede
Opmærksomheden paa den. Af de i Kollegiantervajsenhuset
fundne Manuskripter faa vi derimod Underretning
om den. En af Spinozas yngre Venner, Simon
% --- p. 29 ---
\setcounter{footnote}{0}%
\opage{29}de Vries, der efter hans Raad studerede Kemi og
Medicin i Amsterdam, skriver til ham for at faa Oplysning
om nogle vanskelige Punkter i et Manuskript,
Spinoza havde overladt til ham, og som vistnok var
Udkastet til „Etiken“. Men medens han i den oprindelige
Udgave af Brevene gør denne Forespørgsel i
sit eget Navn, taler han i det fundne Originalmanuskript
paa fleres Vegne. „Jeg har“, siger han, „længe ønsket
at besøge dig, men min Tid og den strenge Vinter
have ikke tilladt det. Jeg beklager min Skæbne, at vi
nu ere skilte af saa stor en Afstand. Hvor lykkelig
er ikke din Husfælle, der kan tale med dig om de
højeste Ting under Maaltiderne og paa Spadsereture.
Men skønt vi ere legemlig fjernede fra hinanden, er
du dog stedse nærværende for min Aand, især naar
jeg sysler med dine Skrifter. Da nu ikke alt i disse
forekommer mine Kammerater klart nok --- hvorfor vi
igen ere begyndte paa vort Kollegium, --- og tillige
for at du ikke skal tro, at jeg har glemt dig, skriver
jeg dette Brev til dig. Vort Kollegium er indrettet
paa følgende Maade. En af os læser op (hver efter
sin Tur), forklarer alt, som han har forstaaet det, og
beviser det efter dine Sætningers Rækkefølge. Naar
vi nu ikke kunne blive enige om Forstaaelsen af et
eller andet Punkt, notere vi det og skrive til dig
derom, for at du kan forklare os det, og for at vi
under din Vejledning kunne forsvare Sandheden mod
de overtroiske fromme og kristne, og saa staa fast
over for den hele Verden.“

Det første Skrift af Spinoza, som blev trykt, var
en systematisk Fremstilling af Descartes's Filosofi, som
% --- p. 30 ---
\setcounter{footnote}{0}%
\opage{30}han havde givet en ung Mand, der vilde undervises af
ham, og hvem han endnu ikke vilde indvie i sin egen
Lære. Det blev udgivet 1663 af hans Ven, Lægen
Ludvig Meyer, under Titel: \textit{Renati des Cartes Principia
philosophiæ more geometrico demonstrata}. Meyer bemærker
i Fortalen, at Skriftet ikke maatte betragtes
som Udtryk for Forfatterens egne Anskuelser, og nævner
som Punkter, hvori disse adskille sig fra Descartes's
Spørgsmaalet om Viljens Frihed, om Menneskets Sjæl
som egen Substants og om Erkendelsens Grænser.
Hvad det sidste Punkt angaar, havde Descartes til
Fordel for Teologien erklæret, at der var meget, der
overgik den menneskelige Fatteevne. Spinoza derimod
mente, at de Problemer, Descartes var stanset ved,
nok kunde løses, naar man blot anvendte den rette
Fremgangsmaade. I et Tillæg, \textit{Cogitata metaphysica},
antyder han noget nærmere Overgangen fra Descartes's
Anskuelser til hans egne.

Foruden den Kreds af yngre Mænd i Amsterdam,
som Spinoza stadig stod i Forbindelse med, havde han
mange Korrespondenter og modtog mange Besøg. Hans
Tid har været stærkt optagen, i det han foruden Glasslibningen,
der var hans Erhvervskilde, fortsatte sine
Studeringer og syslede med kemiske og fysiske Forsøg.
Det var maaske for at være mere uforstyrret, at han
i Aaret 1664 forlod Rhynsburg\footnote[1]{Det lille Hus, hvori han boede i Rhynsburg, staar endnu, og Gaden kaldes undertiden Spinozagaden. Paa Huset findes følgende Indskrift fra Aaret 1667: Ach, waren alle menschen wijs, En wilden darby wel: De aard waar haar een Paradijs, Nu is ze meest een Hel. De tre første Linjer kunne stemme med Spinozas Tanker; den fjerde strider aldeles mod hans optimistiske Opfattelse.} og bosatte sig i
% --- p. 31 ---
\setcounter{footnote}{0}%
\opage{31}Voorburg, en Mil fra Haag. Men da han havde faaet
mange Venner i Haag, flyttede han i Aaret 1669 paa
deres Opfordring til denne By og levede der Resten af
sit Liv. Først boede han hos en Enke, siden hos
Maleren van der Spijk. Den meste Tid tilbragte han
paa sit Kammer med sine Arbejder. Han adspredte
sig med at fange Fluer og sætte dem i Edderkoppespind
eller med at lade Edderkopper kæmpe, et Syn,
hvorover han undertiden kunde briste i Skoggerlatter.
Ofte undersøgte han ogsaa Insekter eller enkelte Dele
af dem i et Mikroskop. Han var en dygtig Tegner,
og man fandt efter ham en hel Samling Portræter,
han havde udført; blandt andet et af ham selv i
Masaniello dragt\footnote[1]{Det Portræt, hvorefter Vignetten til dette Hefte er taget, er malet af Spinozas Vært van der Spijk. Der er stor Ulighed mellem de forskellige Portræter af Spinoza, som haves.}%
. Hyppig gik han ned til sin Værts
Familie og samtalede med den om alt muligt, især
naar der var Sygdom eller Sorg i Huset. Hans Levevis
var i høj Grad tarvelig og hans Fornødenheder
faa. Præsten Colerus anfører i sin Biografi flere Details
i denne Henseende efter de Regninger, der fandtes
efter hans Død. Hverken Nydelser eller Penge spillede
nogen Rolle for ham. Han viste stor Uegennyttighed
over for sine Søstre, skønt de havde villet udelukke ham
fra hans Fædrenearv. Efter at han ved en Proces
havde faaet sin Arveret anerkendt, overlod han dem
% --- p. 32 ---
\setcounter{footnote}{0}%
\opage{32}næsten hele Arven. Hans Ven Simon de Vries vilde
forære ham en Pengesum, for at han kunde leve mere
behagelig; men han afslog den, ligesom han ogsaa
modsatte sig, at Vries indsatte ham til sin Universalarving;
en Pension, Vries ved sin Død havde fastsat
for ham, modtog han først efter at have nedsat den
betydelig.

I Aaret 1670 udgav han sin \textit{Tractatus theologo-politicus},
hvori han undersøgte Religionens Forhold til
Filosofien og viste Religionsfrihedens Nødvendighed.
Dette Skrift er den første Proklamation af ubetinget
Trosfrihed. Under Kampene mellem Katoliker og
Protestanter i det 16de og 17de Aarhundrede havde
vistnok de overvundne ofte talt Tolerancens Sag, men
stedse med store Indskrænkninger. Spinoza hævder
Retten til at udtale sin Overbevisning som en Menneskerettighed,
Staten ej kan berøve Individet, naar dette
ikke kommer andres Ret for nær. Menneskenes aandelige
Natur er saa forskellig, at de ikke alle kunne føle
sig tilfredsstillede ved de samme Meninger. Spinoza
viser fremdeles, at der gives en af Dogmer og Ceremonier
uafhængig Religiøsitet, anlægger en historisk
Kritik paa de bibelske Bøger, hvorved han blev hele
den nyere Bibelkritiks Forgænger, og søger at give en
psykologisk og historisk Forklaring af Profeternes
aandelige Tilstand. Da han nu desuden viste, hvor
farligt det er, naar de gejstlige have for stor Indflydelse
i Staten, er det intet Under, at der rejste sig
et stort Røre over den anonyme Bog. Ved dristige
Udtalelser, der lode Guds Magt falde sammen med
Naturens, vakte han desuden Forargelse ogsaa hos
% --- p. 33 ---
\setcounter{footnote}{0}%
\opage{33}dem, der ikke yndede Hierarkiet. Flere af hans
Korrespondenter bleve fra denne Tid køligere. Bogen
blev forbudt, men udkom snart paa ny under andre
Navne.

Spinozas Stilling var ikke uden Fare, skønt han
søgte at leve saa ubemærket som muligt. Det mægtige
oraniske og klerikale Parti, der havde Almuen paa sin
Side, havde han paa en Maade udfordret ved sit Skrift.
Derhos var han en Ven af den berømte Statsmand
Jan de Witt, der stod i Spidsen for det republikanske
Parti, og mod hvem Præsterne i deres Prædikener
ophidsede Folket. Ulykkerne i den franske Krig styrtede
Witt, og Mængden, der ansaa ham for en Forræder,
sønderrev ham og hans Broder paa den skrækkeligste
Maade. Omtrent paa samme Tid ønskede den
store Condé, der kommanderede i Utrecht, at tale med
den berømte Filosof og sendte ham et Lejdebrev.
Efter sine Venners Ønske fulgte Spinoza Indbydelsen;
men da han kom, var Condé pludselig bleven kaldt
bort. Derimod modtog Marschallen af Luxembourg
ham med stor Opmærksomhed. Man opfordrede ham
til at dedicere et Værk til Ludvig XIV, som saa vilde
give ham en aarlig Pension; men derpaa vilde Spinoza
ikke indlade sig. Da han kom tilbage til Haag, stimlede
Folk sammen uden for hans Bolig paa en truende
Maade, fordi Rygtet gik, at han stod i forræderisk
Forbindelse med Fjenden. Hans Vært blev forskrækket
herover; men Spinoza beroligede ham og erklærede, at
„saa snart Mængden gjorde Mine til at angribe Huset,
vilde han gaa ud til den, selv om han skulde lide
% --- p. 34 ---
\setcounter{footnote}{0}%
\opage{34}samme Skæbne som de stakkels Brødre de Witt; han
vidste med sig selv, at han var en god Republikaner,
der aldrig havde haft andet for Øje end Statens Ære
og Fordel“.

Ikke blot i Frankrig, ogsaa i Tyskland var man
bleven opmærksom paa Spinoza. I Aaret 1673 modtog
han fra Kurfyrsten af Pfaltz, den ved sin Interesse for
Videnskab bekendte Karl Ludvig, gennem Professor
Fabritius en Opfordring til at overtage et Professorat
i Filosofi i Hejdelberg. Der blev lovet ham den
største Frihed til at filosofere, naar han ikke krænkede
den i Staten herskende Religion. Men Spinoza modtog
ikke Indbydelsen. „For det første“, skrev han i
sit Svar, „vil jeg ikke kunne føre Filosofien videre,
naar min Tid er optagen af Undervisning. Dernæst
véd jeg ikke, hvor Grænserne for den omtalte Frihed
til at filosofere ligge, saa jeg ikke faar Udseende af at
ville krænke den herskende Religion. Religiøse Stridigheder
opstaa ikke saa meget af brændende Religionsiver,
som af Menneskenes forskellige Lidenskaber og
deres Modsigelseslyst, der leder dem til at nedsætte
og fordømme alt, selv det rette. Og da jeg allerede
har erfaret dette i det private og ensomme Liv, jeg
nu fører, maa jeg frygte det endnu mere, dersom jeg
modtager en saadan Værdighed.“ Med god Grund
maatte Spinoza være betænkelig ved at optræde som
offentlig Lærer i Filosofi; han vilde næppe kunne fremsætte
en eller to af sine Læresætninger, før man vilde
betragte Statsreligionen som krænket. Paa den Tid
kunde der, trods alle gode Løfter, ikke være Tale om
% --- p. 35 ---
\setcounter{footnote}{0}%
\opage{35}virkelig Lærefrihed\footnote[1]{Endnu i vore Dage er det hos os blevet udtalt, at en Universitetslærer ikke maa angribe Statskirkens Lærdomme. Et af vore Dagblade mente, at afdøde Professor Brøchner i den Henseende havde misforstaaet sin Stilling. Men en offentlig Lærer i en Videnskab kan dog fornuftigvis ikke være bunden til andre Love end dem, der ligge i selve Videnskabens Natur; de ville sikkert indeholde al fornøden Garanti. Det maa netop være hans højeste Pligt at udsondre af Sandhedens Rige alt, hvad der uberettiget har indsneget sig deri. Stats- og Folkelivet vilde til sidst sygne hen, hvis der blev givet visse Lærdomme Privilegium som Undtagelser fra denne Regel.}. --- Det synes, som om Kurfyrsten
kun har kendt Spinoza som Forfatter af Skriftet
om Descartes's Filosofi, ikke som Forfatter til den
teologisk-politiske Afhandling. Alligevel var det Tegn
paa Frisindethed at ville ansætte en banlyst Jøde, der
ikke var traadt over til Kristendommen.

Spinoza opfattede ikke selv sit Forhold til Folkereligionen
som et fjendtligt. Den sande Religiøsitet
afhang efter hans Overbevisning ikke af nogen som helst
Trosbekendelse. De, der ikke ad Tænkningens Vej
kunde naa Erkendelsen af det Guddommelige, maatte
finde deres Hvile i en positiv Religions Lærdomme.
Han samtalede gerne med sin Vært og hans Familie,
naar de havde været i Kirke, og udspurgte dem om,
hvad Udbytte de havde haft deraf. I Voorburg var
han med at underskrive et Andragende om Besættelsen
af et Præsteembede, og i Haag hørte han nogle Gange
sin Værts Præst, hvem han roste meget. Da hans
Værtinde spurgte ham, om han troede, at hun kunde
blive salig i den Religion, hun tilhørte, svarede han:
„Eders Religion er god, I og skal ikke søge nogen
% --- p. 36 ---
\setcounter{footnote}{0}%
\opage{36}anden eller tvivle om, at den vil føre eder til Frelse,
naar I blot beflitter eder paa Fromhed og fører et
fredeligt og stille Liv“. Han havde et altfor klart Blik
for, hvor meget der hører til for selv at danne sig en
Livsanskuelse, naar man har forladt den overleverede,
til at han skulde føre nogen bort fra denne, der endnu
ikke kunde staa paa egne Ben. Men han saa naturligvis
ogsaa, hvor meget de overleverede Lærdomme kunne
hindre Fremgangen i Sandhedens Erkendelse, i det den
menneskelige Ladhed gerne slaar sig til Ro ved det,
man har faaet indpodet som Børn. Størst Ulykke afstedkomme
de dog ved den Kiv og Strid, de volde;
det er de forskellige Meninger, man faar ved Overlevering,
der splitte Menneskene ad, medens de ved at
følge Naturen vilde naa til Enighed, som Børn af en
og samme Moder.

Ligesom Colerus fremhæver hans Agtelse for andres
Tro, siger ogsaa en anden Teolog, der en Snes Aar
efter Spinozas Død samlede Efterretninger om ham,
Sebastian Kortholt, som ellers fordømmer ham i meget
fanatiske Udtryk, at „der aldrig hørtes nogen Ed eller
nogen fræk Tale om Gud fra Spinozas Mund“. Derimod
fortæller den reformerte Teolog Limborch, at han
--- meget uventet --- traf sammen med Spinoza ved et
Gæstebud, og at denne under Bordbønnen viste sit
ugudelige Sindelag, „i det han ved sine Miner syntes at
at spotte over de bedendes Dumhed“. Denne Beretning
strider aldeles mod, hvad der ellers vides om Spinozas
Karakter og Adfærd, og er sikkert kun et Vidnesbyrd
om, hvor oplagt Limborch var til at finde „Tegn paa
Ugudelighed“ \textit{(signa animi irreligiosi)} hos den Fri%
% --- p. 37 ---
\setcounter{footnote}{0}%
\opage{37}tænker, han --- meget mod sin Vilje --- var kommen i
Selskab med. Den gode Teolog kan desuden, som man
med Rette har bemærket, ikke have været meget fordybet
i sin Bøn, naar han har siddet og iagttaget Spinozas
Ansigt.

Spinozas Samtale skal have været aandfuld og livlig
og øvede stor Tiltrækning paa mange. Hans store
Skarpsindighed og fine Ironi gjorde det, han sagde,
baade belærende og underholdende. Hans Optræden
bar altid Præget af Mildhed og Selvbeherskelse; man
saa ham hverken meget bedrøvet eller overvættes glad,
og naar Følelsen blev for stærk, trak han sig hellere
tilbage i Steden for at give den Luft over for andre.
I den sidste Del af hans Liv var hans Helbred i høj
Grad svækket; han led af en tærende Brystsygdom,
som hans anstrengte Studier havde næret og maaske
fremkaldt. I denne Sygelighed kan man maaske søge
Grunden til hans ringe Produktivitet i Slutningen af
hans Levetid. „Etiken“ synes allerede at have været
færdig i Aarene 60, og i disse Aar arbejdede han ogsaa
paa den „Afhandling om Forstandens Forbedring“,
der fandtes ufuldendt efter hans Død. I Aaret 1675
gav han efter for sine Venners Anmodninger og rejste
til Amsterdam for at faa Etiken trykt. Men det Rygte
udbredte sig nu strax, at han vilde udgive en Bog, som
skulde bevise, at der ingen Gud gaves. Teologerne
klagede over ham ved Hoffet og hos Øvrigheden, og
Kartesianerne, som man ansaa for hans Aandsfrænder,
stemte i med for at fralægge sig al Meddelagtighed.
Spinoza besluttede derfor at opsætte Udgivelsen til
bedre Tider; men dem oplevede han ikke. Først efter
% --- p. 38 ---
\setcounter{footnote}{0}%
\opage{38}hans Død udgav Ludvig Meyer Etiken, en ufuldendt
politisk Afhandling, den ufuldendte Afhandling om
Forstandens Forbedring, en ufuldendt hebraisk Grammatik
og en Samling Breve.

„Afhandlingen om Forstandens Forbedring“ har
bl. A. den Interesse, at vi deraf se, hvorledes Spinozas
Tænkning i sin sidste Grund udspringer af et etisk
Motiv. Han begynder dette Skrift paa følgende Maade:
„Efter at Erfaringen har lært mig, at alt, hvad der møder
os i vort sædvanlige Liv, er tomt og intetsigende, og
da jeg har set, at der ikke er noget godt og ondt i de
Ting, vi mest frygte for, uden for saa vidt Sindet lader
sig bevæge af dem, har jeg besluttet at undersøge, om
der gives noget, som er et sandt Gode, som formaar at
meddele os sit Væsen, og som i sig selv kan paavirke
Sjælen, selv om alle andre Goder forkastes, med andre
Ord, om der gives noget, ved hvis Erhvervelse og Besiddelse
jeg kunde nyde en vedvarende, evig og højeste
Salighed.“ I Modsætning til de falske og forgængelige
Goder stilles saa som det sande og højeste Gode Erkendelsen
af det evige og vor Enhed dermed, eller Erkendelsen
af vor Aands Sammenhæng med hele Naturen. For
at naa dette Maal maa vi vinde en sand Erkendelse af
Tingenes Væsen og stræbe efter, at saa mange som
muligt dele den med os. Der maa altsaa foretages en
Lutring og Reformation af Erkendelsen og dannes et
Samfund, i hvilket saa mange som muligt kunne opdrages
og udvikles til det angivne Maal.

I sit Hovedskrift „Etiken“ \textit{(Ethica ordine geometrico
demonstrata)} har Spinoza gennemført den matematiske
Metode, han allerede havde anvendt i sin Fremstilling
% --- p. 39 ---
\setcounter{footnote}{0}%
\opage{39}af Descartes's Lære, og som kun paa en meget tvungen
Maade kan benyttes ved Behandlingen af metafysiske,
psykologiske og etiske Spørgsmaal. Imidlertid vidner
den Maade, hvorpaa Spinoza anvender denne Metode,
ikke blot om hans store Skarpsindighed, men medfører
ogsaa Klarhed og Gennemskuelighed i Tankegangen,
noget, hvorpaa andre spekulative Metoder ellers lide
stor Mangel. Desuden giver Spinoza i Fortaler og Tillæg
til de enkelte Bøger og i Anmærkningerne (Skolierne)
til de enkelte Sætninger sammenhængende Udviklinger,
der til Dels berøre det interessanteste i Bogen,
og hvortil man ved første Læsning især bør holde sig.\footnote[1]{Foruden paa Latin kan man nu ogsaa læse Spinoza i to tyske Oversættelser (af Auerbach og af Kirchmann) og i en fortræffelig fransk Oversættelse (af Saisset).}
„Etiken“ er delt i 5 Bøger. Den første (de deo) handler om
Gud som den absolute Substants, om hans Attributer
og om Tingenes Forhold til ham, og giver altsaa
Grundtrækkene af Spinozas almindelige Verdensanskuelse.
Den anden Bog \textit{(de natura et origine mentis)}
handler om Menneskets Væsen, dets Forhold til Gud
og til andre endelige Væsener, særlig om Forholdet
mellem Aand og Materie og om den menneskelige Erkendelse.
Tredje Bog \textit{(de origine et natura affectuum)}
handler om Lidenskabernes Væsen og Oprindelse og
giver en mesterlig psykologisk Skildring heraf. I fjerde
Bog \textit{(de servitute humana s. de affectuum viribus)} vises,
hvilket Herredømme Lidenskaberne kunne vinde over
Mennesket. Det er Skildringen af Mennesket i hans
Trældomstilstand, som gradvis gaar over til Frihed,
% --- p. 40 ---
\setcounter{footnote}{0}%
\opage{40}efter Haanden som hans sande Væsen udvikler sig.
Endelig skildrer femte Bog \textit{(de potentia intellectus s.
de libertate humana)} Mennesket i hans ideale Væsens
evige Virkelighed. Saaledes giver dette storartede Værk
os, efter Herbarts træffende Lignelse, Billedet af et
Drama i fem Akter. Det har Exposition, Stramning
og Løsning af Knuden. Først kommer Gud og den
menneskelige Aand; derpaa følger det, der nedtrykker
Aanden, Affekterne, og Skildringen af den menneskelige
Trældom, endelig den aandelige Frihed.

Til „Etiken“ slutter sig den „politiske Afhandling“,
som Spinoza arbejdede paa, da han døde. To Ting
forudsattes jo for det højeste Godes Virkeliggørelse:
den sande Erkendelse, som „Etiken“ viser os i dens
Anvendelse paa Menneskets Forhold til Lidenskaberne,
og det rette Samfund, hvis Væsen fremstilles i den
„politiske Afhandling“, der undersøger Naturretten, Statsmagten,
Samfundets Oprindelse og Øjemed.

Men det er ikke blot ved sine dybsindige og frisindede
Skrifter, at Spinoza har virket for Filosofiens
og den frie Forskens Sag. Hele hans Liv, hvis Renhed
og ideale Retning selv Modstandere maatte erkende,
har været et vigtigt Argument i Striden mellem Filosofi
og Teologi. Det laa i de Tiders Aand at anse det
for en Umulighed, at en Mand, der stillede sig uden
for den positive Religion, kunde være et godt Menneske
og føre et sædeligt Liv. Spinozas Liv blev derfor et
teologisk Problem. Den første, som efter Spinoza selv
med Energi udtalte, at det gode og rette har sin
Sanktion i sig selv og ikke er afhængigt af Avtoritetsbud,
var Pierre Bayle. Han henviste netop særlig til
% --- p. 41 ---
\setcounter{footnote}{0}%
\opage{41}Spinoza (hvis sande Betydning han ellers manglede Blik
for) som Bevis for den Sætning, at Sædelighed ej er
afhængig af positiv Religion. I vor Litteratur ere Holberg
og Eilskov paavirkede af Bayles Ideer; dog fremhæve
de navnlig en anden Sætning af ham, nemlig at det
mere er Naturel og Temperament end Trossætninger,
der bestemme Menneskenes Handlemaade. Eilskov
mener, at naar Fritænkere ikke leve ugudelig, maa det
forklares af, „at de mere har fulgt deres Gemyts Beskaffenhed
end Dydens Love; de har levet saa, mere
fordi de ingen naturlig Drift fandt hos sig til at leve
anderledes, end af Frygt for, at de ellers skulde overtræde
deres Pligter. Det er uden Tvivl Grunden til
Spinozas stille og ærbare Levemaade“. Ved denne
Opfattelse blev, som Holberg siger, „den hele Disputats
forgæves“ (Ep. 210), det vil sige, Spørgsmaalet
blev bragt ud af Verden ved Erkendelsen af Karakterens
og Handlemaadens relative Uafhængighed af
Teorier og Anskuelser, baade hos troende og ikketroende.
Det vilde ogsaa være urigtigt at tænke sig
Spinoza som en Helgen; hos ham som enhver anden
er den moralske Karakter sikkert bleven bestemt ved
hans oprindelige Temperament, og ikke uden Ret har
man ogsaa henvist til hans Brystsygdom som Forklaringsgrund
for hans tilbageholdne, resignerede Levevis.
Maaske vil man ogsaa kunne bebrejde ham en vis
Passivitet; hans Begejstring var i hvert Tilfælde den
stille, i sig selv hvilende og førte ham ikke til --- som
en Giordano Bruno eller I. G. Fichte --- at byde Tidsalderen
aaben Trods i Kampen for sine Ideer. Men
det store og smukke i hans Karakter er paa den anden
% --- p. 42 ---
\setcounter{footnote}{0}%
\opage{42}Side netop Harmonien mellem hans Personlighed og
hans Lære; i Troen paa den evige Naturgrund, der
rører sig i alt og i alle, ligger jo ogsaa Forvisningen
om, at det sande er det evig sejrende. Og hans Liv
beholder stadig den Betydning, at det ved sit Vidnesbyrd
har bidraget til, at uvedkommende moralske Hensyn
kunne skydes til Side ved Diskussionen af filosofiske
og videnskabelige Spørgsmaal.

Var det i hine Tider et Spørgsmaal, hvor vidt en
Fritænker kunde leve sædelig, saa var det endnu større
Tvivl underkastet, om han kunde dø rolig. Der var
derfor ogsaa mange forskellige Rygter i Omløb angaaende
Spinozas Død. Nogle vilde vide, at han havde
bedøvet sig med Mandragorasaft, da han mærkede
Dødens Komme, andre, at han havde udstødt fortvivlede
Udraab og forbudt sin Værtinde at lade nogen Præst
faa Adgang til ham, for at ingen skulde blive Vidne
til hans Svaghed. Heldigvis have to Teologer kort
efter Spinozas Død taget sig for at undersøge Sandheden
af disse Rygter, og begge disse „uvillige“ Vidner
erklære dem for aldeles grundløse. Spinozas Vært og
og Værtinde levede endnu og kunde berette alt, hvad
der vidstes om hans Død. Hans Brystsvaghed var stedse
tagen til; men hverken han selv eller hans Omgivelser
ventede, at han skulde dø saa snart. Dagen forud
havde han været nede hos sin Vært og røgt en Pibe,
medens han samtalede med ham om Dagens Prædiken,
og endnu om Morgenen paa den Dag, han døde (den
21de Februar 1677), var han ligeledes nede hos Værten
og hans Familie. Des mere forundrede bleve disse, da
de om Eftermiddagen kom hjem fra Kirke, ved at høre,
% --- p. 43 ---
\setcounter{footnote}{0}%
\opage{43}at han var død Kl. 3. Der havde ingen været hos
ham uden hans Ven Lægen Ludvig Meyer fra Amsterdam,
som rejste bort igen samme Aften. Den ene af
vore to Hjemmelsmænd, Seb. Kortholt, siger udtrykkelig,
at han „drog sit sidste Suk i Fred“, skønt han tilføjer,
at de lærde have været uenige om, hvor vidt en
saadan Død tilkom en Ateist. Og Colerus gendriver
alle Rygterne, skønt han kort efter forarges over, at
Spinoza i Regninger, der indløb til Boet, kaldes „\emph{Salig}
Hr. Spinoza“. Men skønt Colerus betragtede Spinoza
som fordømt i den anden Verden, maa Filosofiens Historie
dog være den brave Præst meget taknemmelig,
fordi han i sin Biografi, der i al sin Naivitet er en
interessant lille Bog, saa ivrig har stræbt for at rense
hans Minde i den nærværende Verden. Det er tillige
et forsonende Træk, at Spinoza, Resignationens og
Tolerancens Talsmand, har fundet en sandhedskærlig
og upartisk Biograf i sine Modstanderes Lejr.

% --- p. 44 ---
\setcounter{footnote}{0}%
\opage{44}\begin{center}Spinoza og hans Korrespondenter.\end{center}


Før end vi gaa over til en sammenhængende Fremstilling
af Spinozas filosofiske Lære, vil det have sin
Interesse at dvæle ved hans Brevvexling. Hans Ideer
træde i denne til Dels frem i en noget friere Form
end i hans Skrifter og blive derved lettere tilgængelige.
Det tjener derhos til Karakteristik af hans historiske Stilling
at se, hvorledes hans nærmeste Samtid har modtaget
hans Ideer. En Fremstilling af Forholdet mellem Spinoza
og hans Korrespondenter danner altsaa en passende
Overgang mellem Biografien og Systemet.

I Indledningen have vi fremhævet Modsætningen
mellem den empiriske Detailforsken og den spekulative,
deduktive Tænkning. Denne Modsætning kommer paa
en interessant Maade frem i de mellem ham og hans
Ven Henrik Oldenburg vexlede Breve.

Oldenburg var født i Bremen og var under Cromwell
og Karl II. den nedersaksiske Kredses Gesandt i
London. Han interesserede sig meget for fysiske og
kemiske Experimenter, var den første Sekretær ved det
1662 stiftede Royal Society og en god Ven af den berømte
Boyle, Grundlæggeren af den analytiske Kemi.
% --- p. 45 ---
\setcounter{footnote}{0}%
\opage{45}Oldenburg var ikke nogen stor Aand, men var begejstret
over sit Bekendtskab med Spinoza, hvem han
besøgte i Rhynsburg, og over det kraftige Opsving, den
experimentale Forsken paa den Tid tog i England.
Han fandt Glæde i at være Mellemmand mellem sine
to geniale Venner, men er ikke fri for at pynte sig selv med
sit Forhold til dem. Igennem ham førtes der en Diskussion
mellem dem, fremkaldt ved et Skrift af Boyle,
som var blevet tilsendt Spinoza. Spinoza opponerede
imod, at Boyle af sine Forsøg vilde slutte, at Forskellen
mellem Legemernes Egenskaber kunde føres tilbage
til Forskelle i Bevægelse, Figur og andre mekaniske
Forhold. Spinoza mener, lige som Descartes, at alt i
Naturen maa forklares efter Mekanismens Love, men
nægter, at denne Forklaringsmaades Berettigelse kan
begrundes ved Experiment. Det Bevis, her kan være
Tale om, anser han for givet af Descartes og Baco.
Naar alle Fænomener i Naturen forudsætte Bevægelse,
og Bevægelsen igen Udstrækning i Rummet, saa er
Materiens Væsen ét med Udstrækningen, og Udstrækningens
Love, Geometriens Sætninger, blive altsaa de
egentlige Naturlove, hvortil alt maa føres tilbage.
Spinoza manglede Blik for den induktive Forskens
Betydning. Hvad han søgte, var Tingenes Væsen, udtrykt
i klare Begreber. Men Erfaringen lærer os, som
han i et af sine andre Breve siger, ikke Tingenes nødvendige
Væsen, men kun det virkelig existerende. Det
sande værende var for Spinoza ikke de enkelte Ting
og Fænomener, men de almindelige Love og det uendelige
Væsen, der ligger til Grund for alt. Erfaringen
kan kun have den Betydning at lede vore Tanker i
% --- p. 46 ---
\setcounter{footnote}{0}%
\opage{46}visse bestemte Retninger, men kan ikke give os en sand
Viden. Boyle derimod gik netop ud paa Undersøgelse
af det enkelte og specielle og lod Tingenes indre Væsen
staa hen. Han samlede Iagttagelser og analyserede de
enkelte Stoffer, men tragtede ikke efter at opstille nogen
Teori. Det var ogsaa netop omhyggelig Iagttagelse, Kemien
trængte til for at komme ud over Alkymien, og den
skylder Boyles skeptisk -e mpiriske Aand meget. Selv
en Ikke-Fagmand kan endnu med Interesse og Belæring
lære hans „\textit{Sceptical Chymist}“ og de andre Afhandlinger,
hvori han bekæmper den peripatetiske og den
alkymistiske Kemi. Oldenburg opfordrer sine Venner
til med Iver at benytte hver sin Metode. „Dig“, skriver
han til Spinoza, „beder jeg særlig, at du med din
matematiske Skarpsindighed vil vedblive at grundlægge
Tingenes Principer, lige som jeg stadig opfordrer min
ædle Ven Boyle til at bekræfte og oplyse Videnskaben
ved gentagne og nøjagtige Forsøg og Iagttagelser.“
Dette var i Virkeligheden ogsaa den bedste Løsning
af en saadan Strid. Ingen af de to Metoder kan
undværes; hvad det gælder om er, at hver finder sit
rette Omraade og sine sande Grænser.

Ikke saa heldig var Oldenburg i en anden Diskussion
med Spinoza. Det viste sig her, hvad der i Indledningen
blev berørt, at den rent empiriske Forskerdrift
kan være saare begrænset; paa det enkelte Omraade,
man har valgt sig, søger man klare Grunde og
Beviser, men kan derfor godt i sin Livsanskuelse følge
den første den bedste Tradition og tage Fornuften
fangen. Sit første Brev til Spinoza skrev Oldenburg
kort efter sit Besøg i Rhynsburg. Han havde følt sig
% --- p. 47 ---
\setcounter{footnote}{0}%
\opage{47}saa tiltrukken af Spinoza, at han strax følte Trang til
at komme i Forbindelse med ham, især da meget af,
hvad de havde talt om, endnu ikke stod klart for ham.
Hvad han især standsede ved, var Spinozas Opfattelse
af Gud som den uendelige Substants, hvis Ytringsformer
alle Ting ere. Da Oldenburg stedse gaar ud fra de
sædvanlige teologiske Forestillinger, er det intet Under,
at han ikke forstaar Spinoza. Det synes, som om denne
sér dette og gerne griber den kemiske Forhandling, han
gennem Oldenburg kom i med Boyle, for at ende den
ufrugtbare Diskussion. Dog fremtræder Oldenburg paa
dette Stadium som den ivrige og begejstrede Ven af
den frie Forsken. Han opfordrer Spinoza til at udgive
sine Skrifter og tage Bladet fra Munden trods Teologernes
Bjæffen \textit{(qvicqvid theologastri ogganire poterint)}.
„Eders Republik“, skriver han, „nyder jo en saare stor
Frihed. Saa bør man ogsaa filosofere i den med Frihed.
Aflæg al Frygt for at ophidse vor Tids Smaamennesker!
Længe nok har Uvidenheden og Gøgleriet
haft frit Spil: lad os nu hejse den sande Videnskabs
Flag og endnu dybere end tilforn udforske Naturens
Hemmeligheder!“ Denne Iver blev imidlertid noget
kølnet ved det nærmere Kendskab til Spinozas Lære
og navnlig ved den „teologisk -p olitiske Afhandling“ og
den store Forargelse, den vakte. Baade Oldenburgs
og Boyles Anskuelser vare ortodoxe. En anden af
Spinozas Korrespondenter, Fysikeren og Filosofen
Tschirnhausen, som paa denne Tid opholdt sig i London,
fandt, at Boyle og Oldenburg havde fattet „underlige
Forestillinger“ om Spinozas Person; dog lykkedes
det ham at betage dem disse Fordomme, saa de igen
fik en gunstig Mening om ham. Efter en halv Snes
% --- p. 48 ---
\setcounter{footnote}{0}%
\opage{48}Aars Afbrydelse genoptages Brevvexlingen, og Oldenburg
begynder med at erklære, at han har haft en
ugunstig Mening om Spinoza, fordi han anlagde den
Maalestok paa dennes Anskuelser, „som de sædvanlige
Teologer og de overleverede Trosbekendelser yde“; nu
derimod har han indsét, at det var en forhastet
Dom. Han ønsker nærmere at indvies i Spinozas Lære
og tilbyder sin Medvirken til at skaffe den Indgang hos
forstandige og gode Mænd og forsvare den mod Fordomme.
Men da han hører, at Spinoza vil udgive sin Etik,
formaner han ham indstændig til ikke at udtale noget,
som paa nogen Maade kunde synes at rokke ved Fromhed
og Dyd --- „især da denne vanslægtede og onde
Tid ikke søger noget med større Begær end Lærdomme,
der kunne forsvare de herskende Laster“. Spinoza udbeder
sig en nærmere Forklaring om, hvad det er for
Lærdomme, der kunne krænke Fromhed og Dyd; ti
hvad der stemmer med Fornuften, anser han ogsaa for
at være Dyden mest tjenligt. Oldenburg nævner da tre
Punkter. Spinoza synes at sammenblande Gud og Naturen.
Han ophæver Miraklernes Sandhed og Gyldighed,
og det er dog paa dem, Aabenbaringens Vished
beror. Endelig er hans Mening om Kristi Person
uklar. Naar han blot i disse Punkter kan tilfredsstille
alle sande kristne, vil hans Sag være uden Fare.
Spinoza svarer omtrent paa følgende Maade. 1) Jeg
opfatter Gud som iboende (immanent) Aarsag; han
virker ikke ud over sit Væsens Grænser, ti hans Væsen
er uendeligt og har intet uden for sig. Ogsaa Apostelen
Pavlus siger jo, at i Gud ere, leve og røres vi. Og
naar man beskylder mig for at sammenblande Gud og
% --- p. 49 ---
\setcounter{footnote}{0}%
\opage{49}Naturen, saa forstaar man ved Natur den ydre Masse,
den legemlige Materie, og misforstaar mig derfor ganske.
2) At grunde Aabenbaringen paa Mirakler er at grunde
den paa Uvidenhed. Ingen véd, hvad Naturens Kræfter
kunne udrette, og fordi vi i et enkelt Tilfælde ikke
kunne finde de naturlige Aarsager, have vi ikke Lov
til at benægte deres Tilværelse. Religionen bør bedømmes
efter Lærens Sandhed; ellers adskiller den sig
ikke fra Overtroen, som bygger paa Uvidenheden.
3) Det er ikke nødvendigt at kende Kristus efter
Kødet; hvad det kommer an paa, er at være delagtig
i Guds evige Søn, det vil sige, den evige, guddommelige
Visdom, der har aabenbaret sig i alt, mest i den
menneskelige Aand, og allermest i Kristus. Men nogle
Kirkers Lære, at Gud har antaget menneskelig Natur,
forstaar jeg ikke; det er for mig lige saa absurd som
at sige, at Cirkelen har antaget Kvadratets Natur.
--- I dette sidste Punkt er Spinoza Forgænger for den
nyere spekulative Religionsfilosofi, og han har særlig
udtalt en Tanke, som Strauss videre udviklede i Slutningsafhandlingen
til sit berømte „Leben Jesu“. --- Oldenburg
kommer nu først frem med den egentlige Hovedindvending:
at Spinoza lærer en evig Nødvendighed i
alt, hvoraf maa følge, at ingen er ansvarlig for sine
Handlinger. Spinoza svarer, at den nødvendige Sammenhæng
i alt netop er Grundlaget for hele hans Lære.
Men, siger han, jeg lærer ikke Fatalisme; den Nødvendighed,
som er Udtryk for et Væsens egen Natur,
er ét med den sande Frihed, er ikke en ydre, men en
indre Nødvendighed og ophæver ikke Forskellen mellem
godt og ondt.

% --- p. 50 ---
\setcounter{footnote}{0}%
\opage{50}Det var Spinozas Særkende og hans Storhed, at
han ej blot vilde Tænkning til en vis Grad, ikke stansede
paa Halvvejen. Hans Stræben efter en sammenhængende
videnskabelig Verdensanskuelse var fuldt
berettiget, selv om han nærede Illusioner om den spekulative
Tænknings Evne. Et filosofisk Problem kan,
naar det én Gang er stillet, ikke trænges tilbage af
Hensyn til hvilke som helst andre Forudsætninger.
Men mange tale med tilsyneladende Anerkendelse om
Fornuft og Tænkning, og naar det saa bliver Alvor,
naar Anerkendelsens Ærlighed skal vise sig i, at man
bøjer sig for Resultatet, saa springe de fra, som Soldater,
der faa Kanonfeber. Det kan være en lang og
møjsommelig Vej for den enkelte, inden han faar bragt
Harmoni mellem sine Følelser og Interesser og den
klare Sandhed, Erkendelsen viser ham. Men Vejen
maa og skal tilbagelægges, tidligere eller senere. Ad
denne Vej gaar Slægten frem, lige som den enkelte i
Stilhed gaar frem ad den. --- Oldenburg hørte til dem,
der vende om, naar Kuglerne begynde at pibe. Og
hos ham er det saa meget mere fremtrædende, som
han begyndte med at være saa rask paa det og optraadte
som den begejstrede Fremskridtsmand.

Mere Sympati føler man i denne Henseende med
en anden af Spinozas Korrespondenter, skønt det gik
ham paa lignende Maade. En Købmand fra Dortrecht,
Willem van Bleyenbergh, syslede i sine ledige Timer
med filosofiske Studier og havde gentagne Gange læst
Spinozas Fremstilling af Descartes's Filosofi. Han
skriver nu til Spinoza og fremstiller sig som en Mand,
der ledes af Trang til ren og ægte Sandhed og søger
% --- p. 51 ---
\setcounter{footnote}{0}%
\opage{51}at tilfredsstille denne Trang ved Videnskaben. Men
han er, lige som Oldenburg, frastødt af Spinozas Determinisme.
Naar Gud virker i alt, siger han, er det
jo ogsaa Gud, der virker i Adams Synd, og altsaa maa
enten han selv være ond eller Synden ikke være noget
virkeligt og positivt. --- Spinoza kommer ham i Møde
med stor Venlighed. „Intet“, siger han, „sætter jeg
større Pris paa end at indgaa Venskab med Mænd,
der oprigtig elske Sandhed. Et saadant Venskab er
grundet paa den Kærlighed, hver enkelt af dem har
til Sandhedens Erkendelse, og derfor kan man lige saa
lidt ophæve det, som man kan undlade at slutte sig til
den Sandhed, man én Gang har erkendt. Intet uden
Sandheden kan bringe dyb Forbindelse mellem forskellige
Sindelag og Karakterer.“ Med Hensyn til det af
Bleyenbergh opkastede Spørgsmaal udtaler Spinoza, at
Begreber som godt og ondt opstaa og anvendes, naar
vi sammenligne Handlinger og Begivenheder indbyrdes
efter en vis Maalestok. Men vore relative Begreber
kunne ikke anvendes paa det absolute Væsen. Gud
kan lige saa lidt karakteriseres ved Begreber, der hentes
fra den menneskelige Natur, som et Menneske ved
Begreber, der hentes fra Elefantens Natur. (I et andet
Brev siger Spinoza, at dersom Trianglen kunde tale,
vilde den sige, at Gud var „triangulær i ophøjet Betydning“
--- \textit{eminenter triangularis,} --- og Cirkelen, at
han var „cirkulær i ophøjet Betydning“ --- \textit{eminenter
circularis.)} Vi tør derfor heller ikke tillægge Gud
Had eller Vrede, som Teologerne gøre, og som Skriften
synes at gøre, da den lemper sig efter Almuens Forestilling,
for hvem Gud staar som Lovgiver og Dommer
% --- p. 52 ---
\setcounter{footnote}{0}%
\opage{52}og ikke som det evige og uendelige Væsen, der ibor
alt, og hvoraf alt udspringer. Fordi alt sker med Nødvendighed,
udslettes Forskellen mellem de gode og de
slette ikke, lige saa lidt som Forskellen mellem Glæde
og Sorg. Den gode lever i den klare Bevidsthed om
sin Harmoni med Ideen om Gud, der bestemmer alle
hans Handlinger, medens den ondes Bevidsthed er opfyldt
endelige og jordiske Ting. --- Bleyenbergh

af. forklarer
nu sit Standpunkt nærmere. Han har to Regler,
hvorefter han filosoferer: den ene er hans Forstands
klare og indlysende Begreb, den anden Guds aabenbarede
Ord og Vilje. Selv om noget staar for ham
som klart og indlysende, mener han dog, at han tager
fejl, naar det strider mod Skriften. Spinoza svarer hertil,
at Bibelen slet intet kan lære os i filosofisk Henseende.
Hvad den lærer, er, naar man anlægger en
filosofisk Maalestok, ofte aldeles uforstaaeligt, hvorimod
man aldrig kan tvivle om Sandheden af det, der bevises
klart og tydelig. Men hvad nytter det at diskutere,
naar den ene Part udtrykkelig erklærer, at en Bevisførelse
kun gør Indtryk paa ham, for saa vidt den
stemmer med, hvad Skriften --- efter hans egen eller
nogle Teologers Fortolkning --- lærer, og at den fornuftige
Erkendelse altsaa egentlig talt er ham ligegyldig?
--- Et Besøg, Bleyenbergh aflagde hos Spinoza
i Voorburg, bragte ikke nøjere Forstaaelse til Veje, og
Spinoza afbrød til sidst Brevvexlingen paa en venlig
Maade og henviste til mundtlig Forhandling. Senere
udgav Bleyenbergh Stridsskrifter mod Spinozas Værker.

Vi kunne ikke her nærmere undersøge, med hvilken
Ret den Beskyldning stadig fremsattes mod Spinoza,
% --- p. 53 ---
\setcounter{footnote}{0}%
\opage{53}at hans Nødvendighedslære (Determinisme) ophæver al
Frihed og Moralitet. I sidste Kapitel af mit Skrift
„Om Grundlaget for den humane Etik“ har jeg udførligere
behandlet de her hen hørende Spørgsmaal. Den
Vanskelighed, her kan være, stammer i hvert Tilfælde
ikke fra Determinismen i sig selv, men fra den spekulative
Opfattelse, der ligger til Grund. Og denne
Vanskelighed er, hvad Spinozas Kritikere stedse oversé,
fælles for Panteismen og Teismen. Enten man
antager en evig Substants, der rører sig i alt, eller en
almægtig Vilje, der skaber alt, bliver det lige vanskeligt
at forstaa, hvorledes et endeligt Væsen kan besidde
en absolut fri Vilje. Derfor besvare konsekvente Teologer
som Augustinus (og efter ham Descartes og Leibnitz)
Spørgsmaalet paa lignende Maade som Spinoza.
Forskellen er kun, at Teologien lægger Vanskeligheden
et Stykke længere tilbage, som i det hele de filosofiske
Vanskeligheder gerne komme igen i Teologien i højere
Potentser. Vanskeligheden er i dette Tilfælde opstaaet
ved den deduktive Metode, der jo ogsaa er fælles for
Teologien og den spekulative Filosofi. Anvende vi
derimod den induktive Metode paa de psykologiske og
etiske Spørgsmaal, saa indrømme alle, at vi kun forstaa
en Handling, naar vi kunne føre den tilbage til Individets
Karakter, tidligere Historie og Livsforhold. Der
kan i Principet ikke være noget i de menneskelige
Handlinger, som ikke paa denne Maade kunde forklares,
naar alle Forhold vare os bekendte. Determinismen
har derfor ubetinget Ret og anvendes uvilkaarlig i
Praxis af alle og enhver. Men deraf følger ikke, at vi
ere henviste til at lade os drive af Strømmen. Vi ere
% --- p. 54 ---
\setcounter{footnote}{0}%
\opage{54}stillede over for vor egen sjælelige og legemlige Natur
som over for den ydre Natur: kunne vi gøre os til
Herrer over Betingelserne, saa kunne vi ogsaa bestemme
Tingenes Gang, aflede eller hæmme Strømmen;
men om vi kunne det, kan i hvert enkelt Tilfælde kun
Erfaringen afgøre. Mirakler kunne vi hverken gøre i
den indre eller i den ydre Natur. --- Denne Betragtningsmaade
er i øvrigt, som vi siden skulle se, ikke
Spinoza fremmed, skønt den deduktive Metode er den,
der behersker hans Anskuelser.

Spinoza stod i venskabelig Forbindelse med Orobio
de Castro, en spansk Jøde, der efter skrækkelige Pinsler
i Inkvisitionens Fængsler var flygtet til Holland
og dér traadt over til sine Fædres Religion. Han
delte ikke Spinozas Anskuelser og skrev senere imod
ham. Han sendte Spinoza et Brev, hvori Lambert van
Welthuysen fra Utrecht, en Tilhænger af den kartesianske
Filosofi, udtaler sig om Spinozas „teologiskpolitiske
Afhandling“. Det er mærkeligt, men ikke
enestaaende, at Welthuysen, der i sin Forfattervirksomhed
viser sig som Rationalist og af Teologerne
blev beskyldt for Socinianisme, udtaler sig om Spinozas
Lære med langt større Heftighed end de ortodoxe
Oldenburg og Bleyenbergh. Forfatteren til det nævnte
Skrift, siger Welthuysen, har aflagt al Religion, i det
mindste gaar han videre end Deisterne, af hvilke disse
ugudelige Tider vrimle, og lærer indirekte (clam)
Ateisme, i det han lader Universet udspringe med
Nødvendighed af Gud, hvilket er det samme som at
gøre Gud til ét med Verden, og i det han lærer, at
Gudsdyrkelsen kun skal tjene til at fremme Dyd og
% --- p. 55 ---
\setcounter{footnote}{0}%
\opage{55}Fromhed, men ikke har nogen absolut Værdi i sig selv.
Saa blive alle Religionsformer jo lige gode, Koranen
stilles lige med Guds Ord, og man kan ikke bevise, at
Muhamed var en falsk Profet! --- Man har fundet
Kladden til Spinozas Svar og kan deraf sé, at han et
Øjeblik er kommen ud af sin kontemplative Ligevægt
og har brugt stærke Udtryk om den Beskyldning, der
her var rettet imod ham. Selv i det Brev, han sendte
Orobio de Castro til Svar, siger han, at han ikke véd,
om Brevskriveren har talt af Ondskab eller af Uvidenhed.
Kan man beskylde en Mand for at have aflagt
al Religion, naar han erkender Gud som det højeste
Gode, der bør elskes af et frit Sind, naar han i denne
Kærlighed finder vor højeste Lykke og vor sande Frihed
og sætter Dydens Belønning i selve Dyden? Er
der da intet i selve Dyden og Erkendelsen, der kan
glæde, men skal man kun se hen til, om der vanker
Løn og Straf, saa man lader sig tvinge som en Træl
og helst fulgte sine Lyster, naar det blot ikke var for
den truende Straf? Hvad Muhamed angaar, saa var
han ganske vist en falsk Profet, ti han nægtede den
Frihed, som den sande Religion tilsteder. Men Kristi
Aand maa dog siges at være tilstede baade hos Tyrker
og Hedninger, naar deres Gudsdyrkelse bestaar i Retfærdighed
og Menneskekærlighed. --- I et senere Brev
til Welthuysen selv anerkender Spinoza hans Stræben
efter Sandheden og hans Oprigtighed.

I Welthuysens Brev se vi Spinoza første Gang
blive kaldt Ateist, skønt indirekte (clam). Det var
den Opfattelse af ham, der i over 100 Aar var den
almindelige. Bayle siger om ham, at han er den første,
% --- p. 56 ---
\setcounter{footnote}{0}%
\opage{56}der har bragt Ateismen i System. Voltaire har givet
den samme Opfattelse Udtryk i et satirisk Digt \textit{(Les
systèmes)}:

Alors un petit juif, au long nez, au teint blème,

Pauvre, mais satisfait, pensif et retiré,

Esprit subtil et creux, moins lu que célébré,

Caché sous le manteau de Descartes son maître,

Marchant à pas comptés, s'approcha du Grand-Être:

„Pardonnez-moi“, dit il, en lui parlant tout bas,

„Je crois, entre nous, que vous n'existez pas“.
(Se bag i).
Det var --- og er til Dels endnu --- almindeligt at
kalde alle for Ateister, som afvige fra de herskende
religiøse Forestillinger. For at vide, hvad der forstaas
ved dette Ord, maa man derfor først vide, hvad man
forstaar ved Gud. I Regelen lægger man i Ordet
Ateisme tillige noget mere end den teoretiske Nægtelse
af Guddommen; da Gud staar som det højeste, man
kan tænke, maa den, der forkaster hans Tilværelse,
nødvendigvis --- saaledes slutter man --- være et lavt
og foragteligt Menneske, en, der fornægter, hvad der
giver Livet dets Værd og højere Betydning. Med
Hensyn til Udstrækningen af dette Begreb hersker der
derfor stor Uoverensstemmelse. Abbé Mersenne, Descartes's
Ven, siger, at Antallet af Ateister er saa
stort, at man maa undre sig over, at Gud lader dem
leve; i Paris alene er der 50,000, --- man kan finde
12 sammen i ét Hus! Men til Ateisternes Tal henregner
han saa ogsaa Charron, Macchiavel, Cardanus,
Campanella, Vanini og Fludd. --- Eilskov bemærker i
sine „Filosofiske Breve“, at de fleste kalde Spinozas
Lære for Ateisme, men mener ikke, man har Ret heri.
„Spinoza sammenblandede vel Gud med Verden, i det
% --- p. 57 ---
\setcounter{footnote}{0}%
\opage{57}han gjorde begge til én Substants; men han indrømmede
dog denne ene Substants Guds fornemste Kendemærke,
at den var til af sin Naturs Fornødenhed [ɔ: Nødvendighed],
følgelig at alt, hvad der skete i Verden,
var en Følge af Tingenes uforanderlige Sammenhæng.“
Eilskov vil derfor hellere kalde Spinoza for Deist end
for Ateist, hellere sige om ham, at han nægtede Forsynet,
end at han nægtede Gud. --- Hvad der adskilte
Spinoza fra de saa kaldte Deister var, at disse vel ikke
antog en Forsynsstyrelse i det enkelte og en overnaturlig
Aabenbaring, fordi alt siden Skabelsen gik til efter
de én Gang indrettede Naturlove, men dog antog en
personlig Gud. Hvad der forargede hos Spinoza, var
nu netop, at han nægtede Gud som et enkelt, bestemt
Væsen, Forestillingens og Fantasiens Genstand. Gud
stod for ham ikke uden for Naturen\footnote[1]{Smlgn. Heibergs „Protestantismen i Naturen“: Den Gud, du søger, er din egen Gud; Hvad vil du mere? Han træder ikke som en Genstand ud Blandt andre flere. I det du søger, du ham alt har fundet, I det du spørger, er alt Svaret vundet.}, men var det i
al Natur virkende og værende. Men en saadan Anskuelse
manglede man et Ord for, skønt det ikke var
første Gang, den var bleven udtalt. Ordet Panteisme
synes først at være bleven indført af den engelske Filosof
John Toland i Begyndelsen af det 18de Aarhundrede.
Det nye Ord var som et Varsel om den
Forandring i Opfattelsen, der skulde foregaa. Hen
imod Aarhundredets Slutning kappedes de største
% --- p. 58 ---
\setcounter{footnote}{0}%
\opage{58}Aander efter Lessings Exempel om at lovprise Spinoza;
den fromme Novalis kaldte ham „en af Gud beruset
Mand“, og Schleiermacher pegede i sine „Taler om
Religionen“ hen til Spinoza som en Religiøsitetens
Heros. Først nu blev den positive Side hos Spinoza
Genstand for almindeligere Anerkendelse; hidtil havde
man kun haft Øje for den negative. Saa længe den
dualistiske Naturopfattelse vedblev at herske, kunde det
heller ikke være anderledes. Spinoza havde ganske
Ret, naar han sagde, at hans Modstandere opfattede
Naturen paa en helt anden Maade end han. For ham
var det væsentlige i Naturen den evige Kraft og de
evige Love, der aabenbare sig i alt. Men en saadan
Værdi kunde hans Modstandere efter deres dualistiske
Anskuelser ikke tillægge Naturen; den stod for dem
som en Magt, der drager Mennesket ned ad, som endelig
og udvortes, som en Skranke, ud over hvilken de
higede.

Det er oplysende med Hensyn til Spinozas Gudsbegreb,
hvad han i et af sine Breve udtaler, at han
har lige saa klart et Begreb om Gud som om et Triangel,
men ikke saa klart et Billede; vi kunne tænke
Gud, men ikke forestille os ham. Forestillingen giver
os et enkelt Billede, afslutter Genstanden i en given
Ramme, men Tænkningen griber --- ikke den enkelte
Genstand for sig --- men Forholdet og Sammenhængen
mellem Genstandene. At Gud tænkes, ikke forestilles,
vil altsaa sige, at han opfattes som Loven eller Kraften
i alt, ikke som et enkelt, for sig bestaaende Væsen,
der har noget uden for sig. Netop fordi den sædvanlige
Opfattelse vil gribe Gud i et Billede, kommer den
% --- p. 59 ---
\setcounter{footnote}{0}%
\opage{59}til at sætte „Verden“ uden for ham. Dog er der næppe
den store Forskel mellem Forestilling og Tænkning,
som Spinoza antager. Skønt Tænkningen overvinder
de Modsætninger og den grelle Udvorteshed, hvori
Forestillingen let hilder sig, kan den dog selv ikke
undgaa at sætte nye Modsætninger; den statuerer Forhold,
Lovmæssighed og Sammenhæng, men forudsætter
derved bestandig en Flerhed af Led, mellem hvilke
Forholdene bestaa og Lovene gælde. En absolut
Enhed er det absolut ufattelige, netop fordi den ikke
vilde staa i Forhold til noget som helst. Vore Begreber
naa ikke ud over Fænomenernes og Relationernes
Verden. Derfor er det ikke sagt, at vor Verden
er den eneste eller den højeste; en konsekvent Erkendelseslære
fører til at antage en Uendelighed af
Muligheder. Hvad vi kunne hævde, er kun, at den
Sandhed og Nødvendighed, vi erkende, maa have sin
Plads og sin Gyldighed inden for en mulig større
Sammenhæng. Hvilken denne Plads er, savne vi ethvert
Middel til at afgøre. Men ved Arbejdet paa,
hvad der er sandt, skønt og godt, ere vi sikre paa at
arbejde i den sande Virkeligheds Tjeneste. For en
saadan monistisk Verdensanskuelse staar Spinoza som
den store Profet, skønt han gik ad spekulative Veje,
vi ikke kunne betræde med samme Tillidsfuldhed
som han.

--- Vi have paa et tidligere Sted omtalt et Brev fra
Spinozas trofaste Ven Simon de Vries, hvori denne
udtaler sin Sorg over at være skilt fra ham og priser
en Husfælle af Spinoza lykkelig, fordi han stedse har
Adgang til ham. I sit Svar siger Spinoza, at de Vries
% --- p. 60 ---
\setcounter{footnote}{0}%
\opage{60}ikke skal misunde hans Husfælle: „Der er ingen, som
er mig mere imod, og som jeg er mere forsigtig over
for. Derfor beder jeg dig og eder alle ikke at meddele
ham mine Anskuelser, før han har naaet en mere
moden Alder. Endnu er han for barnagtig og uden
Stadighed og tragter mere efter, hvad der er nyt, end
hvad der er sandt. Men jeg haaber, at han efter
nogle Aars Forløb vil rette disse Barnefejl, og saa
vidt jeg kan dømme efter hans Karakter, er jeg sikker
derpaa, hvorfor jeg ogsaa elsker ham paa Grund af
hans gode Natur.“ Det er uden Tvivl den samme
unge Mand, om hvem Spinoza i et andet Brev fortæller,
at han ikke vilde indvie ham i sine egne Anskuelser
og derfor havde udarbejdet en Fremstilling af
Descartes's Filosofi for ham. Og endelig genfinde vi
vistnok den samme Person i Spinozas Brevvexling under
Navn af Albert de Burgh. Han skriver til Spinoza fra
Florents og begynder med at fortælle, at han „ved
Guds uendelige Barmhjertighed er bleven Medlem af
den katolske Kirke“. Han omtaler sin tidligere Beundring
for Spinozas Forstand og Skarpsindighed og
for hans ivrige Søgen efter Sandhed; men des mere
sørger han nu over, at han trækkes om og forlokkes
af „de onde Aanders elendige og frække Fyrste“.
Hvad er hele din Filosofi andet end Illusioner og Fantasier?
Og dog bygger du paa den ikke blot din
aandelige Fred i dette Liv, men ogsaa din Sjæls evige
Frelse? Du tror, at du har fundet den bedste Filosofi;
men der har været og er jo saa mange Filosofer,
som ere uenige med dig! I din „teologisk -p olitiske
Afhandling“ har du med djævelsk Kløgt villet skille
% --- p. 61 ---
\setcounter{footnote}{0}%
\opage{61}Teologien fra Filosofien, skønt de hos dig selv falde
sammen. Du paastaar, at hvad du dér siger om
Skriftens Lære, har du lært af Skriften selv. Men
hvor tør du stole paa din egen Fortolkning? Kun ved
hele Kirkens Vidnesbyrd og de apostoliske Traditioner
kan Skriften fortolkes. Du kan ikke beraabe dig paa
Protestanterne, ti de ere lige saa elendige som du og
vandre i lige saa stort Mørke. Hvor tør du ringe
Menneske, du usle Orm, som selv skal blive Ormes
Føde, stole paa din Fornuft over for Vidnesbyrdene fra
de mange hellige og vise Mænd, Kirken kan opvise?
Tør du anse dig for klogere end Patriarker, Profeter,
Apostle, ja Kristus selv? Mener du, at du kan forklare
alt? Er det ikke djævelsk Hovmod, at du vil
bedømme Kristi Livs og Døds skrækkelige Mysterier,
som Kirken erklærer for ubegribelige? I saa mange
Aar har Kirken bestaaet fast og uforanderlig, stadfæstet
ved Mirakler og ved Martyrernes Heltemod, og
i Steden for at bøje dig for den, stoler du hellere paa
din Hjernes Fostre og gør dig elendigere end Dyrene,
saa taabelig opblæst du end er over din Aands og din
--- til
Læres Fortræffelighed! --- Dette Brev har jeg skrevet
dig i kristelig Hensigt, for at vise min Kærlighed til
dig, skønt du er en Hedning, og for at bede dig ikke
vedblive med at fordærve andre.

Spinoza vilde først ikke have svaret paa dette
Brev, fordi han mente, at Grunde her ikke vilde gøre
noget Indtryk. Men nogle fælles Venner, der lige som
han selv havde næret gode Forhaabninger om Burgh,
overtalte ham til at gøre et Forsøg. Han tilbageviser 
paa en værdig Maade den unge Mands Fanatisme.
% --- p. 62 ---
\setcounter{footnote}{0}%
\opage{62}--- Jeg bilder mig ikke ind, siger han, at have fundet
den bedste Filosofi, men jeg tror, at jeg har Indsigt i
den sande Filosofi. Der gives ingen Overbevisning saa
klar og sikker som den, der støtter sig til klare Grunde.
Det sande godtgør sig for Tanken og lærer os derved
tillige, hvor Usandheden ligger \textit{(verum est index sui
et falsi)}\footnote[1]{I et af sine Breve til Bleyenbergh siger Spinoza: Naar jeg har fundet et grundigt Bevis, kan jeg aldrig falde i saadanne Tanker, at jeg tvivler om det, jeg har bevist, og derfor slaar jeg mig til Ro ved det, Forstanden lærer mig, uden Mistanke om, at jeg skuffes. Og selv om jeg kommer til at se, at jeg har taget fejl, betragter jeg det dog som en Lykke, ti jeg gør derved et Skridt fremad, og det kan kun forøge den Glæde og Ro, hvormed jeg stræber at føre mit Liv.}. Men du, som paastaar, at du har fundet
den bedste Religion, hvilke Grunde kan du give derfor,
siden du ikke stoler paa Fornuften? Dersom du beraaber
dig paa Aandens indre Vidnesbyrd, saa kunne
jo alle andre gøre det samme. Dersom du henviser
til de udmærkede og hellige Mænd, der have levet i
Romerkirken, saa kunne alle Sekter blandt Protestanterne
hver for sig føre et lignende Bevis. Alle have
haft deres Martyrer. Der var jo nogle af din egen
Slægt, som under Hertugen af Albas Regimente led
alle mulige Pinsler for deres Tros Skyld. Jeg har selv
kendt en Jøde ved Navn Judas, der halvdød sang en
Psalme paa Inkvisitionens Baal. Peger du hen paa
Romerkirkens lange Fortid og sammenhængende Tradition,
saa overgaas den langt af Jødernes, der føre
deres Slægt og deres Traditioner lige ned til Adam,
og hvis religiøse Samfund bestaar urokket den Dag
% --- p. 63 ---
\setcounter{footnote}{0}%
\opage{63}i Dag, trods alle Forfølgelser i Fortid og Nutid. --- Du,
som en Gang tilbad et uendeligt Væsen, tilbeder nu
en Gud, som Chastillon lod sine Heste spise i Byen
Tienen\footnote[1]{Huguenotten Marechal Chastillon kastede indviede Hostier for sine Kavalleristers Heste.}. Erkend dog, at den sande Kristendom bestaar
i Retfærdighed og Kærlighed og hold op at kalde
absurde Vildfarelser for Mysterier! ---

Medens de hidtil nævnte Korrespondenter\footnote[2]{Endnu et Navn, og det et meget berømt, synes her at maatte nævnes. Hr. A. D. Jørgensen er nemlig ved at sysle med Steno's Biografi og Skrifter kommen til den Overbevisning, at Steno, som studerede i Leyden, medens Spinoza boede i det nærliggende Rhynsburg, og som selv beskæftigede sig med kartesiansk Filosofi, har staaet i nøje Forhold til Spinoza. Hans „\textit{Epistola ad novæ philosophiæ reformatorem}“, som er skrevet efter den berømte Anatoms og Geologs Overgang til Katolicismen, hvor han fandt Fred for sin grublende Aand, synes at være rettet mod Spinoza (særlig mod den „teol.-polit. Afhandling“). Han omtaler sin Modstander med stor Agtelse og hentyder til deres tidligere fortrolige Omgang. Skriftet findes des værre ikke paa vore Biblioteker.} væsentlig
betegne den gamle Verdensanskuelses Kamp mod
den nye Lære, som er den uforstaaelig, have vi endnu
to Grupper af Spinozas Korrespondenter tilbage, der
begge staa i et sympatetisk Forhold til ham; den ene
bestaar af hans egentlige Disciple, der antoge hans
Lære, den anden af dem, der gennemskue den saaledes,
at de formaa at rette en virkelig filosofisk Kritik
imod den og pege ud over den til nye Baner for den
filosofiske Forsken.

Vi have oven for omtalt Simon de Vries og hans
Forhold til Spinoza. I lignende smukt Pietetsforhold
% --- p. 64 ---
\setcounter{footnote}{0}%
\opage{64}stod flere unge Mænd, de fleste Læger og af jødisk
Herkomst, hørende til den lille Kreds, der studerede
hans Lære efter hans Manuskripter. Et Par smukke
Breve har man fundet til en af disse, Dr. med. Bresser.
I det ene beklager Spinoza sig over, at Bresser er
rejst uden at tage Afsked med ham. Han opfordrer
ham til Brevvexling og opmuntrer ham til at lægge
en fast Studieplan og hellige den bedste Del af sit Liv
til Uddannelse af sin Aand. „For at du kan fatte
Mod til at skrive med større Frihed til mig, maa du
vide, at jeg har den Formodning, at du nærer altfor
stor Mistillid til dine Evner og derfor frygter for at
komme med Spørgsmaal eller Udtalelser, der kunde
røbe Ukyndighed.“ --- Det andet Brev handler om den
rette Metode, man skal anvende i sit Studium. Spinoza
fremhæver, at --- afset fra den rent logiske Side ved
Sagen --- gælder det om stadigt Tankearbejde og Fasthed
i Aand og Vilje, hvilket ikke kan opnaas, uden
man fastslaar en bestemt Levevis og foreskriver sig et
bestemt Maal. --- I det Raad, han her giver sin Ven,
ligger der en Antydning af den Vej, han selv var gaaet.

Til den sidste Gruppe høre Leibnitz og Tschirnhausen.
Af disse har Leibnitz et berømt Navn i
Videnskabens Historie; Tschirnhausen har egentlig først
faaet Interesse for Filosofiens Historie, da man ved
det tidligere omtalte Fund af Spinozas Manuskripter
opdagede, at han var den hidtil ubekendte Forfatter
til de skarpsindigste af de til Spinoza rettede Breve.
Saa meget beslægtet disse Forskere end have med
Spinoza, saa gaa de dog ud over hans Standpunkt.
De betegne tillige Reaktionen mod ham, i det navnlig
% --- p. 65 ---
\setcounter{footnote}{0}%
\opage{65}Leibnitz igen lader Filosofien indgaa et Kompromis
med Teologien, der ingenlunde var til Fordel for hans
Tankes Klarhed og Grundighed. I det hele var Leibnitz
den fuldstændige Modsætning til Spinoza. I Steden
for dennes koncentrerede Ro og Fordybelse i én Tanke
træder hos hin en forbavsende Alsidighed: i Matematik
og Fysik, i Filosofi, Historie og Retsvidenskab er han
hjemme, og overalt griber han ind paa original Maade
og gør nye Opdagelser. Medens Spinoza holder sig
tilbage fra det offentlige Liv og tilbringer sin Tid med
at slibe sine Glas og dyrke sine Studier i det tarvelige
Kammer, kaster Leibnitz sig ind i en Hvirvel af praktiske
Forretninger, færdes ved Hofferne, filosoferer med
Dronninger og Prinsesser og teologiserer med høje
kirkelige Dignitarier. Til Fordel for hans Karakter
som Tænker og Menneske var denne Mangesidighed og
Udadvendthed ikke. Han gik let paa Akkord og tog
Hensyn, han ikke burde tage. Hans Adfærd over for
Spinoza viser ham særlig fra en lidet smuk Side. I
Spinozas Brevvexling findes kun ét Brev fra Leibnitz
til Spinoza og ét fra denne til hin, og begge handle
om optiske, ikke om filosofiske Spørgsmaal. Men af
de i sin Tid fundne Manuskripter faa vi ogsaa her nye
Oplysninger. Tschirnhausen, der traf Leibnitz i Paris,
skildrer ham som „fri for de sædvanlige teologiske
Fordomme“, og fortæller, at han satte stor Pris paa
Spinozas „teologisk-politiske Afhandling“; Leibnitz
havde endog, efter sit Udsagn til Tschirnhausen, skrevet
et bifaldende Brev til Spinoza om denne Bog; men
det er des værre gaaet tabt. Senere har han antydet,
at han en Stund har heldet til Spinozismen. Tschirn
% --- p. 66 ---
\setcounter{footnote}{0}%
\opage{66}hausen beder Spinoza om Tilladelse til at gøre Leibnitz
bekendt med Manuskriptet til Etiken. Men Spinoza
vil endnu ikke give sit Minde hertil. Vel kan han se
af Leibnitz's Breve, at han er en Mand med et frit
Blik og bevandret i alle Videnskaber; men han vil dog
vente med at meddele ham sine Skrifter, til han faar
nærmere Oplysninger om hans Karakter og om hans
Virksomhed i Paris\footnote[1]{Leibnitz var dengang i Kurfyrsten af Mainz's Tjeneste og var rejst til Paris med store politiske Planer.}. Den Mistillid, Spinoza saaledes
antyder, viste sig her --- lige som over for Albert de Burgh

--- ikke at være ubegrundet. Da Leibnitz efter sit
Besøg i Paris og London over Holland vendte tilbage
til Tyskland, besøgte han, som han selv fortæller, flere
Gange Spinoza og talte meget længe med ham. Men
officielt vilde han senere ikke ret vedkende sig dette
Bekendtskab. I sin \textit{Théodicée} fortæller han i Forbigaaende,
at han har besøgt Spinoza, men giver det
Udseende af, at deres Samtale drejede sig om „Anekdoter
fra den Tid“. Efter det Forhold, han dengang
stod i til Spinoza, og efter de store Problemer, hvormed
begge Mænd beskæftigede sig, at dømme, have de
sikkert talt om andet end Anekdoter. Senere betegnede
Leibnitz i et Brev Spinoza som „\textit{novateur trop
connu}“. --- Bekendtskabet vilde have været kompromitterende,
især for den, der vilde have Udseende af at
være ortodox og vilde forlige Tro og Viden. Men
hvad opnaaede Leibnitz ved al sin Leflen med Kirkelæren?
Man troede ham ikke; Folket i Hannover forandrede
hans Navn til „Løwenichts“ ɔ: Glaubenichts,
og da han døde, fulgte ingen gejstlig hans Baare.

% --- p. 67 ---
\setcounter{footnote}{0}%
\opage{67}Dog maa det ikke glemmes, at Leibnitz virkelig
udformede en filosofisk Anskuelse, hvis Styrke netop
laa, hvor Spinozas svage Side havde været. Spinoza
førte alt tilbage til den ene Substants; men Leibnitz
gik ud fra, at det ene og samme Univers opfattes
uendelig forskellig paa de forskellige Udviklingstrin,
paa hvilke de existerende Væsener befinde sig; der
er altsaa egentlig uendelig mange Verdener eller Substantser,
som han kaldte Monader. Denne Betragtningsmaade
er ikke ganske fremmed for Spinozas System,
men spiller dér kun en underordnet Rolle og
vilde ved bestemt at fremdrages føre til Strid med
hans Grundanskuelse. Mellemmanden mellem Spinoza
og Leibnitz, Tschirnhausen, har netop i sine Breve til
Spinoza henpeget paa dette Punkt.

Walther Ehrenfried von Tschirnhausen var en
saksisk Adelsmand, der i sin Ungdom rejste til Holland
for at studere. Han kom i Forbindelse med
Kredsen af Spinozas unge Venner og gennem dem
med Spinoza selv. Han er bekendt som Matematiker
og Fysiker (er bl. a. Opfinder af Brændglasset) og
stod som Filosof Spinoza og Leibnitz nær. Hans
Skrift \textit{Medicina mentis}, et Slags Logik, har saa
meget fælles med Spinozas Afhandling „om Forstandens
Forbedring“, at man endog har beskyldt ham for Plagiat.
I hvert Tilfælde ere hans vigtigste Lærdomme
en videre Udvikling af, hvad allerede Spinoza lærte,
og han bruger ofte endog de samme Ord og Vendinger.
Enkelte Steder hentyder han til Spinoza, men nævner
ham aldrig ved Navn (kun som \textit{qvidam} eller paa
lignende Maade), naturligvis af samme Grund som
% --- p. 68 ---
\setcounter{footnote}{0}%
\opage{68}den, der bestemte Leibnitz til saa stor Forsigtighed.
Tschirnhausen havde tilmed dediceret sit Skrift til
Ludvig XIV for at faa en Pension. --- Her interesserer
han os som Forfatter til de Breve i Spinozas Korrespondance,
der røbe det skarpeste Blik for Systemets
Grundbegreber og derved ogsaa for dets svage Sider.
Det er betydningsfuldt, at disse Breve ere skrevne paa
den Tid, da Tschirnhausen i Paris omgikkes Leibnitz;
de Punkter, de berøre, ere netop særlig komne til
deres Ret i Leibnitz's Filosofi.

Forhandlingen mellem Tschirnhausen og Spinoza
angaar to Hovedpunkter. Hvorledes kan man af Udstrækningens
almindelige Begreb udlede de mangfoldige
virkelig existerende Legemer med deres bestemte Figur
og Bevægelse? Ved dette Spørgsmaal rammes hele
den deduktive og abstrakt matematiske Metode, som
Descartes havde indført i Filosofien, og som Spinoza
havde bragt til sin Fuldendelse. I Udstrækningen i
sig selv ligger hverken Figur eller Bevægelse. Descartes
var endda, som Tschirnhausen bemærker, heldigere
stillet i dette Punkt; han saa, at Bevægelsen
ikke kunde udledes af den blotte Udstrækning, og
antog, at Gud ved Skabelsen havde bibragt Materien
den Bevægelse, den havde. Men Spinoza, der jo har
opgivet denne teologiske Forestilling, kan ikke forklare
den virkelige Bevægelse\footnote[1]{Denne Indvending mod Spinozas Lære fremførtes senere vidtløftigere i en skarpsindig Afhandling af John Toland i hans Letters to Serena (1704).}. Spinoza svarer hertil, at
Descartes's Begreb om Materien som hvilende Masse
var utilfredsstillende, men at han selv endnu ikke ret
% --- p. 69 ---
\setcounter{footnote}{0}%
\opage{69}havde gennemtænkt Sagen. Han udtaler sig i denne
Debat med en Mangel paa Sikkerhed, som maa forklares
ved, at han her stod ved sin Tankes historiske
Skranke. Et nyt Begreb var nødvendigt i Naturopfattelsen,
et Begreb, som Leibnitz særlig har fremdraget
baade i Fysiken og i Filosofien: Kraftens Begreb.
Spinoza har vel, som vi skulle se, sat dette Begreb
i nøjeste Forbindelse med sit Grundbegreb, Substantsen,
men det gennemtrænger endnu ikke hans
Naturopfattelse og viser sig ved nærmere Betragtning
at staa i Disharmoni med Substantsens Begreb, som
Spinoza opfatter det. --- Men hvad der stødte Tschirnhausen
i Spinozas Opfattelse af Materien, var ikke blot
den ensidige Bestemmelse af den som Udstrækning i
Rummet, hvorved det egentlig fysiske blev tilsidesat;
det var ogsaa den Fordring, at de materielle Fænomener
skulde udledes af ét eneste Begreb. For at
drage Slutninger, siger han, behøve vi stedse flere Begreber;
af ét enkelt kunne vi ikke udlede en Flerhed
af Bestemmelser. Naar jeg f. Ex. alene betragter
Cirkelens Omkreds, kan jeg ikke deraf slutte andet
end, at den overalt er sig selv lig og ensformig, hvorved
den adskiller sig fra andre Kurver; men andre
Egenskaber kan jeg ikke udlede. Naar jeg derimod
tillige tager Hensyn til andre Ting, nemlig til Radierne,
som drages fra Centrum, til to eller flere hverandre
skærende Linjer o. s. v.; saa kan jeg deraf udlede flere
Egenskaber. Spinoza svarer, at dette gælder i Matematiken
og Logiken, men ikke hvad reale Begreber
angaar, og derfor heller ikke for Substantsens Begreb;
men han angiver ikke, hvorfra denne Forskel skulde
% --- p. 70 ---
\setcounter{footnote}{0}%
\opage{70}skrive sig. Tschirnhausens Indvending er af stor Interesse,
fordi den peger hen paa Relativitetens Lov, der
mere og mere anerkendes i Filosofien. Enhver Fornemmelse,
enhver Forestilling og enhver Tanke staar i
Forbindelse med andre Fornemmelser, Forestillinger
og Tanker og er kun mulig i og ved en saadan Forbindelse.
Over for et i og for sig absolut og afsluttet
staar Tanken derfor magtesløs; den kan ikke tænke
det, fordi der intet er, hvormed den kan sammenligne
og maale det. Tschirnhausens Indvending, som kun
er rettet mod Materiens Begreb, rammer selve Grundbegrebet
hos Spinoza; ti hans „Substants“ er det, som
bestaar i sig selv og skal forstaas ved sig selv; men
alt hvad vi forstaa, forstaa vi stedse ved Hjælp af
noget andet. Nødvendigheden af at slaa ind paa en
anden Vej og benytte en anden Metode end Spinozas
bliver her indlysende.

Tschirnhausens andet Spørgsmaal angaar Spinozas
Lære, at Substantsen aabenbarer sig under uendelig
mange Attributer eller Grundformer, men at vi kun
kende to af dem (Aand og Materie). Spinozas Tankegang
var denne: Jo fuldkomnere et Væsen er, des
flere Grundformer maa det ogsaa have at aabenbare
sig under, og da Gud er det uendelige Væsen, maa
hans Væsen udtrykke sig i en Uendelighed af Grundformer.
Men, siger Tschirnhausen, da enhver Ting
hører med til Guds eller Substantsens Væsen, maa enhver
Ting udtrykkes under uendelig mange Grundformer,
og der bliver altsaa egentlig uendelig mange
Verdener. Og selv om det er én Substants, der bestaar
under de uendelig mange Attributer, ét Univers,
% --- p. 71 ---
\setcounter{footnote}{0}%
\opage{71}der opfattes paa uendelig mange Maader, saa kende
vi jo kun to af disse Maader, og de andre Grundformer 
blive altsaa som fremmede Verdener for os. --- Her
føres vi nær hen til den Ændring, Leibnitz, som nylig
antydet, indførte i Substantsens Begreb. Leibnitz's
Monader svare til den spinozistiske Substants's uendelig
mange Attributer, ti Verden forestilles i hver
Monade paa en ejendommelig Maade. Ved at sætte
Monaden i Steden for Substantsen har Leibnitz banet
Vejen for den kritiske Filosofi, der gaar ud fra, at
det, vi vide om Verden, er betinget ved vor Bevidstheds
Natur og Vilkaar og derfor kun gælder fra vort
Standpunkt. Spinoza gik dristig ud fra Ideen om det
absolute Væsen og hævdede en absolut Sandhed; han
saa endnu kun ufuldkomment, at „det sete afhænger
af Øjnene, der se“, og at vi kun kende Lyset som
brudt i en Mangfoldighed af Straaler. Det er en
Sandhed, som den engelske Filosofi, Leibnitz, Kant,
Sansefysiologien og Udviklingshypotesen hver paa sin
Maade have udtalt.

Saaledes se vi Spinoza i hans Brevvexling stillet
mellem Fortid og Fremtid. Vi se hans Tanke i dens
berettigede Kamp mod Fordommene, saa vel som i
dens Begrænsning, hvor den mødes med Fremtidens
Spirer.

% --- p. 72 ---
\setcounter{footnote}{0}%
\opage{72}\begin{center}Spinozas Filosofi.\end{center}


\begin{center}I.\end{center}


Spinoza begynder sit Hovedværk, Etiken, med en
Række Definitioner, der allerede giver os hele hans
System i kort Begreb. „Ved Aarsag til sig selv“ ---
saaledes lyder den første Definition --- „forstaar jeg
det, hvis Væsen medfører Existents, eller, med andre
Ord, det, der ikke kan begribes uden som existerende.“
Hertil slutter den tredje Definition sig: „Ved Substants
forstaar jeg det, der er i sig selv og begribes ved sig
selv, det vil sige, det, hvis Begreb ikke trænger til
nogen anden Tings Begreb, af hvilket det skal dannes.“
Det, som bestaar i og ved sig selv, maa ogsaa være
Aarsag til sig selv. Ti hvad der udspringer af noget
andet eller har sin Aarsag uden for sig selv, kan ikke
forstaas i og ved sig selv alene, og kan altsaa ikke
være Substants. Som Aarsag til sig selv er Substantsen
evig. Ti ved Evighed forstaar jeg, siger Spinoza,
selve Tingenes Existents, for saa vidt den fremgaar
af Tingenes Begreb. Substantsen er fremdeles
uendelig. Ti Endelighed er Begrænsning og Negation,
% --- p. 73 ---
\setcounter{footnote}{0}%
\opage{73}og hvis Substantsen var endelig, vilde dens Existents
bero paa noget andet end dens eget Væsen.

Substantsens Væsen og Indhold finder sit Udtryk
i Attributerne. Ved Attribut forstaar Spinoza, hvad
Tanken begriber om Substantsen som udgørende (konstituerende)
dens Væsen. Attributerne ere Definitioner
af Substantsen, ere hver for sig selve Substantsen fra
én bestemt Side. Attributet maa derfor lige som Substantsen
begribes i og ved sig selv: hvert Attribut for
sig er uendeligt og evigt, og det ene Attribut kan ikke
forklares ved Hjælp af det andet, det ene er ikke opstaaet
af det andet.

Medens Substantsen er det, der bestaar i sig selv
og forklares ved sig selv, forstaar Spinoza ved Maade
\textit{(modus)} det, der har sin Bestaaen i noget andet og
kun forstaas ved dette. Alle Maader ere derfor nærmere
Bestemmelser \textit{(affectiones)} ved Substantsen. Alt
værende falder altsaa i to Klasser: enten har det
Bestaaen i sig selv (Substants) eller i et andet
(modi), --- enten har det en ubetinget eller en betinget
Existents.

Det er de tre Grundbegreber hos Spinoza. Søge
vi danske Benævnelser for dem, kunne vi maaske kalde
Substants for Grundvæsen, Attribut for Grundform,
Modus for Enkeltvæsen. Enkeltvæsenerne have altsaa
deres Bestaaen i Grundvæsenet, og dette lægger sig for
Dagen i Grundformerne. Jo højere et Væsen staar,
des flere Grundformer har det. Det ene og uendelige
Grundvæsen, der bestaar under uendelig mange Grundformer,
kalde vi Gud. Gud existerer med Nødvendighed,
hans Tilværelse er grundet i hans Væsen. Eller
% --- p. 74 ---
\setcounter{footnote}{0}%
\opage{74}antag, at Gud ikke existerer: der maatte da være
noget, som hindrede og udelukkede hans Existents;
men hvor skulde dette noget være, da Gud er det
uendelige Grundvæsen, uden for hvilket der intet er?
Foruden Gud kan intet Grundvæsen gives eller tænkes.
Alt, hvad der existerer, har altsaa kun Existents som
enkelt Form for det ene Grundvæsen.

Som det, der begribes i og ved sig selv, er Grundvæsenet
Tankens Udgangspunkt, kun ved det forstaas
alt andet. Den rette Maade at filosofere paa er, siger
Spinoza, at gaa ud fra den højeste Aarsag. Men da
al Existents indbefattes i Grundvæsenet, har Spinozas
Filosofi ikke blot sit Udgangspunkt i dette Begreb;
den kan, naar den vil være konsekvent, ikke komme
ud over det. Dersom Spinoza antager en Existents,
der ligger uden for Grundvæsenet, modsiger han sig
selv: der vilde saa være noget uden for det ene
værende.

Af de uendelig mange Grundformer for det uendelige
Væsen kende vi kun to: Tænkning og Udstrækning.
Ved Udstrækning forstaar Spinoza, som vi vide,
hvad vi kalde Materie, i det han anser de geometriske
Egenskaber for de væsentlige. Udstrækningen er en
Form for Guds Væsen. Vel har man Ret i, at Gud
ikke er legemlig; ti et Legeme er en bestemt og begrænset
Figur. Men den udstrakte eller legemlige
Natur er ikke uden for Gud eller skabt af Gud. Grundformen
(Attributet) er jo selve Grundvæsenet (Substantsen),
evig og uendelig som dette. Man anser
det for uværdigt at gøre det legemlige til Udtryk for
Guds Væsen; men det er kun, fordi man stedse tænker
% --- p. 75 ---
\setcounter{footnote}{0}%
\opage{75}paa de endelige Legemer, der kunne deles og opløses.
Grundvæsenet er som saadant udeleligt og evigt.
Saaledes kan Vandet som Vand deles, og dets Dele
skilles fra hverandre, men ikke for saa vidt det er
Substants. --- Spinoza staar her som Forløber for Læren
om Materiens Bestaaen, som den nyere Kemi experimentalt
har bekræftet. De forskellige enkelte Stoffer
kunne opløses og omsættes til hverandre, men intet
Stof gaar til Grunde, intet Atom kommer til og intet
gaar tabt. --- Medens Spinoza fører et særligt Forsvar
for, at han tillægger Gud Udstrækningens Grundform,
gør han det ikke, hvad Tænkningens Grundform angaar,
da han her kan støtte sig til den almindelige Opfattelse
af Gud som Aand.

Vi kende altsaa i Følge Spinoza Gud som den uendelige
Kraft, der virker i Tankens saa vel som i Materiens
Verden. Men vi komme her naturlig til at spørge,
hvorfor vi kun kende disse to Grundformer blandt de
uendelig mange. Aarsagen ligger ikke fjern: vi kende
os selv som bevidste og materielle Væsener, og andre
Grundformer end Bevidsthed og Materie ere os ikke
tilgængelige i Erfaringens Verden. Dette fører os til
at se, hvorledes Spinoza i det hele er kommen til Begrebet
om Attribut eller Grundform. Han fastslaar
dette Begreb ligefrem ved en Definition uden at omtale
dets Oprindelse. Men det er egentlig vundet ved en
Analyse af det i Erfaringen givne. Alt, hvad vi kende,
kan henføres til to store Grupper: Bevidsthedsfænomener
og materielle Fænomener. Dette er den dybest liggende
Forskel, vi finde i hele Erfaringens Verden.
At Spinozas „Attribut“ egentlig kun er Karaktermærket
% --- p. 76 ---
\setcounter{footnote}{0}%
\opage{76}for en af disse Grupper, kan ses af, at han (lige i Begyndelsen
af „Etiken“) bruger Udtrykket „Art“ i samme
Betydning som „Attribut“. I samme Betydning bruger
han fremdeles ogsaa Udtrykket „Natur“, og i et tidligere
Skrift „Egenskab eller Kvalitet“ \textit{(proprietas sive qvalitas
sive attributum. Princ. phil. Cartes. I. def. 5)}.
Attributet er altsaa egentlig en Kvalitet, en Ejendommelighed,
der giver en vis Gruppe eller Klasse
Fænomener sit Grundpræg. Spinoza opstiller som 5te
Grundsætning i „Etikens“ anden Bog: „Vi fornemme
og iagttage ingen andre enkelte Ting end Legemer
og Tankeformer“; og støttende sig hertil beviser han
saa, at Udstrækning og Tænkning ere guddommelige
Attributer. Her ligger Begrebets virkelige Oprindelse
klart for Dagen.

Spinoza opfatter Attributet som Udtryk for Substantsens
Væsen. Hvorfra faar han da Ideen om Substantsen?
Den er hos ham umiddelbart given som det
i sig klareste og mest indlysende i vor Erkendelse.
Han bestemmer den derfor lige strax ved en Definition.
Dens Væren (Guds Existents) er en evig Sandhed.
Men se vi nøjere til, ville vi ogsaa her kunne spore
den virkelige Oprindelse. I sin Fremstilling af Descartes's
Filosofi siger Spinoza (I. def. 5): „Enhver Ting,
hvori eller hvorved noget, som vi opfatte, d. v. s. en
Egenskab, en Kvalitet eller et Attribut, har sin Bestaaen,
kaldes Substants“. „Det værende som saadant,
som Substants i og for sig, paavirker os ikke, og derfor
maa det forklares ved et Attribut, fra hvilket det
dog kun logisk kan sondres“ \textit{(Cogit. metaph. I, 3, 1)}.
Her se vi Sporene af den Metode, ved hvilken Erken
% --- p. 77 ---
\setcounter{footnote}{0}%
\opage{77}delsen i Virkeligheden stiger op ad for at naa det substantielle
i Tingene: Erfaring og Analyse af Erfaringer.\footnote[1]{Hos \emph{Descartes} udtales det endnu bestemtere, at Substantsens Begreb er naaet paa denne Maade. „Substantsen kan ikke oprindelig erkendes deraf alene, at den er en existerende Ting, ti det paavirker os ikke i og for sig alene; men vi lære den let at kende af et af dens Attributer. Deraf nemlig, at vi se, at et eller andet Attribut er til Stede, \emph{slutte} vi, at der nødvendigvis ogsaa maa være en existerende Ting eller en Substants til Stede, hvem dette Attribut kan tillægges“ (\textit{Princ. phil. 1. 52}). --- I sit Stridsskrift mod Descartes lægger \emph{Gassendi} med Rette stor Vægt paa, at vi naa Substantsens Begreb ved Slutning. Derfor er, siger han, dette Begreb saa abstrakt, utydeligt og problematisk, og det er ikke Egenskabernes (Accidentsernes) Erkendelse, der beror paa Substantsens Erkendelse, men omvendt Substantsens Erkendelse, der beror paa Erkendelsen af Egenskaberne. (\textit{Disqvisitio metaphysica. Amstelod. 1644 p. 121---123}).}


Men Erfaringen kan, siger Spinoza, kun lære os de
endelige Tings Existents at kende, ikke Grundvæsenet
og Grundformerne. De enkelte Ting, vi kunne lære af
Erfaringen, forstaa vi kun ret, naar vi i Forvejen kende
Tingenes indre Væsen.\footnote[2]{I en Note til den ufuldendte Afhandling „om Forstandens Forbedring“ siger Spinoza, at han har i Sinde udførlig at undersøge den empiriske Metode, som de nyere Filosofer anvende. Des værre har han ikke faaet udført denne Plan.} Han overspringer derfor Erfaringens
og Analysens Vej som unyttig og begynder
med at definere Grundvæsenets Ide, som han anser for
umiddelbart klar og indlysende. Det vilde for ham
være en Selvmodsigelse, at det, der i sig indbefatter al
Væren, skulde erkendes ved noget andet. Erfaringen
er for ham en Stige, der kun naar Halvvejen, og som
det saa meget des mindre falder ham ind at anvende,

% --- p. 78 ---
\setcounter{footnote}{0}%
\opage{78}som han fra Begyndelsen af føler sig i Højden. Naar
en Matematiker definerer, er hans Definition ét med en
Konstruktion, i det den angiver, hvorledes det definerede
bliver til. I Definitionen af Cirkelen som en Linje,
hvori ethvert Punkt har samme Afstand fra et vist
givet Punkt, ligger Anvisningen til at konstruere en
Cirkel ved at lade en ret Linje bevæge sig om sit ene
Endepunkt. Men en saadan Konstruktion have vi intet
tilsvarende til, hvad de abstrakte metafysiske Begreber
angaar. For Spinoza stod det som en Modsigelse, at
det i sig selv værende ikke skulde existere. „Det at
kunne existere er en Magt, og jo mere Realitet der
tilkommer en Tings Natur, des flere Kræfter har det i
sig til at existere: altsaa har det absolut uendelige
Væsen eller Gud absolut Magt i sig til at existere, og
existerer derfor absolut. Mange ville maaske ikke let
se det indlysende i dette Bevis, fordi de ere vante til
kun at betragte de Ting, der udspringe af ydre Aarsager.
Men jeg taler her kun om Grundvæsenet, der
ikke kan frembringes af nogen ydre Aarsag. Ti de
Ting, der opstaa ved ydre Aarsager, faa al den
Fuldkommenhed eller Realitet, de have, fra den ydre
Aarsags Kraft, og deres Existents udspringer alene af
den ydre Aarsags Fuldkommenhed, ej af deres egen.
Derimod, den Fuldkommenhed, Grundvæsenet besidder,
skyldes ingen ydre Aarsag, og derfor maa dets Existents
følge af dets Natur alene. Der er derfor intet,
om hvis Existents vi kunne være sikrere, end om det
absolut uendelige og fuldkomne Væsens Existents.“
Vi forstaa nu, hvorfor Spinoza begynder sit System
med Definitionen af Begrebet „Aarsag til sig selv“.
% --- p. 79 ---
\setcounter{footnote}{0}%
\opage{79}Vi forstaa fremdeles, at dette Begreb ikke kunde
begrundes ved Erfaringen, der stedse viser os det ene
betinget og bevirket ved det andet. Det er altsaa en
Fornuftidé, som Spinoza forudsætter, og lige som
enhver, der har Evne til at danne og anskue Figurer
i Rummet, kan tvinges til at anerkende Cirkelens Realitet,
saaledes mener Spinoza, at enhver, der gennemtænker
Begrebet „Aarsag til sig selv“, maa antage en
saadan Aarsags Existents.

Ud over og bag ved dette Begreb kunne vi hos
Spinoza ikke komme. Kun ved at tage et helt andet
Udgangspunkt og anvende en anden Metode kan man
komme til at sætte Begrebet om det i sig selv værende
under Debat. Det sker i den empiriske Filosofi i England
og i den kritiske Filosofi i Tyskland. Hvad der
for Spinoza stod som det sikreste af alt, Tankens første
og dybeste Udgangspunkt, bliver nu netop det mest
problematiske, det vanskeligste Begreb af alle. Kants
Begreb om „Tingen i og for sig“ (das Ding an sich)
minder ligefrem om Spinozas Definition af Substantsen;
men Kant betegner derved netop det, der stadig unddrager
sig vor Erkendelse, Grænsebegrebet i Tankens
Verden.

I sin „kortfattede Afhandling om Gud og Mennesket“
omtaler Spinoza (I. 7) den Indvending, at Gud
ej kan defineres, fordi enhver Definition sker ved at
henføre det definerede til en vis Art eller Klasse, men
Gud jo ikke kan henføres til nogen Klasse. Han indrømmer,
at dette er den sædvanlige Logiks Lære om
Definitionen, men betænker sig ikke paa at forkaste
den, hvad Grundvæsenet angaar. Substantsen, siger
% --- p. 80 ---
\setcounter{footnote}{0}%
\opage{80}han, er jo et ved sig selv bestaaende Væsen, som erkendes
i og ved sig selv. Det er kun de Ting, der
ikke bestaa ved sig selv, om hvilke det gælder, at de
defineres ved at henføres til en vis Klasse. Her hævder
altsaa Spinoza, at den Erkendelse, vi have af Substantsen,
er af en hel anden Art end vor øvrige Erkendelse.
Spørgsmaalet er da, om der virkelig gives
en saadan dobbelt Erkendelsesmaade, eller om vor
Erkendelses Art og Maade ikke er væsentlig den
samme helt igennem. Dette Spørgsmaal ville vi
komme tilbage til ved Behandlingen af Spinozas Erkendelseslære.

Som vi oven for have set, var allerede Tschirnhausen
stanset ved den Vanskelighed, at det ene Grundvæsen
skulde aabenbare sig under uendelig mange Grundformer,
hvoraf vi dog kun kende to. Der brydes her,
som Kuno Fischer klart har vist, to Bestræbelser hos
Spinoza. I Følge hans Grundtanke, at alt er ét, maa
det, der gælder for Substantsen, ogsaa gælde for Enkeltvæsenerne;
men mægtigere end denne Stræben efter
at opfatte det absolute som iboende er hans Forudsætning
om dets uendelige Ophøjethed over ethvert endeligt
Væsen. Denne Ophøjethed lægger sig for Dagen
ved, at medens vi som endelige Væsener kun have to
Attributer, har det uendelige Væsen uendelig mange.
Her have vi saaledes et Punkt, hvor det ses, med hvor
megen Ret jeg i Indledningen har hævdet, at Spinozas
spekulative Princip, Ideen om den absolute Substants,
kun ret forstaas, naar man gaar ud over Filosofien,
tilbage til den religiøse Bevidsthed. I en Anmærkning
til den nylig nævnte „kortfattede Afhandling“
% --- p. 81 ---
\setcounter{footnote}{0}%
\opage{81}(I, 1)\footnote[1]{Det maa dog bemærkes, at det ej er ganske sikkert, om Anmærkningerne til dette Skrift hidrøre fra Spinoza selv.} hedder det: „Efter at vi have betragtet Naturen,
have vi i den hidtil ikke kunnet finde mere end to Grundformer,
som tilhøre det alfuldkomne Væsen selv. Disse
ere os nu ikke nok; de kunne ikke være det eneste,
hvoraf dette fuldkomneste Væsen bestaar. Vi finde
tvært imod i os selv noget, som aabenbart kundgør os
ej blot flere, men uendelig mange fuldkomne Grundformer,
der tilkomme hint Væsen, for at det skal kunne
kaldes fuldkomment. Men dette noget kan ikke udspringe
af de to Grundformer; ti to give kun to og ikke uendelig
mange. Hvorfra da? Ikke fra mig selv; ti jeg
kan ikke give, hvad jeg ikke har. Hvorfra da ellers,
uden fra de uendelige Grundformer selv, der sige os,
at de ere, uden tillige at sige os, \emph{hvad} de ere; ti kun
om to alene vide vi, hvad de ere.“ Dette er den
eneste Antydning, der forekommer i Spinozas Skrifter
i denne Henseende, og den har næsten en mystisk
Karakter. Ogsaa denne Side af Sagen komme vi tilbage
til ved Spinozas Erkendelseslære. At der dog
ligger en stor Sandhed skjult i Spinozas Idé om Grundvæsenets
uendelig mange Grundformer, have vi allerede
oven for set. Netop naar vi indse vor Erkendelses Relativitet,
dens Afhængighed af vor Natur og af vor Stilling
i Tilværelsen, og altsaa give Afkald paa en absolut
Erkendelse, kunne vi ikke andet end indrømme Mulighednn
af uendelig mange andre Synspunkter for den
samme Tilværelse.

Skønt Spinoza ved at kalde Grundvæsenet Gud
% --- p. 82 ---
\setcounter{footnote}{0}%
\opage{82}antyder dette Begrebs Slægtskab med den religiøse
Tros Genstand, giver han dog nærmere Bestemmelser,
der fjerne de populære Forestillinger om Gud. Intet,
lærer han, kan være eller begribes uden Gud. Men
af den guddommelige Naturs Nødvendighed følger en
Uendelighed af Virkninger. Gud virker efter sin Naturs
Nødvendighed; der er intet uden for ham, som kunde
bestemme eller hindre hans Virken. Dette er den sande
Frihed, som er ét med den indre Nødvendighed. Gud
er den eneste frie Aarsag: de endelige Væsener have
alle noget uden for sig, som hindrer dem i at følge
deres Natur. Og da det er grundet i Guds Natur, at
uendelige Virkninger udspringe af den, saa maa disse
Virkninger være evige. Gud har fra Evighed af frembragt
alt, hvad han frembringer. Stod der noget tilbage,
som han endnu ikke havde faaet frembragt, vilde
hans Almagt ikke være evig. I Gud er der ikke Forskel
mellem Mulighed og Virkelighed. Hans Væsen
og Virken er en evig Virkelighed. Tid, Tal og Maal
finde ingen Anvendelse paa Grundvæsenet. I det
evige er der ingen Fortid eller Fremtid. Naar man
tilskriver Gud Forstand og Vilje, maa disse Ord tages
i en ganske anden Betydning, end naar de bruges om
Menneskene; Ligheden bliver kun en Navnelighed, som
den mellem Hunden, det gøende Dyr, og Stjernebilledet
af samme Navn. Særlig urigtigt vilde det være at tænke
sig Gud handlende efter Formaal; det vilde være at
underordne ham under en højere Magt, at sætte noget
uden for ham, hvorefter han rettede sig i sin Handlen;
altsaa det vilde være at ophæve selve Guds Begreb.
Ordet Personlighed, som Teologerne bruge til at betegne
% --- p. 83 ---
\setcounter{footnote}{0}%
\opage{83}Guds Væsen, er altfor ubestemt og uklart til at kunne
bruges med Nytte.

Det Spørgsmaal fremstiller sig her, hvorledes den
guddommelige Virken er at forstaa, naar al Tidsforskel
og al Forskel mellem Mulighed og Virkelighed skal
være udelukket. Oplysning herom ligger for det første
i selve Begrebet „Aarsag til sig selv“. Grundvæsenets
Virken frembringer ikke andet end det selv: dets Frembringen
af alt falder sammen med dets Frembringen af
sig selv. „Tingenes Væsen og Existents“, siger Spinoza,
„er en nødvendig Følge af den guddommelige Natur.
Gud maa kaldes Aarsag til alle Ting i samme Betydning,
i hvilken han kaldes Aarsag til sig selv. De endelige
Væsener ere ikke andet end Bestemmelser \textit{(affectiones)}
ved Guds Attributer, eller Former \textit{(modi)}, i hvilke Guds
Attributer udtrykkes paa en vis bestemt Maade.“ Lige
som vi ved Forstand baade forstaa Tankernes Kilde
og Indbegrebet af alle Tankerne, saaledes er Gud alle
Tings Kilde og alle Tings Indbegreb, Aarsag til alt
og Totaliteten af alt. Gud er med andre Ord iboende
(immanent, reflexiv) Aarsag, ikke overskridende (transscendent)
Aarsag; hvad han frembringer, er ikke noget
uden for ham, men en Form for hans eget Væsen.
--- Men for det andet maa det fremhæves, at Spinoza ofte
som Omskrivning for „Guddommens Virken“ bruger
Udtrykket „Følge eller Konsekvents af Guddommens
Natur“ og endog antyder, at denne sidste Betegnelse
er den rigtigste. I samme Retning viser den Sammenligning,
Spinoza saa ofte bruger, at alt følger af Guds
Væsen med samme Nødvendighed, som det følger af
Trianglens Væsen, at Summen af dens Vinkler er lig
% --- p. 84 ---
\setcounter{footnote}{0}%
\opage{84}to rette. Forholdet mellem Grund og Følge træder
altsaa i Steden for det mellem Aarsag og Virkning.
I Forholdet mellem Grund og Følge spiller Tiden
ingen Rolle; Følgen vorder ikke, men er evig ét med
Grunden. Spinozas Opfattelse bliver altsaa den, at alt
er som i et evigt Nu, fordi det har sin Grund i den
uendelige Substants, hvor ingen Tidsforskelle gælde.

Lige som Grundvæsenets Begreb ved konsekvent
Betragtning faldt sammen med Guds Begreb, saaledes
falder fremdeles Guds Begreb sammen med Naturens
Begreb. „Hint uendelige Væsen, som vi kalde Gud
eller Naturen, virker med samme Nødvendighed, med
hvilken det existerer.“ Men Spinoza skelner nu mellem
to Former for Guddommens eller Naturens Væsen: den
\emph{frembringende} Natur \textit{(natura naturans)}, Gud som
fri Aarsag, og den \emph{frembragte} Natur \textit{(natura naturata)},
Enkeltvæsenernes Totalitet, der hverken kan
være eller begribes uden Gud. Og inden for den frembragte
Natur skelnes igen mellem det, som \emph{umiddelbart}
udspringer af Guds Væsen, og det, som kun
\emph{middelbart} udspringer deraf. Hvad der umiddelbart
følger af Guds evige og uendelige Natur, maa selv
være evigt og uendeligt. Herved er at tænke paa
Naturordenen i sin Helhed, de i hele Universet gældende
Love saa vel inden for Aandens som inden for
Materiens Verden. Men det, som er endeligt og har
en begrænset Existents, kan ikke umiddelbart frembringes
af Guds absolute Natur. Det endelige kan
derfor kun frembringes af Gud, for saa vidt han fremtræder
i begrænset og endelig Form: det ene endelige
bestemmes altsaa af det andet, dette af et tredje, og
saaledes i det uendelige. Hvert Enkeltvæsen er Led
% --- p. 85 ---
\setcounter{footnote}{0}%
\opage{85}i Naturens Kæde og bestemmes ved de andre Led i
denne Kæde. --- Her indføres et helt andet Aarsagsforhold
end det, der gælder for Grundvæsenet. Enkeltvæsenerne
staa selvstændige uden for og over for hverandre,
virke som mekaniske Aarsager, medens Grundvæsenet
kun virker som iboende og reflexiv Aarsag,
aldrig paa andet end sig selv, aldrig mødende Modstand.

Spinoza siger i et af sine Breve, at Nødvendighedens
Begreb er Hovedgrundlaget for hans Lære.
Han nægter, at der gives noget tilfældigt. „Dersom
Menneskene klart forstod Naturens hele Orden, vilde
de finde alt lige saa nødvendigt som alle Matematikens
Sætninger.“ Men Nødvendigheden fremtræder altsaa i
hans System under en dobbelt Form: som indre Nødvendighed
i Grundvæsenets Natur og Virken, og som
ydre Nødvendighed i Enkeltvæsenernes Vexelvirken og
gensidige Bestemmelser. Da nu Grundvæsenet er Spinozas
første og sidste Tanke, maatte han konsekvent
være stanset ved den indre Nødvendighed; men saa
vilde Enkeltvæsenernes Verden, den konkrete Virkelighed,
ingen Realitet have faaet. Enhver enkelt og
endelig Ting er i Følge Spinoza en Begrænsning og en
Bestemmelse af Substantsens Væsen; og, føjer han til,
enhver Bestemmelse er en Nægtelse. Bestemmelsen
hører nemlig ikke nødvendig med til Tingens Natur:
der, hvor Grænsen sættes, hører jo Tingen op at være.
Enhver bestemt Figur er saaledes en Nægtelse, en
Ophævelse af den uendelige Udstrækning. At sætte
en enkelt endelig Ting som existerende er at ophæve
Guds Existents. Vel lærer Spinoza udtrykkelig, at
uagtet enhver enkelt Ting bestemmes af en anden Ting
% --- p. 86 ---
\setcounter{footnote}{0}%
\opage{86}til at existere paa en vis Maade, følger dog den Kraft,
hvormed den bestaar i sin Existents, af Guds Naturs
evige Nødvendighed. Men denne Kraft maa dog stedse
selv være en Bestemmelse og altsaa efter Spinozas
Opfattelse en Nægtelse; den maa forholde sig til den
uendelige Guddomskraft, som Figuren forholder sig til
den hele Materie, og af Substantsens Væsen kan man
ikke udlede det, der ophæver Substantsen. Der kan i
Spinozas System egentlig ikke blive Plads for en ydre
Nødvendighed, en Vexelvirkning mellem reale Enkeltvæsener.
Naar disse hvert for sig ere Negationer,
kunne de ikke paavirke og bestemme hverandre; ti af
intet kommer intet. Naar Spinoza nu alligevel lærer
en saadan ydre Nødvendighed og fremstiller Naturen
som en Verden af Elementer, der forme sig til større
og mindre Helheder efter Lovene for deres Vexelvirkning,
saa er denne Lære ikke vunden ad deduktiv
Vej. Den hviler paa en Forudsætning, som ikke er
indeholdt i Grundvæsenets Idé, og han har selv udtalt
i et af sine Breve, at det er Erfaringen, han her bygger
paa. Til Dobbeltheden inden for hans System svarer
altsaa en Dobbelthed i hans Metode: Grundvæsenet og
Enkeltvæsenerne forholde sig til hverandre som Spekulation
og Erfaring. Spinoza lægger Spekulationen til
Grund og bruger kun Erfaringen, hvor Spekulationen
ikke slaar til. Her træder han i en tydelig Modsætning
til Nutidens Filosoferen, som i Regelen lægger
Erfaringen til Grund og lader Spekulationen bygge sine
Hypoteser paa Erfaringserkendelsens Grundlag. Dette
er en saa meget mere naturlig Fremgangsmaade, som
det kun herved bliver muligt at kontrollere selve de
% --- p. 87 ---
\setcounter{footnote}{0}%
\opage{87}Principer, hvorfra Spekulationen gaar ud. Men det
viser Spinozas Dybsindighed, at selv en saaledes forandret
Metode fører Tænkningen til en Verdensanskuelse,
der er beslægtet med hans. Saaledes ville Arbejdere,
der grave en Tunnel gennem et Bjerg hver fra
sin Side, mødes i Midten, naar de kun gaa den rette
Vej, hver fra sit Udgangspunkt.

Jakobi ansaa Spinoza for den eneste konsekvente
Filosof og gav Valget mellem hans Lære og et \textit{salto
mortale} over i Troen. Senere opdagede Kritiken til
manges Glæde den dobbelte Strømning og for saa vidt
Inkonsekventsen hos Spinoza. Jernringen syntes brudt.
Men det maa vel bemærkes, at den dobbelte Strømning
ikke først fremkommer paa et vist Punkt inde i Systemet,
men er til Stede i det første Anlæg, i de Definitioner
og Grundsætninger, hvormed det begynder.
Naar Spinoza ikke blot begynder med Definitionen af
„Aarsag til sig selv“ og af „Grundvæsen“, men i samme
Aandedræt definerer „Enkeltvæsenet“, og naar hans
første Grundsætning er: alt, hvad der er, er \emph{enten} i
sig selv \emph{eller} i et andet, --- saa begynder han i Virkeligheden
ikke med én, men med to Forudsætninger,
øser fra først af fra to forskellige Kilder. Og Dobbeltheden
er i Virkeligheden kun paa Overfladen; de to
Strømninger lyde en og samme Grundlov: det er Ideen
om den lovbestemte Sammenhæng, der behersker alt
hos Spinoza, og som godtgør en dybere Enhed bag den
Dobbelthed, Kritiken kan paapege.

Den Sætning, at enhver Bestemmelse er en Nægtelse
\textit{(omnis determinatio est negatio)}, er ofte og med
Rette anført som karakteristisk for Spinozas Standpunkt.
% --- p. 88 ---
\setcounter{footnote}{0}%
\opage{88}For ham var det positive det i sig bestaaende, der
er hævet over alle Forhold og alle Skranker. Men
denne Opfattelse rammes af Goethes bekendte Ord:
„In der Begrenzung zeigt sich erst der Meister.“ Det
er netop vor Opgave at vinde Bestemmelser, at gribe
Nuancer, at bestemme Forhold. Vi finde Tænkningens
Positivitet i det rige, mangfoldige og artikulerede Indhold,
den formaar at vinde. Først ved Bestemmelsen
bliver Indholdet positivt. Enhver Bestemmelse, kunne
vi sige, er en Relation; hver Gang vi tænke noget,
bestemme vi dets Relationer og lære det først derved
at kende \textit{(omnis determinatio est relatio)}. Vor Erkendelse
er relativ, og kun derved er den ogsaa positiv.
Man vil ogsaa let se, at det absolute og uendelige, det
i sig bestaaende, kun kan tænkes og betegnes paa
negativ Maade, ved Negation af de endelige Bestemmelser.
Det i \emph{sig} bestaaende kunne vi ikke forklare
anderledes end som det, der \emph{ikke} behøver \emph{noget}
\emph{andet} for at bestaa. Hvad Spinoza ansaa for det
sande, eneste positive, er altsaa netop negativt, er et
Grænsebegreb, der kun tænkes ved Hjælp af de Bestemmelser,
som i det negeres.

% --- p. 89 ---
\setcounter{footnote}{0}%
\opage{89}\begin{center}II.\end{center}


Den ejendommelige og konsekvente Gennemførelse
af Substantsens Begreb er det første og mest epokegørende
Træk i Spinozas Filosofi. Men hertil slutter
sig nøje et andet Hovedpunkt. Naar ét Grundvæsen
aabenbarer sig i alt, kan Forholdet mellem Aand og
Materie ikke være saaledes, som det almindelig forestilles.
Spinoza fastholder, lige som Descartes, at Forskellen
mellem Tankernes indre Verden og Udstrækningen
i Rummet for os er uovervindelig; Tanke og
Udstrækning ere for ham to Grundformer, hvori Grundvæsenet
aabenbarer sig, helt og fuldt i enhver af dem.
Af Grundformens Definition følger det da ligefrem,
hvorledes Forholdet mellem Aand og Materie maa
være. Ved Grundform eller Attribut forstaar Spinoza
jo, hvad Tanken udsiger om Substantsen som udgørende
dens Væsen. Men som Udtryk for Substantsen
eller Grundvæsenet maa enhver Grundform begribes i
og ved sig selv; Grundformen er jo kun Grundvæsenet
selv, set fra en vis Side. Den ene Grundform kan
derfor ikke være frembragt af den anden: Aanden ikke
af Materien, Materien ikke af Aanden; alle Grundformer
ere oprindelige og samtidige. Et og samme
Grundvæsen aabenbarer sig paa én Gang under forskellige
Former, af hvilke den ene ikke udelukker den
% --- p. 90 ---
\setcounter{footnote}{0}%
\opage{90}anden. Hver enkelt Ytringsmaade af en Grundform maa
derfor forstaas og forklares ved andre Ytringsmaader af
samme Grundform og ikke ved Bestemmelser, som høre
ind under en anden Grundform: Tanker kunne kun forklares
ved Tanker, materielle Bevægelser kun ved materielle
Bevægelser. Aarsagen maa overalt være ensartet
med Virkningen. Men da det er et og samme
Grundvæsen, der aabenbarer sig i de to Grundformer,
vil der være en Parallelisme til Stede: Tankernes Orden
og Forbindelse er den samme som Tingenes Orden og
Forbindelse. Guds Tankekraft og Handlekraft svare
til hinanden, og alt, hvad der følger af Guds uendelige
Natur inden for Tankens Verden, følger i samme Orden
og samme Forbindelse inden for Materiens Verden.
Hvad enten vi derfor opfatte Naturen som Materie eller
som Aand, saa er det en og samme Aarsagsrække, vi
have for os. Og hvis andre Grundformer vare os tilgængelige,
vilde det forholde sig paa samme Maade
med dem; hver for sig vilde være et i sig sammenhængende,
i en uendelig Række af Fænomener aabenbaret
Udtryk for det ene Grundvæsen.

Det er karakteristisk for Spinozas Standpunkt og
Fremgangsmaade, at han begynder med at bestemme
Forholdet mellem Aand og Materie i dets Almindelighed,
og først derpaa gaar over til Forholdet mellem Sjæl og
Legeme. Dette Forhold er et særeget Exempel paa
hint almindelige Forhold. --- Opfattet under Aandens
Grundform lægger Grundvæsenet sig for Dagen i en
Verden af Tanker. Den menneskelige Aand er en af
disse uendelig mange Tanker. Naar vi altsaa sige, at
den menneskelige Aand opfatter dette eller hint, betyder
det, at Grundvæsenet, ikke for saa vidt det er
% --- p. 91 ---
\setcounter{footnote}{0}%
\opage{91}uendeligt, men for saa vidt det finder sit Udtryk i den
menneskelige Aands Natur eller udgør (konstituerer) den
menneskelige Aands Væsen, har denne eller hin Tanke.
Til hver enkelt Tanke svarer der inden for Udstrækningens
Grundform en materiel Bevægelse. De forskellige
Bevidstheder og de forskellige Tanker afvige
fra hverandre i samme Grad, som de tilsvarende materielle
Bevægelser og Legemer afvige fra hverandre.
Jo flere Paavirkninger Legemet er i Stand til at modtage,
des flere Ideer er ogsaa Bevidstheden skikket til
at opfatte; og jo mere selvstændigt og aktivt Legemet
er, des klarere og tydeligere formaar Aanden at tænke.
Da Grundformerne hver for sig ere uendelige, maa
dette gælde for alt virkeligt: alt i Naturen er altsaa
besjælet, skønt i forskellige Grader. Ethvert existerende
Væsen har --- som Led i Grundvæsenet --- en
dobbelt Side, hvorfra det kan sés, er baade aandeligt
og materielt. --- Da det her bliver af Interesse at kende
Legemernes forskellige Natur, indskyder Spinoza nogle
Sætninger af Naturlæren.

Legemer adskille sig kun fra hinanden ved det
forskellige Forhold mellem Bevægelse og Hvile. Ethvert
Legeme bestemmes til Bevægelse eller Hvile ved
et andet Legeme, dette af et tredje, og saaledes fremdeles.
Den Maade, hvorpaa et Legeme bestemmes,
beror dels paa dets egen Natur, dels paa det Legemes
Natur, som paavirker det. Et og samme Legeme bevæges
forskellig efter de paavirkede Legemers Forskellighed,
og forskellige Legemer bevæges paa forskellig
Maade af et og samme Legeme.

Af de simpleste, mest enkelte Legemer danner der
% --- p. 92 ---
\setcounter{footnote}{0}%
\opage{92}sig sammensatte Legemer. Naar nemlig Legemer af
samme eller forskellig Størrelse holdes saaledes sammen
af andre Legemer, at de føje sig sammen indbyrdes,
eller naar de indbyrdes meddele hverandre deres Bevægelser
i et vist bestemt Forhold, siges de at være
indbyrdes forbundne og danne alle til sammen ét Legeme
eller ét Individ, som adskilles fra andre Individer
ved den bestemte Forbindelse mellem de enkelte Legemer,
hvoraf det bestaar. Et saadant Individ vedbliver
at være det samme, selv om nogle Dele udskilles og
andre optages, og selv om Delenes Størrelse og Bevægelsernes
Retning forandres, naar kun den hele
Form for Sammensætningen vedbliver at være den
samme.

Flere saadanne enkelte Individer kunne nu tænkes
forenede til et Individ af anden Orden; flere af denne
sidste Art igen til et Individ af tredje Orden, og saaledes
fremdeles. Til sidst opfatte vi da hele Naturen
som et eneste stort Individ, hvis enkelte Dele vexle i
det uendelige, uden at det hele Individ undergaar nogen
Forandring.

Den Maade, hvorpaa Spinoza her udvikler Individets
Begreb, er ganske i den nyere Naturvidenskabs
Aand. I Steden for at søge indre, mystiske Enheder
søger han den Lov, som gælder for Individets Tilbliven.
Kun derved kan Individualitetens Begreb fastholdes
over for den Vanskelighed, at Erfaringen ingen absolute
og afstukne Grænser viser os mellem de enkelte
Væsener. Det er en nødvendig Konsekvents af den
nyere Naturforsken, at man enten helt maa opgive
Individualitetens Begreb eller ogsaa gøre det relativt,
% --- p. 93 ---
\setcounter{footnote}{0}%
\opage{93}bruge det i et videre Omfang, saa at vi, lige som
Spinoza, faa Individer af forskellige Grader.

I det Spinoza lærer en opadstigende Række af
Individualitetsgrader i Naturen, har han tillige udtalt
Udviklingens Lov i en af dens mest karakteristiske
Former: Overgang fra mindre til større Sammenhæng.
Hans Fremstilling er en Forløber for Herbert Spencers
\textit{First Principles.}\footnote[1]{Se mit Skrift: „Den engelske Filosofi i vor Tid“. Kbhvn. 1874. P. 136 ff.}%


Erindre vi nu, at der til hver materiel Form, hvert
legemligt Individ svarer en ideel Form, en individuel
Tanke, i det alt er besjælet, skønt i forskellige Grader\footnote[2]{Man vilde aldeles misforstaa Spinozas Lære om Naturens almindelige Besjæling, dersom man (som f. Ex. \emph{John Toland} gør i den tidligere nævnte Afhandling) mente, at den skal være et Forsøg paa at forklare, hvorledes Bevægelsen opstaar i Materien. Spinoza nægter jo netop paa det bestemteste, at Tanke og Materie kunne virke paa hinanden. Hans Tankegang er derimod den, at Materien er ensartet i alle sine Former og Sammensætninger; naar vi altsaa paa et vist Punkt finde Bevidsthedsfænomener forbundne med materielle Bevægelser, saa maa alle materielle Bevægelser, skønt i meget forskellig Grad, være ledsagede af saadanne. (\textit{Et. II, 13. Schol.}) Denne Teori er opstaaet hos Spinoza ved den resolute Maade, hvorpaa han gennemfører den materielle Verdens Enhed og Sammenhæng. Den staar ogsaa nu som en naturlig Hypotese, og er i vor Tid udtalt ej blot af Filosofer, men ogsaa af Naturforskere som Zøllner, Haeckel o. fl.}, ---
saa fremtræder Spinoza som Forgænger for Leibnitz's
Monadelære. Lige som det menneskelige Legeme
bestaar af mange sammensatte Individer af forskellig
Natur, saaledes er ogsaa den Tanke, der udgør den
menneskelige Aands Natur, sammensat af mange Tanker.

% --- p. 94 ---
\setcounter{footnote}{0}%
\opage{94}Hele Naturen bestaar af en Uendelighed af individuelle
Væsener, hvert med sin Grad af Forestillingsevne.
Og saadanne Væsener er det netop, Leibnitz kalder
Monader\footnote[1]{Fra en anden Side viser --- som tidligere fremhævet --- Spinoza fremad til Leibnitz ved Læren om de uendelig mange Grundformer for det ene Grundvæsen.}.

Ogsaa i dette Punkt er Spinozas Anskuelse afgjort
Monisme. Lige saa lidt som han kunde indrømme et
dualistisk Forhold mellem Grundvæsenet og Enkeltvæsenerne,
lige saa lidt kunde han det mellem de
aandelige og de materielle Fænomener. Han søger
overalt Enheden, som ligger bag Forskellighederne.
Og denne hans Stræben træder især stærkt frem, naar
man sammenligner ham med hans Forgænger i Filosofien,
Descartes. I den kartesianske Filosofi vare Sjæl
og Legeme to helt forskellige Substantser, som egentlig
kun ved et Mirakel kunde forbindes med hinanden.
Skønt Descartes, der med Iver havde dyrket Anatomien,
havde en for hans Tid saa klar og træffende Opfattelse
af Nervesystemets Betydning, at den nyere Nervefysiologi
i ham har sin Forgænger\footnote[2]{Han var, i Følge \emph{Huxley}, for Nervesystemet, hvad Harvey var for Blodomløbet.}, var det dog paa rent
udvortes Maade, han tænkte sig Sjælen forbunden med
Legemet. I den saa kaldte „Koglekirtel“ i Hjernen
saa han det Sted\footnote[3]{At „Koglekirtlen“ ikke har en saadan Betydning, havde \emph{Steno} allerede godtgjort i en Afhandling 1669.}, hvor Sjælen virker paa det legemlige
og atter paavirkes af dette. Spinoza bemærker
i sin Kritik af denne Lære, at det aandelige og det
materielle ikke have noget fælles Maal, ikke kunne
% --- p. 95 ---
\setcounter{footnote}{0}%
\opage{95}sammenlignes indbyrdes og derfor heller ikke bestemme
hinanden. Hverken kan Legemet bestemme Sjælen til
at tænke, ej heller Sjælen Legemet til Bevægelse eller
Hvile eller noget som helst andet. Hvad der bestemmer
Sjælen til Tænkning, maa selv være en Tanke, og hvad
der sætter Legemet i Bevægelse eller Hvile, maa selv
være et Legeme. Sjæl og Legeme ere et og det
samme, som snart betragtes under Tænkningens, snart 
under Udstrækningens Grundform. Vi have ikke
--- som hos Descartes --- først en legemlig Bevægelse, der
sætter Sjælen i Bevægelse, og saa en sjælelig Virksomhed;
men da den legemlige Bevægelse i sig selv
(i Grundvæsenet) er et med den sjælelige Akt, falder
Tidsforskellen bort.

Spinoza føler, at det deduktive Bevis (ud fra Grundformens
Begreb) her ikke vil være nok. „Folk ere“,
siger han, „fast overbeviste om, at Legemet bevæges
eller hviler paa Sjælens blotte Vink, og at det udfører
meget, som beror paa Sjælens Vilje og Tankekløgt.
Derfor vil man næppe overveje Sagen med Ro, naar
jeg ikke henter Beviser fra Erfaringen.“ Han bemærker
da først, at ingen kan vide, hvor langt Legemets Kraft
naar. Ingen har endnu fundet Grænsen for, hvad der
kan bevirkes i Følge den materielle Naturs Love. Det
er altsaa ganske vilkaarligt, naar man her vil sætte bestemte
Skranker, ud over hvilke intet skulde være muligt
uden ved Sjælens Indgriben. --- Men hvad mener
man i Grunden, naar man taler om, at Sjælen sætter
Legemet i Bevægelse? Ingen kan danne sig et klart
og tydeligt Begreb derom. At sige, at Sjælen bevirker
en legemlig Bevægelse, er altsaa det samme som at
% --- p. 96 ---
\setcounter{footnote}{0}%
\opage{96}sige, at man ikke véd, hvorledes denne Bevægelse
bliver til. Man pynter her, som saa ofte, sin Uvidenhed
med velklingende Ord. --- I Dyrenes instinktmæssige
Handlinger, i Søvngængeres og Somnambulers Færd
have vi Tilfælde, hvor den legemlige Natur virker alene,
og hvor den udretter Ting, som alle beundre. --- Fremdeles
viser Erfaringen, at Sjælens Livlighed og Virksomhed
staar i et bestemt Forhold til den legemlige
Tilstand. Det vilde visselig staa langt bedre til, hvis
Menneskene havde Evne til at styre deres Legemer.
Men de have jo ikke engang Tungen i deres Magt.
Og dog tror den snakkesalige, at han handler fuldkommen
frit, hvor tit han end kommer til at fortryde sin
Aabenmundethed. Da Mennesket ikke kender de Aarsager,
der bestemme hans Handlen, anser han sig for
absolut fri, lige som Stenen, naar den kunde tænke,
vilde mene, at den faldt til Jorden med fri Vilje. --- Vi
kunne jo ikke handle efter vor Beslutning, uden vi
erindre, hvad den gaar ud paa. Men Erindringen ere
vi ikke altid Herre over. Enhver maa indrømme, at
vi ikke i ethvert Øjeblik kunne fremkalde for vor
Tanke, hvad det skal være. Tanker og Beslutninger
opstaa i Bevidstheden med samme Nødvendighed som
Forestillingerne om de ydre Genstande. Og den nødvendige
Lov, der gælder for Bevidstheden, svarer til
den, der gælder for den tilsvarende legemlige Proces:
„Aandens Beslutning eller Attraa og den legemlige
Bestemmelse ere i Følge deres Natur samtidige, eller
rettere, det er en og samme Ting, som vi kalde Beslutning,
naar den betragtes under Tænkningens Attribut
og forklares ved det, og legemlig Bestemmelse, naar
% --- p. 97 ---
\setcounter{footnote}{0}%
\opage{97}den betragtes under Udstrækningens Attribut og udledes
af Lovene for Bevægelse og Hvile.“

Hele denne Lære hører til det, der sikrer Spinozas
Skrifter en blivende Betydning. Han har her paa genial
Maade foregrebet en Erkendelse, til hvis erfaringsmæssige
Gennemførelse først vor Tid er naaet. Nogen
Grund havde han at bygge paa, nemlig Descartes's
mekaniske Nervefysiologi. Men Descartes blev her,
som paa flere Punkter, staaende paa Halvvejen. Spinozas
Blik for Enhed og Sammenhæng førte ham til
den Overbevisning, at de legemlige Aarsagers Række
ikke paa noget Punkt kan tænkes afbrudt. Naar det
nu alligevel staar fast, at der er en (for os) uovervindelig
Forskellighed mellem Bevidsthed og Materie
(eller i Spinozas Sprog: Tænkning og Udstrækning),
saa bliver den eneste mulige Opfattelse den, at tænke
sig disse to Tilværelsesformer som samtidige og parallele
Udtryk for et og samme Grundvæsen. Til en
lignende Opfattelse synes ogsaa vor Tids Fysiologi og
Psykologi at føre\footnote[1]{Dette haaber jeg ad Aare at kunne godtgøre i et særligt Skrift.}%
.

For Spinoza som spekulativ Filosof stillede det sig
saaledes, at Grundvæsenet, der ligger bag ved Aand
og Materie, er det for Tanken klareste og sikreste.
For os svimler det paa denne Højde. Vi glæde os ved at
se Traadene nærme sig hverandre, jo mere vi formaa at
følge dem, ved at se selv de for os uovervindelige Modsætninger
i vore Erfaringers Kreds vise tilbage til en
dybere Enhed; men selve denne Enhed formaa vi ej
% --- p. 98 ---
\setcounter{footnote}{0}%
\opage{98}at fremdrage; hver Gang Tanken søger at holde den
fast, viser det sig, at vi ere i Færd med at forlade
vor Erkendelses egentlige Omraade. Fra de spekulative
Hypotesers vildsomme Regioner søge vi da tilbage
til, hvad Kant har kaldt „Erfaringens frugtbare Dyb“.
Vi undre os derfor ikke over, at vi ved Spinozas Stræben
efter at bygge Filosofien paa det absolute Væsens
Idé føres til saadanne Vanskeligheder som den, hvorledes
Grundvæsenet, sin Enhed ubeskaaren, kan fremtræde
i en Flerhed af Grundformer, og hvilken Betydning
disse have for det og for hverandre indbyrdes.
Livets og Tilværelsens Gaade møder os stedse paa ny
ved Enden af alle de Forsøg, den menneskelige Aand
i sine religiøse og spekulative Systemer har gjort paa
at løse den.

Spinoza er den første, som har udviklet en Teori,
der hjælper os ud over Spiritualismens og Materialismens
Ensidigheder, ikke ved at gaa paa Akkord, men ved
virkelig at tænke Sagen til Ende og modig tage alle
Konsekventser. Men som man ikke forstod hans religiøse
Teori og udlagde hans Panteisme som Ateisme,
saaledes forstod man ikke hans psykofysiske Teori, men
udlagde hans Monisme som Materialisme. Og selv om
man ikke gaar saa vidt, fremhæves det dog ofte som
en Inkonsekvents hos Spinoza, at han ser Sagen vel
meget fra den legemlige Side, naar han f. Ex. lærer,
at Bevidsthedens Fuldkommenhed beror paa det tilsvarende
Legemes Fuldkommenhed. Men det ligger i
Sagens Natur, at vi ere nødte til at begynde med det
legemlige, naar vi ville sammenligne de forskellige
% --- p. 99 ---
\setcounter{footnote}{0}%
\opage{99}Trin af sjæleligt Liv. Med Undtagelse af vor egen
Bevidsthed kende vi kun Sjælelivet gennem det legemlige
som Mellemled og finde det overalt knyttet til dette.
Og hvad der er Hovedsagen: vor Erkendelse af det
legemlige kan naa en Fuldkommenhed og Nøjagtighed, 
som Erkendelsen af det aandelige aldrig vil naa. Det
sikreste Udgangspunkt har man derfor, naar man begynder
med det legemlige. At Sagen dog efter Spinozas Opfattelse
lige saa godt kunde ses fra den anden Side,
viser en Sætning som den (Et. V, 1): „Lige som
Tankerne og Forestillingerne ordnes og forbindes i Bevidstheden,
nøjagtig saaledes ordnes og forbindes Legemets
Tilstande i Legemet.“ Der kan ikke ske nogen
Tankeforbindelse, uden at der tillige sker en Forbindelse
af legemlige Bevægelser (det moderne Udtryk vilde
være: Molekularbevægelser i Hjernen). Derfor er det
heller intet Brud paa Teorien, naar Spinoza viser,
hvorledes Mennesket kan opnaa Herredømme over sine
Drifter og aandelig Frihed. Til det højere og det
lavere inden for Bevidsthedslivet svare naturligvis forskellige
Kombinationer af legemlige Processer. Den
nyere Fysiologi viser os i Forholdet mellem højere og
lavere Nervecentra, i den Maade, hvorpaa hine kunne
regulere disse, en Parallel til den psykologisk-etiske
Modsætning mellem Tanke og Drift („Aanden“ og
„Kødet“). Man vilde ikke saa ofte have fremført den
omtalte Indvending mod Spinoza, naar man havde
erindret hans Lære om Grundformerne som hver for
sig uendelige, saa at der gives legemlige Former, der
svare til de højeste aandelige Funktioner, saa vel som
% --- p. 100 ---
\setcounter{footnote}{0}%
\opage{100}aandelige Former, der svare til de simpleste legemlige
Bevægelser, og naar man tillige havde betragtet Sagen
fra et fysiologisk Synspunkt. Spinoza henviser udtrykkelig
til de uvilkaarlige Handlinger (Instinkthandlinger,
Somnambulisme etc.) og finder, lige som den
nyere Psykologi, i dem et Vink til Forstaaelsen af de
vilkaarlige Bevægelser. Det er Sammenhængens, Kontinuitetens
Ide, der ogsaa her leder hans Tanke til et
lykkeligt Greb.

Inden vi forlade dette Punkt i Spinozas Lære, maa
det fremdrages, at han her ikke altid udtaler sig klart
og konsekvent. Som Exempel paa, at en vis materiel
Form og den tilsvarende Tanke ere et og det samme,
anfører han et Sted (\textit{Et. II, 7. Schol.}). Cirkelen, der
existerer i Naturen, og Cirkelens Begreb. Dette Exempel
anføres ogsaa stadig i Fremstillingerne af Spinozas
Filosofi. Men jeg kan ikke tro andet, end at Spinoza
her har gjort sig skyldig i en øjeblikkelig Forvexling
og Uklarhed. Forholdet mellem den virkelige Cirkel
og Cirkelens Begreb er et Forhold mellem Objekt og
Subjekt, ej mellem Materie og Bevidsthed. Den Form
inden for Udstrækningens Attribut, der svarer til en
vis Form inden for Tænkningens Attribut, er ikke det,
som tænkes og forestilles, men er den materielle Bevægelse
(H jerneproces), der falder sammen med den tænkende
Virksomhed. Naar vi forestille os en Genstand,
sker det, siger Spinoza, derved, at et andet Legeme
paavirker vort Legeme og frembringer en vis Tilstand
i det: vor Forestilling om den ydre Genstand er saa
den Tankeform, som svarer til den Tilstand, vort Le%
% --- p. 101 ---
\setcounter{footnote}{0}%
\opage{101}geme kommer i ved Paavirkning af det andet Legeme.
Vor Bevidsthed opfatter intet ydre Legeme som virkelig
existerende uden gennem Forestillingerne om vort
eget Legemes Tilstande, og derfor kan vort Legemes
Tilstande kaldes Tingenes Billeder, skønt de ikke fremstille
deres Figur. Til min Forestilling om Bordet
svarer inden for Udstrækningens Attribut --- ikke Bordet,
men --- den Tilstand, hvori mit Legeme (særlig
Hjernen) sættes ved Paavirkning af Bordet. Til Begrebet
om Cirkelen svarer paa samme Maade, den
materielle Bevægelse (i Hjernen), som foregaar, naar
jeg tænker Cirkelen. --- Hvad der har foranlediget
Skævheden i det anførte Exempel hos Spinoza, er vistnok
hans underlige Sprogbrug, i det han kalder det
menneskelige Legeme for den menneskelige Bevidstheds
\emph{Objekt}. Den legemlige Proces, der svarer til Bevidsthedsprocessen,
er, som Anatomiens og Fysiologiens Historie
viser, først meget sent og ad mange Omveje
begyndt at blive Objekt for Bevidstheden, og den
umiddelbare Bevidsthed aner maaske slet ikke, at en
saadan materiel Proces er nødvendig. Skønt de to
Problemer: om Forholdet mellem Subjekt og Objekt
og om Forholdet mellem Bevidsthed og Materie berøre
hinanden og gaa over i hinanden, er det dog af
Vigtighed ikke at sammenblande dem. Den spekulative
Filosofi i Tyskland, som i Begyndelsen af dette Aarhundrede
fornyede Spinozismen, foretog netop den
væsentlige Forandring helt igennem at sætte Subjekt
og Objekt i Steden for Bevidsthed og Materie og fik
netop derved en abstrakt logisk-idealistisk Karakter,
% --- p. 102 ---
\setcounter{footnote}{0}%
\opage{102}som aabnede Vejen for mange Fantasterier. Schelling
erklærede Spinozas Filosofi for „ensidig Realisme“;
men hvad vi savne hos Schelling selv og hos
hele hans Retning er netop Spinozas geniale Blik og
dybe Respekt for Tingene reale, lovmæsssige Sammenhæng.


% --- p. 103 ---
\setcounter{footnote}{0}%
\opage{103}\begin{center}III.\end{center}

Til enhver Tanke svarer der, i Følge Spinozas Lære,
en legemlig Bevægelse. Men Bevægelserne i vort Legeme
opstaa ikke af sig selv; de fremkaldes ved Bevægelser,
der udgaa fra andre Legemer. Vort Legeme
kan slet ikke bestaa, naar det ikke bestandig genfødes
ved Paavirkninger fra andre Legemer. Heraf følger
fremdeles, at den Tilstand. hvori vort Legeme i ethvert
givet Øjeblik befinder sig, ikke blot er bestemt ved
dets egen Natur, men ved de andre Legemers Natur,
som paavirke det. De Tanker, der svare til de forskellige
legemlige Tilstande og Bevægelser, ere derfor Udtryk,
ej blot for vort eget, men ogsaa for flere andre
Legemers Natur. Derved opstaa, hvad vi kalde sanselige
Forestillinger \textit{(imaginationes)}. Vi mene i dem at
have Billeder af Tingene; men de svare kun til den
Maade, hvorpaa vort Legeme paavirkes af Tingene. Ofte
indbilde vi os derfor at have Sanseiagttagelser af Ting,
som ikke ere nærværende for os, i det vort Legeme af
indre Aarsager kan bringes i samme Tilstand som af
ydre Aarsager. Den Fejltagelse, vi herved begaa, udspringer
af, at der ikke opstaar nogen anden Forestilling
hos os, som udelukker den formentlig iagttagne
Genstands Nærværelse. Vildfarelse er altsaa noget
negativt og bestaar i Manglen af andre Forestillinger,
% --- p. 104 ---
\setcounter{footnote}{0}%
\opage{104}der kunde overbevise os om, at Genstanden ikke er til
Stede. Den, som ikke kender Hestens Natur, kan
meget godt tro, at han ser en bevinget Hest, og nærer
ingen Tvivl om dens Realitet.\footnote[1]{Denne Teori stemmer ganske med den nyere Psykologi. Taine har givet den et epigrammatisk Udtryk i den Definition: la sensation est une hallucition vraie. Dog var denne Synsmaade allerede tidligere fremsat af den dygtige, men lidet kendte Forsker \emph{Louis Peisse}.} % PRINTED AS IS: „hallucition“ (Rettelser: „S. 10“ [sic, for p. 104], l. hallucination)

Naar to eller flere Legemer have paavirket vort
Legeme samtidig, vil Bevidstheden, naar den senere
iagttager det ene af dem, ogsaa komme til at tænke
paa det andet. Herpaa beror Erindringen, en Forbindelse
eller Sammenkædning mellem Forestillingerne,
som danner sig i Bevidstheden i Overensstemmelse
med de legemlige Tilstandes Orden og Forbindelse.
I denne Forbindelse have vi ingen Erkendelse af Tingenes
Natur; den er af rent subjektiv Natur. Derfor
vil det ogsaa hos de forskellige Mennesker være forskellige
Forestillinger, der paa denne Maade forbindes.
Naar en Soldat ser Hestespor i Sandet, tænker han
strax paa en Rytter, dernæst paa Krig etc.; hos en
Bonde derimod fører Forestillingen om Hesten til den
om en Plov, en Mark etc. Forbindelserne bero paa,
hvorledes enhver er vant til at se Tingene følges ad. ---
Denne Lov, Ideassociationens Lov, nævner Spinoza et
Sted som en aandelig Naturlov, et Sidestykke til de
fysiske Bevægelseslove. Der findes altsaa Love saa vel
paa det aandelige som paa det legemlige Omraade.
Læren om Ideassociationerne var kort før Spinoza bleven
udviklet i klar og koncis Form af Englænderen
% --- p. 105 ---
\setcounter{footnote}{0}%
\opage{105}Thomas Hobbes; men Spinoza er en af de ældre Forfattere,
som har behandlet denne vigtige psykologiske
Lære paa den skarpsindigste Maade, og som har set
dens store og gennemgribende Betydning, hvilket vil
vise sig, naar vi komme til hans Lære om Lidenskaberne.
Dersom Spinoza havde været mere kendt i
England, vilde han sikkert være bleven nævnt som
en af den moderne engelske Psykologis vigtigste Forgængere.


Saa længe de sanselige Forestillinger og Ideassosiationerne
herske, forholde vi os passive og naa ingen
virkelig Erkendelse. Ingen Forestilling og ingen Ideassociation
kan forklares ud af vort eget Væsen alene.
De svare jo stedse til den Maade, hvorpaa vort Legeme
paavirkes eller er blevet paavirket af andre Legemer, og
kunne derfor kun forstaas, naar man tager Hensyn til alle
disse Legemer og deres Vexelvirkning. Holde vi os
alene til den menneskelige Aand, ere Sanseiagttagelserne
og Ideassociationerne Konsekventser uden Forudsætninger.
Det menneskelige Væsen er kun ét Led i en
uendelig Række og kan ikke forstaas, uden man tager
Hensyn til de andre Led i Rækken. Eller for at se
Sagen fra dens højeste Synspunkt: en fuldkommen
(adækvat) Erkendelse af de ydre Legemer findes ikke
i Gud, for saa vidt han betragtes som udtrykt i den
menneskelige Aand alene, men kun for saa vidt han
tillige betragtes som udtrykt i mange andre Væsener.

Saa længe vi blive paa dette Erkendelsestrin, naa
vi kun en forvirret og ufuldstændig Erkendelse, beroende
paa den tilfældige Maade, hvorpaa Tingene paavirke
os. Denne første Erkendelsesart, der af Spinoza
% --- p. 106 ---
\setcounter{footnote}{0}%
\opage{106}betegnes som Mening eller sanselig Forestilling (Imagination),
har ikke andre Kilder end hvad man hører
af andre eller tilfældigvis erfarer \textit{(experientia vaga)}.
Tingenes Væsen undgaar os her bestandig; det er kun
deres ydre Former og Forhold, vi kunne opfatte. Men
Tingenes ydre Former og Forhold bero paa deres indre
Væsen, og det gælder derfor at finde en Erkendelsesart,
der kan aabne os Adgang til dette.

Det er karakteristisk for Spinoza, at han med saa
stor Hast forlader Erfaringsmetoden og ikke undersøger,
om man ikke kunde raade Bod paa de Mangler, hvoraf
den ved første Øjekast synes at lide. Dens Ufuldkommenhed
beror i Følge Spinoza paa, at den ser Tingene
fra det enkelte Væsens Standpunkt, som dog kun er et
Led i Rækken. Men det har netop været hele den
nyere Experimentalforskens Bestræbelse at komme ud
over dette isolerede og tilfældige i Udgangspunktet, og
med hvilket Held den har stræbt derefter, derom bærer
Naturvidenskabens Historie glimrende Vidnesbyrd. Allerede
Baco stræbte at formulere en Metode, der rakte
ud over \textit{experientia vaga}. At man ad Induktionens
Vej kan naa til almindelige Love, derpaa har Spinoza
selv afgivet Bevis; ti det er netop ad denne Vej han
er kommen til at opstille sin Lære om Individualitetens
Grader i Naturen, sin Erkendelseslære, sin Psykologi,
sin Politik. At han i sin Metodelære staar den nyere
Forsken saa fjernt, har sin Grund i hans spekulative
Udgangspunkt, der atter er et Udtryk for hans Higen
efter at gribe det absolute, det i sig værende, der ej
kan være givet i nogen Erfaring. --- Paa den anden
Side kunde det synes, som om netop Spinozas Grund%
% --- p. 107 ---
\setcounter{footnote}{0}%
\opage{107}anskuelse maatte føre ham til at lægge en principiel
Vægt paa Erfaringsmetoden og paa Opfattelsen af det
enkelte og endelige. Naar Dualismen forkastes, og alt
i Naturen sés som Ytringsform for ét Grundvæsen, saa
bliver Erkendelsen af den i Erfaringen givne Virkelighed
ogsaa en Erkendelse af selve det uendelige, af
Grundvæsenet. Spinoza siger derfor ogsaa et Sted:
„Erkendelsen af Virkningen er ét med en fuldkomnere
Erkendelse af Aarsagen. Vi kunne derfor ikke vinde
Indsigt i Naturen, uden tillige at udvide vor Kundskab
om den første Aarsag.“ Men Spinoza opfatter, lige som
Descartes, den rette Metode som den, der gaar ud fra
den højeste Aarsag. Kendskabet til det i Erfaringen
givne faar kun den Betydning at lede vor Tanke mod
visse bestemte Ting. Selve Erkendelsen af Tingenes
Væsen maa øses af en anden Kilde.

Den anden Erkendelsesart er hævet over det isolerede
og subjektive Standpunkt, som var den sanselige
Forestillings (I maginationens) Ufuldkommenhed. Alle
Legemer have noget til fælles: de høre alle ind under
samme Attribut, og de almindelige Bevægelseslove gælde
for dem alle. Af dette fælles kan der haves en fuldkommen
Erkendelse; ti det udtrykker ikke blot det
enkelte Individs Tilstand eller Natur, men Erkendelsen
af det er almengyldig. Ved at fatte Tingenes almindelige
og lovmæssige Sammenhæng hæver Mennesket sig
over Imaginationens Tilfældighed og Udvorteshed. Denne
Erkendelse angaar ikke noget, som kun gælder for et
vist bestemt Tidspunkt, lige saa lidt som for en enkelt
Ting, men noget, som gælder uden Hensyn til Tiden.
At se Tingene i deres lovmæssige Sammenhæng er
% --- p. 108 ---
\setcounter{footnote}{0}%
\opage{108}altsaa at se dem under Evighedens Synspunkt. Tingenes
Nødvendighed er ét med Nødvendigheden i Guds
egen Natur: Gud frembringer jo alt, i det han frembringer
sig selv. Tingene miste her deres ydre Selvstændighed
over for hverandre; vi erkende, at den Kraft,
hvormed enhver af dem bestaar i sin Existents, følger
af Naturnødvendigheden i Guds Væsen. Paa to forskellige
Maader kunne vi altsaa betragte Tingene som
virkelige, enten saaledes at vi betragte dem som existerende
i et vist Tidspunkt og paa et vist Sted i Rummet,

--- hvilket er Imaginationens, den første Erkendelsesarts,
Betragtningsmaade, --- eller saaledes, at vi se dem
indeholdte i Gud og følgende af den guddommelige
Naturs Nødvendighed, --- og dette er den anden Erkendelsesart,
som Spinoza kalder Fornuften (ratio).

Jo mere Mennesket betragter Tingene som tilfældige,
des mere er han selv underkastet Tilfældighed.
Hvad vi erfare ved Sanseiagttagelse og ved Overlevering,
raade vi ikke selv for. Derimod bero de klare og tydelige
Begreber, vi danne, ene paa vor Natur og dens faste
og bestemte Love. Vi behøve med Hensyn til disse
Begreber ikke at hente nogen Bekræftelse og Bestyrkelse
uden fra: Sandhed og Vished falde sammen. Med
den sande Idé følger umiddelbart Visheden om dens
Sandhed. Lige som Lyset hjælper os at skelne mellem
Lys og Mørke, saaledes er det Sandheden, ved hvis
Hjælp vi opdage det usande \textit{(verum est index sui et falsi)}.
Tvivl, Usandhed og Indbildning opstaa kun ved, at vi
betragte Tingene paa forvirret Maade og ikke danne
os tydelige og klare Begreber ved at gaa tilbage til de
simple Elementer og de almindelige Love. Ideen om
% --- p. 109 ---
\setcounter{footnote}{0}%
\opage{109}en firkantet Cirkel er forvirret; ti saa snart man gaar
til dens Elementer, Begrebet Cirkel og Begrebet Firkant,
ser man strax, at de ikke kunne forbindes paa
denne Maade. Og i hvert Tilfælde vil en falsk og opdigtet
Idé snart vise sig i sin Usandhed, naar vi kun
omhyggelig drage alle dens Konsekventser.

For ret at opfatte, hvad Spinoza mener med Fornufterkendelse,
er det nødvendigt at skelne mellem
det fælles og det abstrakte. Fornufterkendelsen har
til sin Genstand Tingenes almindelige Love og derved
Tingenes Væsen. De Love, der udspringe af Tingenes
Væsen, gælde saa vel for det hele som for hver enkelt
Del og ere altsaa Udtryk for Enheden i alt. De svare
til Begrebet om Attributer eller Grundformer, som jo
hver for sig ere ét med Substantsen eller Grundvæsenet;
hver Gang vi betragte et Legeme efter Udstrækningens
almindelige Lov eller en Tanke efter Tænkningens
almindelige Lov, opfatte vi dem som Udtryk og Former
for Grundvæsenet. --- De abstrakte Begreber dannes
derimod, i Følge Spinoza, af Imaginationen eller den
sanselige Forestillingsmaade og bero paa en forvirret
Opfattelse af en Mængde Enkeltheder. Her hen høre
alle Almenbegreber, som egentlig ere skematiske Billeder,
der opstaa hos os, naar vi ofte have haft Forestillinger
om legemlige Genstande (Heste, Mennesker o.
s. v.), og nu disse Forestillinger lige som afslibe deres
Forskelligheder, saa der kun bliver det tilbage, hvori
de alle stemme overens. Hvert Menneske danner disse
Almenbegreber paa sin Maade, efter de Træk hos Genstandene,
der særlig have gjort Indtryk paa ham. ---
Forskellen mellem det abstrakte og det fælles beror
% --- p. 110 ---
\setcounter{footnote}{0}%
\opage{110}altsaa paa, at hint giver et utydeligt Skema, dette en
klar og tydelig Loverkendelse og i Loverkendelsen en
Erkendelse af det, som virker i alt. I den sande Erkendelse
følge vore Tanker Tingenes evige Orden; vor
Aand er et \textit{automaton spirituale}, hvis Virken følger
samme Love som den guddommelige Natur.

Spinozas Forbillede ved Skildringen af denne Erkendelsesform
er den matematisk -mekaniske Naturopfattelse.
Han vilde have opstillet et nøjere Forhold
mellem den første og den anden Erkendelsesform, dersom
han havde haft større Sympati for og klarere
Forstaaelse af den experimentale Metode. Ti Vejen,
ad hvilken man naar til Anvendelse af den matematiske
Metode i Naturvidenskaben, er den nøjagtige Experimenteren,
som gør kvantitative Bestemmelser mulige. ---
Fornufterkendelsen angaar jo i det hele Love og Forhold;
men de Elementer, de virkelige Existentser,
mellem hvilke disse Forhold finde Sted, og for hvilke
Lovene gælde, kunne vi kun kende ved Erfaringen. Som
tidligere vist, ere selve Attributets og Substantsens
Begreber udsprungne ved Analyse af det i Erfaringen
givne. --- Atter her se vi, at der ikke, som Spinoza
mente, er nogen fuldstændig Modsætning eller nogen
skarp Grænse mellem Erfaringserkendelse og Fornufterkendelse.


Men Fornufterkendelsen har før Spinoza endnu det % PRINTED AS IS: „før“
utilfredsstillende, at den er af diskursiv Karakter,
drager Slutninger fra en almindelig Lov til de enkelte
Tilfælde. For at oplyse Forholdet mellem Erkendelsesarterne
bruger han følgende Sammenligning. Naar tre
Tal ere givne, og der skal findes et fjerde, som forholder
% --- p. 111 ---
\setcounter{footnote}{0}%
\opage{111}sig til det tredje, lige som det andet forholder sig til
det første, multiplicerer den, der blot har lært Regula
de tri det andet og tredje Tal sammen og dividerer
Produktet med det første. Hans Kundskab er rent
empirisk, beror paa Overlevering og Erfaring og hører
altsaa ind under den første Erkendelsesart. Den, der
kan Matematik, foretager derimod den samme Operation,
fordi den er en Anvendelse af Læren om Proportionerne.
Dette svarer til den anden Erkendelsesart.
Men naar det er meget simple Tal, behøver man slet
ingen Operation at foretage. Man kan f. Ex. umiddelbart
indsé, at 1 forholder sig til 2 som 3 til 6. Den
almindelige Regel bliver her overflødig; i Steden for
Beregning og Ræsonnement træder umiddelbar Skuen
af Forholdet; i Steden for det diskursive træder det
intuitive.

Den tredje og højeste Erkendelsesform er den
intuitive Viden. Den skuer Tingenes Væsen umiddelbart,
uden Hjælp af almindelige Sætninger. Som Exempel
paa saadan intuitiv Viden anfører Spinoza, at
naar vi vide noget, have vi dermed ogsaa en umiddelbar
Erkendelse af, hvad det er at vide. I „Etikens“
første Bog har han efter eget Sigende anvendt den
anden Erkendelsesform. Han opstillede der en Række
Definitioner og Grundsætninger og udledte af dem sin
Lære om, at alt, baade hvad Væsen og Existents angaar,
beror paa Grundvæsenet. Først ved „Etikens“
Slutning (i femte Bog) bliver denne Lære Genstand
for intuitiv Viden, i det den menneskelige Aand umiddelbart
ser sig som Element i Grundvæsenet. Her
have vi ikke længer almindelige Regler, men en
% --- p. 112 ---
\setcounter{footnote}{0}%
\opage{112}umiddelbar Skuen af det enkeltes indre Sammenhæng
med det hele.

I den „kortfattede Afhandling“ skildres
denne
højeste Erkendelsesform med Udtryk, der have et vist
mystisk og religiøst Præg. Vi mindes derved om det
oprindelige praktiske Motiv for Spinozas Tænkning.
„Fornuftslutningen er ikke det fortrinligste i os, men
er kun som en Stige, ad hvilken vi svinge os op til
det forønskede Punkt, eller som en god Aand, der
uden Falskhed og Bedrag melder os om det højeste
Gode for derved at opfordre os til at søge det og forene os dermed, hvilken Forening er vor højeste Velfærd
og Lykke.“ „At denne [tredje] Erkendelsesart
ikke er en Følge af noget andet, men er umiddelbar,
indlyser deraf, at Gud er Aarsagen til al Erkendelse
og kun erkendes ved sig selv, ikke ved noget andet,
og deraf, at vi af Naturen ere saaledes forbundne med
ham, at vi uden ham hverken kunne bestaa eller forstaas.
Af denne nøje Forbindelse mellem Gud og os
følger, at han kun kan forstaas umiddelbart.“ --- Den
Maade, hvorpaa vi saaledes umiddelbart opfatte Grundvæsenet
eller Guddommen, betegnes dels som en umiddelbar
Aabenbaring for Forstanden, dels som en Fornemmelse
og Nydelse af Tingene i sig selv. --- Ogsaa
i sine senere Skrifter betegner Spinoza den tredje
Erkendelsesform som den højeste Dyd og den højeste
Salighed.

Der er den Vanskelighed ved Spinozas Skildring
af den tredje Erkendelsesform, at han nogle Steder
beskriver den som en aldeles umiddelbar Anskuen,
men andre Steder bestemmer den som en Tanke%
% --- p. 113 ---
\setcounter{footnote}{0}%
\opage{113}bevægelse, der gaar fra Grundformernes Erkendelse til
Erkendelsen af Tingenes Væsen, eller som en Erkendelse
af Tingen ved dens nærmeste Aarsag. Ved denne
sidste Bestemmelse bliver den ikke ganske forskellig
fra Fornufterkendelsen, som den dog skal hæves over;
den viser uvilkaarlig sit Slægtskab med denne, lige som
vi oven for saa et naturligt Slægtskab mellem Fornufterkendelsen
og Erfaringserkendelsen, der heller ikke
vare saa fremmede for hinanden, som Spinoza antog.

Motivet til at søge ud over Fornufterkendelsen
ville vi let finde, naar vi erindre Spinozas Grundbegreb
og dets Definition: „Substants er det, som er i sig
selv og begribes ved sig selv, d. v. s. det, hvis Begreb
ikke trænger til nogen anden Tings Begreb, af hvilket det
skulde dannes.“ Dersom et Begreb skal opfattes i og
ved sig selv, uden Benyttelse af andre Begreber, saa er
det klart, at der ingen Tankeoperation eller Slutning
kan finde Sted. Hvor hen skulde Tankebevægelsen gaa,
naar der kun er ét Begreb? Hvad der forstaas „i og
ved sig selv“, er Genstand for Skuen, ej for Tænkning.
Spinoza kunde derfor ikke fastholde sit Substantsbegreb
uden at opstille en fra den almindelige Erfaring og
Tænkning forskellig og over dem hævet Erkendelsesform.
Det samme gælder om alle de Filosofer, der
have hævdet en direkte og positiv Erkendelse af det
absolute. For Plato var den højeste Viden „en Skuen  
af det bedste af det værende“, nemlig af „den første
forudsætningsløse Grund til alt“. Forudsætningsløs skal
jo ogsaa Spinozas Substants være. Men hvad der skal
være aldeles uden Forudsætninger, kan ikke tænkes;
det maa anskues. Schelling satte dette paa Spidsen i
% --- p. 114 ---
\setcounter{footnote}{0}%
\opage{114}sin Lære om den „intellektuelle Anskuen“, der er en
umiddelbar Opfattelse af det absolute som den evige
Enhed af Subjekt og Objekt, af alle Modsætninger.
Denne Anskuen var for Schelling et særeget Organ,
som man enten havde eller ikke havde. De, der
manglede den, manglede ligefrem Betingelsen for at
filosofere, lige som den vilde mangle Betingelsen for at
være Matematiker, der ikke kunde konstruere og anskue
Figurer i Rummet. --- Lige som den spekulative Filosofi
maa appellere til et særligt Organ, naar den vil
hævde en positiv Erkendelse af det absolute, saaledes
er det ogsaa Tilfældet med den positive Religion og
Teologien. Ogsaa disse pege ud over Erfaring og Tænkning
til en højere Erkendelseskilde. Ved Genfødelsen
og Troen aabnes Adgangen til en højere Natur end den,
som er tilgængelig for den ikke troende. --- Vi se atter
her, hvorledes den spekulative Filosofi minder om den
positiv-religiøse Bevidsthed. Naar Filosofien derimod
anerkender Relativitetens Lov og gennemfører den fuldstændig,
maa den opgive Appellen til en særegen Erkendelseskilde
og stille sig paa lige Trin med al anden
videnskabelig Forsken. Videnskaben har, lige som al
Kultur, oprindelig udfoldet sig under Religionens beskyttende
Dække, og Filosofien bærer endnu mange
Mærker og Rester af dette Hylster til Skue. Den vil
uden Tvivl ogsaa afkaste dem, og den vil kunne det
uden at krænke nogen virkelig menneskelig Interesse.

Hermed skal ikke være sagt, at Ideen om en
intuitiv Viden ganske er at afvise. Det urigtige bestaar
kun i at betragte den som en særegen Erkendelseskilde.
Erkendelsens Udvikling naar ikke sit Højdepunkt i det
% --- p. 115 ---
\setcounter{footnote}{0}%
\opage{115}blotte Ræsonnement. Enhver véd, at man først fuldstændig
har tilegnet sig en Kundskab, naar man kan
overspringe eller overskue alle eller dog mange Mellembestemmelser
og se den hele Sammenhæng i ét Blik.
Intuition eller Forstandsanskuelse er altsaa et naturligt
Resultat af Anstrengelse og Øvelse med videnskabelige
Tankerækker. I denne Betydning gælder det, at kun
den intuitive Viden er Kronen paa vort Erkendelsesarbejde.
Men den staar saa i den nøjeste Forbindelse
med de andre Erkendelsestrin og voxer paa naturlig
Maade frem af dem, lige som gode Vaner og Instinkter
voxe frem af Viljens Anstrengelser. Spinozas eget oven
nævnte Exempel viser i denne Retning. Ti at 1 forholder
sig til 2, som 3 til 6, sés kun umiddelbart af
den, der er øvet i at operere med Tal; Barnet, der
endnu ikke kan sin Tabel, skal først møjsommelig arbejde
sig frem til en saadan Intuition.

--- Der er en Klasse af Begreber, som Spinoza frakender
absolut Gyldighed, fordi de lide af de Ufuldkommenheder,
der ere ejendommelige for den første
Erkendelsesart. Det er saadanne Begreber som godt
og ondt, skønt og hæsligt, Orden og Uorden, Hensigtsmæssighed
o. s. v., naar de anvendes paa Tingenes
Væsen. Imaginationens eller den sanselige Forestillings
Ufuldkommenhed beroede jo paa, at den betragter alt
fra Menneskets individuelle og tilfældige Standpunkt og
maaler alt efter dette. Mennesket bedømmer fra først
af alt praktisk, betragter alle Ting omkring sig som
Midler og Redskaber og kalder godt, hvad han kan
bruge, ondt, hvad der skader og hindrer ham. Han
vænner sig saaledes til først og fremmest at spørge ved
% --- p. 116 ---
\setcounter{footnote}{0}%
\opage{116}enhver Ting, han sér, hvortil den tjener, i Steden for,
hvorledes den er bleven til. Naar han saa har faaet at
vide, at Øjnene ere til at se med, Planter og Dyr til at
spise, Solen til at lyse, har han Forklaring nok. Herfra
stammer Forestillingen om Naturens Hensigtsmæssighed.
Nu er der paa den anden Side meget skadeligt
og fordærveligt i Naturen; det mener man saa
skyldes Gudernes Vrede over Menneskenes Synder.
Og naar Erfaringen viser, at Ulykkerne lige saa vel
ramme de gode som de onde Mennesker, erklærer man,
at Gudernes Raadslutninger overgaa den menneskelige
Fatteevne. I Steden for at afhjælpe sin Uvidenhed beraaber
man sig altsaa paa den som et Argument.
Dersom Matematiken ikke havde anvist Menneskene
en anden Sandhedsregel, vilde Menneskeslægten være forbleven
i sin Uvidenhed. Men de gamle Fordomme ere
endnu saa stærke, at den, der søger videnskabelig Indsigt
i Naturen i Steden for at vedblive med taabelig
Undren, anses for ugudelig. --- Lige som Menneskene
tillægge Naturen Øjemed efter deres egne Ønsker,
saaledes finde de ogsaa Orden og Uorden, Skønhed og
Hæslighed i Naturen, alt efter som Tingene ere lette
eller vanskelige at anskue og opfatte eller ikke, tiltale
deres Opfattelsesevne eller ikke. Men disse Begreber,
i det hele alle Begreber om Fuldkommenhed og Ufuldkommenhed,
udspringe af Sammenligninger, som vi anstille,
og udsige intet om Naturen i sig selv. Vi danne
os et Almenbegreb, en vis Type, der passer for alle
Individer af en vis Art, og jo mere det enkelte Individ
nærmer sig denne Type, des fuldkomnere sige vi, at det
er. Men dette Almenbegreb, denne Type tilhører vor
% --- p. 117 ---
\setcounter{footnote}{0}%
\opage{117}Forstand, og ikke Naturen. Det uendelige Væsen,
som vi kalde Gud eller Naturen, virker med samme
Nødvendighed, som det existerer. Vore Begreber om
Fuldkommenhed og Ufuldkommenhed kunne derfor ikke
anvendes paa det. Kun naar man ved Fuldkommenhed
forstaar Realitet, Væren, kan dette Begreb anvendes
paa Naturen; Fuldkommenheden strækker sig da lige
saa langt som Tingenes sande Væsen. Men Begrænsning
tilkommer kun det endelige, ikke det uendelige
Væsen. Naar vi finde Ufuldkommenhed i Naturen, er
det altsaa kun, fordi vi betragte den fra et endeligt
Synspunkt: i sig selv er alt fuldkomment, naar man
ikke stiller det over for noget andet, med hvilket man
sammenligner det, d. v. s. naar man ikke betragter det
som endeligt. Naturen virker ikke efter den menneskelige
Fornufts Lov og retter sig ikke efter de menneskelige
Formaal og Ønsker, men efter uendelig mange
andre Love, der udspringe af den evige Naturorden,
i hvilken Mennesket kun er et enkelt Led. Naar noget
derfor synes os latterligt, urimeligt eller ondt, kommer
det af vor begrænsede Viden. Hvad vor Fornuft erklærer
for ondt, er kun ondt efter \emph{vor} Naturs Love, ej
efter \emph{hele} Naturens Orden og Lov. Den, der stræber
efter Erkendelse af Tingenes sande Væsen, vil ikke
dadle eller spotte noget i Naturen, fordi han véd, at
alt skér i Følge den guddommelige Naturs evige og
nødvendige Love.

Naar æstetiske og etiske Begreber ikke kunne anvendes
paa Naturen i sig selv, fordi de kun opstaa ved
Sammenligninger og Vurderinger, vi anstille fra vort
endelige Standpunkt, saa synes det at maatte gælde
% --- p. 118 ---
\setcounter{footnote}{0}%
\opage{118}for endnu flere Begrebers Vedkommende, og det har
Spinoza ogsaa indsét. Som vi allerede have set, ere
Tid, Tal og Maal Bestemmelser, der ikke gælde for
den sande Væren. Tiden er ikke en Bestemmelse ved
Tingene selv, men blot en Forestillingsmaade hos os,
hvorved vi bestemme Varigheden af en Tings Existents.
Men i Gud er der intet tidligere eller senere; hans
Existents --- og altsaa al Existents --- er en nødvendig
eller evig Følge af hans Væsen.

Denne Betragtningsmaade burde sikkert være anvendt
paa alle andre Begreber ogsaa. Spinoza vilde
have gjort det, naar han ikke som god Dogmatiker
havde troet paa Muligheden af, at vi kunne naa til at
betragte Tingene fra et andet Synspunkt end vort eget.
Saadanne Begreber som Substants og Aarsag, hvilke
Spinoza uden Betænkning anvender paa det absolute,
ville ved nærmere Betragtning ogsaa vise sig at være
relative. Dog er det mærkeligt at se, i hvor høj en
Grad det er lykkedes Spinoza at emancipere sig fra
Antropomorfismen. Den naive Betragtning anvender
uden videre alle Bestemmelser og Begreber paa Tingene.
Den nyere Filosofi er i sin successive Udvikling
ført til at vise alle disse Begreber tilbage til den
menneskelige Bevidsthed og se dem som Symboler,
der udtrykke Tilværelsen fra et bestemt Synspunkt,
uden at vi i dem have den eneste mulige Løsning af
den store Gaade. Man har ofte skarpt bebrejdet Spinoza
hans Lære, at Begreber som godt og ondt ikke
gælde for Tingenes Væsen og kun have relativ Gyldighed.
Deri er han dog i sin gode Ret. Hans Behandling
af de æstetiske og etiske Begreber er en nødvendig
% --- p. 119 ---
\setcounter{footnote}{0}%
\opage{119}Konsekvents af Kopernikus's Astronomi og Galilæi's
og Descartes's Fysik. Naar Jorden ikke længer er
Midtpunktet i Universet, kan Mennesket ikke heller
være det, og hans Domme gælde kun fra et Synspunkt,
der er forsvindende i Forhold til det store, uendelige
Hele. Og naar alt i Naturen underligger den matematiske
Mekaniks Love, kan man ikke forklare dens
Fænomener ved at indskyde ideale Kræfter. Hvad
man maa indvende mod Spinoza er tvært imod, at han
ikke er gaaet vidt nok, og desuden, at han til Dels
antager, at Begrebernes Relativitet er ens betydende
med deres Ugyldighed. At dette ikke forholder sig
saa, indrømmer han i Virkeligheden selv, hvad de etiske
Begreber angaar, ti i sin etiske Lære opstiller han „et
idealt Forbillede for den menneskelige Natur“ og kalder
det godt, som fører os dette Forbillede nærmere,
det ondt, som fjerner os derfra. Naar en saadan Type
virkelig udtrykker det menneskelige Liv, som det i Følge
sine Naturlove vilde være, naar det naaede sin højest
mulige Udvikling, har den og de andre etiske Begreber, 
der knytte sig til den, al den Gyldighed, som kan fordres,
og denne Gyldighed lider intet Skaar ved, at vi
ikke ere i Stand til at afgøre, hvilken Betydning dette
Ideal og vor Stræben efter dets Virkeliggørelse har for
Universet i dets Helhed.\footnote[1]{Sammenlign mit Skrift „Om Grundlaget for den humane Etik.“ Kbhvn. 1876, P. 84---86.}

% --- p. 120 ---
\setcounter{footnote}{0}%
\opage{120}\begin{center}IV.\end{center}


Det er en Fordom, som er ganske almindelig, at
en Etik er ganske umulig paa Determinismens Grund.
Naar en Filosof opfatter de menneskelige Følelser og
Handlinger som underlagte Naturlove, er en vis Art
Kritik strax til rede med den Bemærkning, at han kun
ved en lykkelig Inkonsekvents kan opstille en Etik, en
Lære om Formaalene for den menneskelige Handlen
og om Midlerne til at naa dem. Uden absolut fri
Vilje, siger man, ingen Etik. Ved nærmere Betragtning
vil man let se, paa hvilken Side den „lykkelige
Inkonsekvents“ finder Sted. Etiken vilde netop være
umulig, naar der var en absolut fri Vilje hos Mennesket;
da vilde enhver sammenhængende Lære om den
menneskelige Handlen være umulig. Kæden vilde hvert
Øjeblik kunne brydes. Om bestemte Formaal, begrundede
i den menneskelige Natur, og om bestemte Midler
til deres Virkeliggørelse kunde der da ikke være Tale.
Spinoza var derfor netop som Determinist vel skikket
til at fremstille en etisk Lære. I grundigt psykologisk
Kendskab til den menneskelige Natur og dens Love
er der næppe nogen Etiker, der kommer ham nær; og
naar Etiken skal grundes paa den menneskelige Naturs
Love, er grundig Psykologi den første Betingelse for
en Etiker.

% --- p. 121 ---
\setcounter{footnote}{0}%
\opage{121}En anden Sag er det, at Etiken, lige som enhver
Undersøgelse af Enkeltvæsenernes Udvikling og Love,
i Spinozas System kun faar principiel Betydning, for
saa vidt den experimentale Understrøm formaar at
hævde sig over for den spekulative Overstrøm, der har
sit Udtryk i Spinozas Grundtanke, Substantsens Begreb.
Gaar man ene ud fra dette Begreb, bliver den etiske
Udvikling som al konkret og endelig Udvikling uforklarlig
og illusorisk. Men tage vi begge Strømninger i
Spinozismen, saaledes som de allerede lægge sig for
Dagen i den første Grundsætning i Etikens første Bog
(se oven for P. 87), saa er Spinoza fuldstændig konsekvent,
naar han opstiller en etisk Lære.

Spinoza hævder gentagne Gange med stor Energi,
at de menneskelige Følelser og Lidenskaber skulle
undersøges med samme Ro og Upartiskhed som andre
Naturfænomener. „De fleste“, siger han i Fortalen til
„Etikens“ tredje Bog, „som have skrevet om Lidenskaberne
og om Menneskenes Levevis, synes ikke at
tale om Naturfænomener, der følge Naturens almindelige
Love, men om Ting, der staa uden for Naturen.
De synes at betragte Mennesket i Naturen som en
Stat i Staten. Ti de tro, at Mennesket snarere forstyrrer
Naturens Orden end følger den, og at han har
en absolut Magt over sine Handlinger og ikke kan bestemmes
andre Steder fra end fra sig selv. Naturen
er overalt den samme og overalt én, dens Evne og
Virkekraft, d. v. s. Naturens Love og Regler, efter
hvilke alt sker og gaar over fra en Form til en anden,
ere overalt og altid de samme, og der er derfor kun
én Maade, hvorpaa hvilken som helst Tings Natur kan
% --- p. 122 ---
\setcounter{footnote}{0}%
\opage{122}forstaas, nemlig ved Naturens Love og almindelige
Regler. Hadets, Vredens og Misundelsens Lidenskaber
udspringe altsaa i sig selv af samme Naturnødvendighed
og Naturkraft som andre endelige Ting; de have bestemte
Aarsager, ved hvilke de forklares, og bestemte
Egenskaber, der ere lige saa værdige til vor Erkendelse
som en hvilken som helst anden Tings Egenskaber,
hvis Betragtning glæder os. Jeg vil derfor afhandle
Lidenskabernes Natur og Kræfter og Aandens
Magt over dem efter samme Metode, som jeg har anvendt
i det foregaaende, og jeg vil betragte de menneskelige
Handlinger og Drifter paa samme Maade som
jeg betragter Linjer, Flader og Legemer.“ I Indledningen
til sin politiske Afhandling erklærer Spinoza
ligeledes, at han vil behandle disse Spørgsmaal med
samme aandelige Frihed, som vi medbringe til matematiske
Undersøgelser. „Jeg har alvorlig beflittet mig
paa ikke at spotte, begræde eller forbande de menneskelige
Handlinger, men at forstaa dem, og derfor har
jeg betragtet de menneskelige Lidenskaber, Kærlighed,
Had, Vrede, Misundelse, Ærgerrighed, Medlidenhed og
de øvrige Sindsbevægelser, ikke som Fejl ved den
menneskelige Natur, men som Egenskaber, der høre
til den, lige som der til Luftens Natur hører Varme,
Kulde, Uvejr, Torden og andet lignende, som vel medføre
Ulemper, men dog ere nødvendige og have bestemte
Aarsager, ved hvilke vi stræbe at forstaa deres
Væsen, en Forstaaelse, der bringer Tanken lige saa
megen Glæde som Erkendelsen af de Ting, der medføre
behagelige Følger for os.“ --- Man har undertiden
anført disse Ord med Gysen og betragtet dem som
% --- p. 123 ---
\setcounter{footnote}{0}%
\opage{123}Tegn paa en kold og hjerteløs Tænkemaade. Hertil
er kun at sige, at skal i det hele en Psykologi --- en
Erkendelse af det aandelige Livs Love --- være mulig,
saa maa man opfatte de aandelige Fænomener som
rene Erkendelsesgenstande, uden anden Følelse end
Ønsket om at forstaa og forklare dem saa godt som
muligt. Ingen Forsken er mulig uden Forskerstemning,
og denne Stemning fortrænges, saa snart æstetiske og
moralske Følelser gøre sig gældende. Netop fordi
Spinoza saa klart har set dette, og fordi denne Forskerstemning
var den herskende hos ham, er han bleven
en af de vigtigste Forgængere for den nyere videnskabelige
Psykologi.

Kuno Fischer har træffende vist, hvorledes Spinoza
her har ydet for Filosofien, hvad før ham Shakespeare
allerede havde ydet i Poesien: „Enhver virkelig Karakter
har sine Lidenskaber, og enhver Lidenskab har
sin nødvendige Udviklingsgang, der er dramatisk, fordi
den bestaar i bestemte Handlinger; ethvert sandt Drama
er ikke andet end menneskelige Lidenskabers Naturhistorie.
--- I Filosofien var Spinoza den første, der
begreb og med Grunde beviste denne Sandhed; men
den verdenshistoriske Fortjeneste først over Hovedet
at have opdaget det dramatiske Menneskeliv i denne
Betydning tilkommer en tidligere Aand. Hinsides Filosofien
hæver sig paa Grænseskellet mellem to Tidsaldre
den gigantiske Shakespeare, for hvis digteriske
Anskuelse det først blev klart, at den menneskelige
Verdens Begivenheder fremgaa af Handlingerne, Handlingerne
af Lidenskaberne, Lidenskaberne af Karaktererne
og Karaktererne af Tingenes naturlige og
% --- p. 124 ---
\setcounter{footnote}{0}%
\opage{124}historiske Sammenhæng. --- Mennesket handler i Shakespeares
Fantasi og i Spinozas Filosofi som en Natur,
der opfylder sin Lov, og Lidenskaben, der bevæger
den menneskelige Kraft og driver den fremad, kender
kun én Vej, Strømmens Vej til Faldet.“ --- Fremdeles
gør Kuno Fischer opmærksom paa, at Shakespeare og
Spinoza fejrede deres Genopstandelse paa samme Tid
og hos den samme Kreds af Mænd. Det var særlig
Lessing, der atter drog dem begge frem og banede
Vejen for deres Forstaaelse. (Geschichte der neuern
Philosophie. I. Mannheim 1854. P. 416---420).\footnote[1]{I Fortalen til sit første Skrift henviser \emph{Taine} til Spinoza, og denne Henvisning er karakteristisk for den berømte Kritikers hele Standpunkt og Metode.}

Spinozas Lære om Lidenskaberne (Affekterne),
der danner Underlaget for hans etiske Lære, kan kun
forstaas, naar man gaar ud fra hans Lære om Enkeltvæsenet.
Hvert Enkeltvæsen er en endelig Form for
det uendelige Grundvæsen; alle Enkeltvæsener staa i
Vexelvirkning med hverandre og danne i Forening det
store Naturindivid. For saa vidt Enkeltvæsenet paavirkes
og bestemmes af de andre Enkeltvæsener, er det
indskrænket, begrænset og passivt. Derfor er det kun
delvis Aarsagen til de Virkninger, som udgaa fra det.
Fuldstændig (adækvat) Aarsag vilde et Enkeltvæsen
kun være, naar Virkningen ganske kunde forklares af
dets Væsen alene. Vi siges at handle, naar vi ere
fuldstændige Aarsager, men at lide, naar vi ere ufuldstændige
Aarsager. Ved Affekt eller Lidenskab forstaa
vi en Tilstand (Affektion), hvorved vor Virkekraft
% --- p. 125 ---
\setcounter{footnote}{0}%
\opage{125}forøges og formindskes, og tillige Forestillingen om
denne Tilstand. Lidenskaben er af aktiv Karaktér,
naar vi selv ere fuldstændig Aarsag til Tilstanden, men
af passiv Karaktér, naar vi selv kun delvis ere Aarsag.

--- Det er altsaa som Led i Naturvæsenernes Række,
at Mennesket er passive Lidenskaber underkastet.
Han kan lige saa lidt som noget andet Enkeltvæsen
unddrage sig de ydre Aarsagers Indflydelse. Lidenskaberne
opstaa ved Brydningen mellem de ydre Aarsagers
Virken og Enkeltvæsenets egen Grunddrift.
Enhver Ting søger, for saa vidt den har Bestaaen i sig
selv, at vedblive at bestaa. Kun ydre Aarsager kunne
ophæve en Tings Existens; for saa vidt det kommer an
paa den selv alene, vil den vedblive at bestaa. Denne
Selvopholdelsesdrift er i sig selv ikke andet end selve
den menneskelige Natur, hvoraf det med Nødvendighed
følger, der tjener til dens Bevarelse, saa at Mennesket
naturlig bestemmes til at søge det. --- Spinoza synes
her at have overført Inertiens Lov fra det legemlige
til det aandelige Omraade. Galilæi, og efter ham Kepler
og Descartes, havde vist, at et Legeme vil blive i sin
Tilstand (af Hvile eller Bevægelse), naar andre Legemers
Indflydelse ikke bringer det ud af den. Denne
Lov betragter Spinoza nu som en almindelig Lov, ikke
indskrænket til de fysiske Legemer. Den er i sig selv
ogsaa kun en Form for det almindelige Kavsalitetsprincip,
der fordrer en bestemt Aarsag til enhver
Forandring. For Spinoza laa denne Almindeliggørelsenær.
Alle Enkeltvæsener ere endelige Former for Grundvæsenet,
og Grundvæsenets Kraft virker i dem. Vel
er Enkeltvæsenet endeligt og begrænset, men som Form
% --- p. 126 ---
\setcounter{footnote}{0}%
\opage{126}for Grundvæsenet har det Del i dettes Selvstændighed.
„Naturtingenes Kraft til at existere og til at virke kan
ikke være nogen anden end Guds evige Magt.“ „Skønt
enhver Ting bestemmes af en anden enkelt Ting til at
existere paa en vis Maade, følger dog den Kraft, hvormed
enhver Ting vedvarer i sin Existents, af Guds
Naturs evige Nødvendighed.“

Lidenskaben opstaar nu ved, at denne Grunddrift
hæmmes eller fremmes i sin Virken. Naar Aanden
paa Grund af ydre Paavirkninger gaar over til højere
Tankekraft, opstaar Glæde; naar den gaar over til
mindre Tankekraft, opstaar Sorg. Attraa, Glæde og
Sorg ere de tre oprindeligste Lidenskaber, af hvilke alle
andre udspringe. De Love, hvorefter de udvikle sig,
ere Ideassociationens Love. Naar Glæden ledsages af
Forestillingen om den Aarsag, der har fremkaldt den.
opstaar Kærlighed. Naar Sorgen ledsages af Forestillingen
om dens Aarsag, opstaar Had. Aanden søger
saa vidt muligt at fastholde det, der forøger og fremmer
Virkekraften, og vender sig bort fra det, der hæmmer
den. Den, som elsker, stræber derfor med Nødvendighed
at have den elskede Genstand nærværende og at
bevare den, hvorimod den, der hader, søger at fjerne
og ødelægge Hadets Genstand.\footnote[1]{Descartes havde (i sin Traité des passions) defineret Kærligheden som Vilje til at forbinde sig med den elskede Genstand saaledes, at man sammen med den udgør et Hele. Spinoza indvender mod denne Definition, at den vel udtrykker en væsentlig Egenskab ved Kærligheden, men ikke selve Kærlighedens Væsen. --- En lignende Bemærkning gør Sokrates i Plato's „Symposium“ i Anledning af Aristofanes's mytiske Skildring af, hvorledes Menneskene i Kærligheden søge deres oprindelige Halvdel.}

% --- p. 127 ---
\setcounter{footnote}{0}%
\opage{127}Naar vi paa én Gang have haft to Lidenskaber,
ville vi, dersom den ene af dem atter opstaar, ogsaa
komme til at føle den anden. Paa denne Maade kan
enhver Ting indirekte blive Aarsag til Glæde, Sorg
eller Attraa, i det nemlig den Følelse, Tingen i sig
selv fremkalder, vækker en anden Følelse, som vi derfor
ofte ikke selv kunne forklare, fordi vi have glemt
den oprindelige Forbindelse mellem de to Følelser.
Mange uforklarlige Sympatier og Antipatier have denne
Oprindelse. --- En og samme Ting kan fremkalde blandede
Følelser, i det den direkte er Aarsag til én Følelse,
indirekte til en anden.

Paa Grund af den Egenskab ved Kærligheden, at
den forbinder den elskende og den elskede Genstand
til et Hele, vil den elskede Genstands Følelser ogsaa
opstaa hos den elskende. Kærligheden bliver altsaa
som en Udvidelse af vort Væsen, hvorved nye Aarsager
til Følelser af forskellig Art opstaa. Hvad der vækker
Glæde eller Sorg hos den elskede Genstand, vil fremkalde
Kærlighed eller Had hos den elskende. --- Sorg
bevirket ved andres Ulykke, er Medlidenhed, hvorimod
Spinoza savnede et Ord i Sproget til at udtrykke Glæden
over andres Lykke. Kærligheden til den, der gør
en anden vel, kaldes Billigelse \textit{(favor)}, Hadet til den,
der har paaført en anden ondt, Harme \textit{(indignatio)}.
--- Paa den anden Side føle vi Kærlighed til den, der
paafører én ondt, som vi hade, og Had til den, der
gør ham vel.

Naar vi forestille os, at nogen elsker, attraar eller
hader noget, som vi selv elske, attraa eller hade, ville
vi derved komme til at elske, attraa eller hade det
% --- p. 128 ---
\setcounter{footnote}{0}%
\opage{128}endnu mere, hvorimod vi ville føle en Vaklen i vort
Sind \textit{(animi fluctuatio)}, naar vi forestille os, at han
afskyr, hvad vi attraa, eller omvendt. Enhver vil derfor
af Naturen stræbe efter, at de andre komme til at
leve efter hans Sind. Og i det nu alle stræbe herefter,
hindre de i lige Grad hverandre. Enhver vil roses
eller elskes af alle andre, og derfor hade alle hverandre
indbyrdes. Ærgerrigheden bliver derfor den stærkeste
af alle Lidenskaber; den drager Næring fra mange
Sider. Den søger Menneskenes Bifald, men vil paa
den anden Side ogsaa, at alle skulde bifalde, hvad den
selv bifalder.

Tanken om, at et Væsen, der ligner os, og over
for hvilket vi ej have nogen forudfattet Følelse, har en
eller anden Følelse, vækker den samme Følelse hos os.
For saa vidt nemlig det andet Væsen er os ligt, kommer
det ud paa ét, om vi forestille os det i en vis Tilstand
eller vi forestille os selv i den samme Tilstand. Vi
glæde os og sørge med andre, blot fordi de ere Væsener,
der ere os lige. Selv vor Fjendes Undergang vækker
derfor Sorg hos os, fordi han var et Menneske som vi.
Her hen hører fremdeles al uvilkaarlig Efterligning, som
naar én trækker Haanden til sig, fordi han ser en
anden brænde sig. Efter samme Lov vækker en andens
Glæde ved en Genstand vor Attraa efter den, og vi
søge at berøve ham Genstanden. Det gaar altsaa saaledes
med den menneskelige Natur, at vi ynkes over
dem, der have det ilde, og misunde dem, der have det
godt. Barmhjertighed, Misundelse og Ærgerrighed
udspringe af en og samme Egenskab ved den menneskelige
Natur. Dette vil man især sande, naar man
% --- p. 129 ---
\setcounter{footnote}{0}%
\opage{129}iagttager Mennesket i de første Aar af dets Liv. Det
græder og ler, efter som det ser andre græde og le,
og det føler Attraa efter, hvad det ser andre glæde
sig ved og attraa. Den Rolle, som Misundelsen og
Ærgerrigheden spille i det menneskelige Liv, næres
ved Opdragelsen, i det Forældrene særlig benytte disse
Lidenskaber for at anspore Børnene til Dyd.

I ethvert enkelt Tilfælde faar Lidenskaben en forskellig
Karakter efter den særlige Genstand, der vækker
den, og efter det Individs Ejendommelighed, hos hvilket
den opstaar. Et og samme Individ paavirkes forskellig
af forskellige Genstande; samme Genstand paavirker
forskellige Individer forskellig, og et og samme Individ
forskellig til forskellige Tider. Lidenskaben er jo ikke
andet end selve Enkeltvæsenets Natur, og den enkelte
bestemte Lidenskabs Væsen beror paa, hvorledes det
paavirkes af andre Enkeltvæsener, andre Led i Naturens
Kæde.

Ved Sammensætning og Vexelvirkning af de forskellige
Lidenskaber kan der opstaa en Mængde afledede
Lidenskaber. De danne sig med samme Nødvendighed
som de oprindelige Lidenskaber. Vi sættes paa
mange Maader i Bevægelse af ydre Aarsager og omtumles
lige som Havets Bølger, der bevæges af forskellige
Vinde, uden at vi kende Maal og Udgang.

--- Spinozas Skildring af Lidenskabernes Udspring og
Udvikling er ikke blot karakteristisk for hans Aandsretning
og hans Standpunkt, for hans Søgen efter den
nødvendige Sammenhæng i alt, men er ogsaa i sig selv
et Mønster paa psykologisk Undersøgelse og har stor
blivende Værdi. I Fastlandets psykologiske Litteratur
% --- p. 130 ---
\setcounter{footnote}{0}%
\opage{130}er denne Skildring endnu uovertruffen; England har
derimod i Hobbes, Hartley, Hume og i vor Tid især
Bain en Række af Mænd, der have behandlet samme
Emne efter samme Metode som Spinoza og kunne
stilles ved Siden af ham.

Den Grundbestemmelse af Glæde og Sorg (Lyst
og Smerte), hvorfra Spinoza gaar ud, stemmer ganske
med den nyere Psykologis. Glæde, siger han, er\footnote[1]{Rigtigere vilde det have været at sige „opstaar ved“ i Steden for „er“. Disse simple psykiske Fænomener kunne kun defineres ved det, der fremkalder dem.}
Overgang til højere, Sorg til mindre Fuldkommenhed;
og han gør udtrykkelig opmærksom paa, at Glæden
ikke er selve Fuldkommenheden eller Sorgen selve
Ufuldkommenheden. „Ti dersom Mennesket blev født
med den Fuldkommenhed, til hvilken han gaar over,
vilde han besidde den uden Følelse af Glæde. Endnu
klarere viser dette sig ved Sorgen. At Sorgen bestaar
i Overgangen til mindre Fuldkommenhed, ikke i selve
den mindre Fuldkommenhed, kan ingen nægte, da
Mennesket ikke kan sørge, for saa vidt han er delagtig
i nogen Fuldkommenhed.“ --- Vi have her Relativitetens
Lov i dens Anvendelse paa Følelsernes Omraade. Lyst
og Smerte føle vi kun ved en Forandring i vor Tilstand,
en Fremgang eller en Tilbagegang i en vis Henseende.
Dog bestaa de, som Spinoza ligeledes udtrykkelig fremhæver,
ikke i en Sammenligning mellem de to Tilstande,
der afløse hinanden. En saadan Sammenligning vilde
føre os ud over Følelsens Omraade til det rent intellektuelle.
--- Vi ville dog faa at se, at Spinozas Psykologi
% --- p. 131 ---
\setcounter{footnote}{0}%
\opage{131}netop lider af den Ensidighed, at denne Forskel mellem
de forskellige sjælelige Livsytringer ej fastholdes, i det
den intellektuelle Side egentlig hos ham bliver ét
og alt.

Lyst og Smerte opstaa ved Fremgang og Tilbagegang
i vort Væsen, og vort Væsen ser Spinoza som
aktivt, som Stræben efter Selvhævdelse, --- en endelig
Form for den uendelige, i alt virkende Grundkraft, der
er Aarsag til sig selv. Dette minder om den biologiske
Grundlov, der i vore Dage er bleven udtalt og gennemført
paa saa genial en Maade, og der ser Livet som
en Kamp for Tilværelsen. Spinoza har i denne Opfattelse
Forgængere --- ej blot i Descartes og de andre
Fysikere, der opstillede Inertiens Lov, --- men ogsaa i
Thomas Hobbes, der byggede hele sin etiske og politiske
Lære paa Selvopholdelsesdriften. ---

Vi have hidtil kun taget Hensyn til de passive
Lidenskaber, som opstaa hos Mennesket ved Paavirkninger
ude fra og derfor ikke ganske kunne forklares
ud af hans eget Væsen alene, i det de ydre Omgivelsers
Natur spiller en væsentlig Rolle med. Men Lidenskaben
er ej altid af passiv Natur, skønt der i den danske Betegnelse
kun er taget Hensyn til den.\footnote[1]{Spinozas almindelige Betegnelse er \textit{affectus}, og denne kan være enten passio eller actio.} Foruden den
Glæde og Attraa, der er af passiv Art, gives der andre
Følelser af Glæde og Attraa, der tilhøre os som handlende
Væsener.

Vor Handlekraft falder for Spinoza sammen med
vor Tankekraft. Han ser Aandens Væsen i Tanken
% --- p. 132 ---
\setcounter{footnote}{0}%
\opage{132}(lige som han ser Materiens Væsen i Udstrækningen),
Beslutningen eller Viljesakten er for ham kun Tilslutningen
til og Bekræftelsen af den erkendte Sandhed.
Men Tanken er ikke et stumt Billede; den er Virksomhed
og Liv. Man kan handle mod et Billede eller et Ord,
der foresvæver Bevidstheden, men ikke imod, hvad der
virkelig er blevet Aandens Ejendom og hvad den har
sluttet sig til. Forstand og Vilje ere ét. Vi handle, for
saa vidt vi have fuldkomne Ideer, d. v. s. Ideer, der
ganske kunne forklares af vort eget Væsen og derfor
ere klare og tydelige. Vi lide, for saa vidt vi have
dunkle og lemlæstede Ideer, og dette er Tilfældet, naar
vi bestemmes ude fra. Vor Handle- eller Tankekraft
kan forøges eller formindskes. Sorg er Tegn paa, at
den formindskes, og hører derfor til de passive Følelser.
Derimod \emph{kunne} Glæde og Attraa være af
aktiv Natur.

Alle Handlinger, der tilhøre Aanden som tænkende,
udspringe af Aandskraften \textit{(fortitudo)}, der fremtræder
dels som Livsmod \textit{(animositas)}, dels som Højmod
\textit{(gneerositas)}. Livsmodet er den Attraa, hvormed enhver % PRINTED AS IS: „gneerositas“ (Rettelser: l. generositas)
stræber at bevare sit Væsen efter Fornuftens Bud
alene. Ved Højmod forstaar Spinoza den Attraa, hvormed
enhver efter Fornuftens Bud stræber at hjælpe
andre Mennesker og forbinde dem med sig i Venskab.
Altsaa de Handlinger, der have den handlendes
eget Vel for Øje, henhøre under Livsmodet; de, der
ogsaa have andres Vel for Øje, til Højmodet. Til
første Art høre Dyder som Maadeholdenhed og
Aandsnærværelse, til sidste Art Beskedenhed, Mildhed
o. s. v.

% --- p. 133 ---
\setcounter{footnote}{0}%
\opage{133}I Modsætningen mellem den lidende og den handlende
Tilstand eller mellem Tankens Dunkelhed og
Klarhed have vi Udgangspunktet for en Vurdering af
Affekterne. Skønt Begreberne godt og ondt ikke betegne
noget positivt i Tingene selv, faa de dog en
relativ Betydning for os, i det vi sammenligne den
virkelige menneskelige Natur med et Forbillede, vi
danne os af den fuldkomne Menneskenatur. Godt
kalde vi saa det, der fører os dette Forbillede nærmere,
ondt, hvad der fjerner os derfra. Dyden er
Kraft --- den er nemlig selve Menneskets Natur, for
saa vidt Mennesket formaar at frembringe noget, der
kan begribes ved selve hans Naturs Love. Selvstændighed,
aandelig Frihed, aandelig Kraft er Dydens
Væsen, der, som sagt, fremtræder dels som Livsmod,
dels som Højmod.

Da Menneskets Væsen bestaar i Tanken, kan Fornuftens
Bud ikke være forskelligt fra selve Grunddriften.
Lidenskabens Kilde var jo Enkeltvæsenets
Stræben efter at hævde sig selv. Men netop denne
Stræben er Dydens første og eneste Grundlag. Der
kan ikke tænkes noget Princip, som skulde ligge dybere
end denne Stræben, der er ét med selve Grundnaturen
i Enkeltvæsenet. Enhver stræber efter det for ham
nyttige. Denne Stræben hæmmes kun ved ydre Aarsager
og er i sig selv ét med Dyden. Fornuftens Bud
er, at enhver skal søge sin sande Nytte. Fornuften er
ikke andet end selve vor Aand, for saa vidt denne har
en klar og tydelig Erkendelse. Derfor gaar Fornuftens
Bud ikke ud paa andet end paa Erkendelse, og Stræben
efter sand Erkendelse er den eneste Dyd. Vor Magt
% --- p. 134 ---
\setcounter{footnote}{0}%
\opage{134}og vor Frihed bestaar i den sande Erkendelse; vi
handle jo kun, naar vi have fuldkomne Ideer, ellers
lide vi. Men det højeste, Aanden fatter, er Gud, det
absolut uendelige Væsen, uden hvilket intet kan være
eller tænkes. At udvikle sin Erkendelse er at erkende
Gud, hans Attributer og hans Virken, der følge af hans
Natur. Saligheden bestaar i den Aandens Fred, der
udspringer af den anskuende Erkendelse af Gud. Den
højeste Attraa hos den, der ledes af Fornuften, er den,
der fører ham til klar Erkendelse af sig selv og af alle
Ting, der frembyde sig for hans Forstand. I denne
Stræben finder den menneskelige Natur sin højeste
Fuldendelse. Den udvikler sig her efter sine egne
indre Love. Der kan derfor ikke vederfares Mennesket
noget ondt uden fra ydre Aarsager eller for saa vidt
Mennesket er en Del af Naturen og paavirkes af andre
Dele. Og dette maa han som endeligt Væsen nødvendigvis
være. Men naar andre Enkeltvæsener stemme
overens med hans Natur, ville de kunne være ham til
Hjælp i hans Stræben. Intet er saa nyttigt for os som
det, der ganske stemmer med vor Natur. Naar to
Individer af samme Natur indbyrdes forbindes, danne
de et Individ, der er dobbelt saa stort som det enkelte
Individ. Mennesket kan til sit Væsens Bevarelse ikke
ønske noget bedre, end at alle stemme saaledes overens
i alt, at de danne lige som én Aand og ét Legeme,
og at alle paa én Gang søge alles fælles Nytte. De,
der ledes af Fornuften, ønske derfor intet for sig selv,
som de ikke ogsaa ønske andre Mennesker.

For saa vidt Menneskene ledes af de passive Lidenskaber,
afvige de fra hverandre, og for saa vidt er ogsaa
% --- p. 135 ---
\setcounter{footnote}{0}%
\opage{135}hvert enkelt Menneske omskiftende og ubestandigt.
Naar de derimod ledes af Fornuften, stemme de overens.
Enhver virker da, hvad der er godt for den
menneskelige Natur i det hele, og derfor ogsaa, hvad
der stemmer med hvert enkelt Menneskes Natur. Ved
at søge sin egen sande Nytte befordrer hver enkelt altsaa
andres Nytte. Det højeste Gode, Erkendelsen af
det evige og den Fred, dermed følger, er fælles for
alle; fra Adgang til det er ingen udelukket. Og jo
mere enhver for sig bliver delagtig i dette Gode, des
mere vil han ogsaa ønske at gøre andre delagtige deri.
Dels er det en psykologisk Lov, at det, vi elske, vinder
endnu større Værdi for os, naar vi se, at ogsaa andre
elske det, dels ville vi fremmes i vor Stræben, naar
andre stræbe med os. Der er ingen Maade, hvorpaa
et Menneske bedre kan lægge for Dagen, hvilke aandelige
Kræfter han raader over, end ved at opdrage andre
til at leve efter deres egen Fornufts Bud. Desuden
kunne Menneskene kun ved forenede Kræfter skaffe
sig, hvad de behøve, og modstaa Farerne. Ganske vist
leve de færreste efter Fornuftens Bud, men alligevel
gælder det --- trods alt, hvad Satirikere, Melankolikere
og Teologer indvende, --- at der af det menneskelige
Samfundsliv udspringer langt flere Goder end Onder.

Fra det nu vundne Standpunkt bedømmer Spinoza
de forskellige Lidenskaber. Glæden er i og for sig
god, da den er Udtryk for en Fremgang i Fuldkommenhed.
Sorgen derimod er i og for sig ond, da den betegner
en Tilbagegang. Alt, hvad vi gøre, skal sigte
til Fremgang og Forbedring. Men saa længe vi sørge,
gøre vi os selv uskikkede til at handle saaledes: Sorgen
% --- p. 136 ---
\setcounter{footnote}{0}%
\opage{136}er en Tilbagegang i Handlekraft. --- Spinoza vurderer
derfor mange Følelser paa en anden Maade end den
sædvanlige, og han træder her navnlig i en karakteristisk
Modsætning til den teologiske Opfattelse. Af hans
klare Indsigt i den nødvendige Sammenhæng og hans
Tro paa, at denne nødvendige Sammenhæng er ét med
Guddommens Væsen, maatte følge en anden Vurdering
af forskellige Stemninger og Følelser end hos dem, der
betragte Naturen som fordærvet ved et oprindeligt Fald.
I en Udtalelse af mere personlig Karakter, end man
ellers finder i hans Skrifter, siger han: „Kun en vild
og mørk Overtro kan forbyde Glæden. Hvorfor bør
man snarere stille Sult og Tørst end fordrive Melankolien?
--- Saaledes tænker jeg, og saaledes stiller
Sagen sig for mig: Ingen Guddom, og i det hele
ingen, som ikke er misundelig, glæder sig ved min
Svaghed og mit Uheld og ansér Taarer, Hulken, Frygt
og lignende Tegn paa aandelig Svaghed for Tegn paa
Dyd. Tvært imod, jo større Glæde vi føle, des større
Fuldkommenhed gaa vi over til, og des mere have vi
Del i den guddommelige Natur. Den vise Mand vil
derfor bruge Tingene og glæde sig ved dem, altsaa
styrke og forfriske sig ved passende og behagelig Spise
og Drikke, ved Vellugt, smukke, friske Planter, pyntelig
Dragt, Musik, Legemsøvelser, Skuespil og andre saadanne
Ting, som enhver kan nyde uden Skade for
nogen anden.“

At Hadet er ondt, følger af sig selv. Det falder
bort ved den Erkendelse, at Bestræbelsen for det
menneskelige Ideals Virkeliggørelse ikke udelukker
nogen. Had forøges ved gensidigt Had, men slukkes
% --- p. 137 ---
\setcounter{footnote}{0}%
\opage{137}ved Kærlighed og Højmod og gaar selv over til en
Kærlighed, der bliver des større, jo stærkere Hadet
har været. --- Haab og Frygt ere heller ikke i sig selv
gode. De ere forbundne med Uvished og Sorg og ere
Tegn paa usikker og mangelfuld Erkendelse. Den, der
er sikker i sin Sag, hverken haaber eller frygter.
--- Medlidenhed er i sig selv slet og unyttig hos et Menneske,
der lever efter Fornuftens Bud. Med Medlidenhed
er der forbundet Sorg, og det gode, der kan udspringe
af den, udspringer ogsaa af Fornuftens Bud.
Den, der har indsét, at alt følger af den guddommelige
Naturs Nødvendighed og sker efter Naturens evige Love
og Regler, vil ikke ynkes over nogen. En saadan Stemning
fører let til Handlinger, der siden efter vise sig
at være urigtige, og man vil let blive skuffet af falske
Taarer. --- Ydmyghed og Anger ere heller ikke Dyder,
men Tegn paa Ufuldkommenhed. Ydmyghed er en
Sorg, der opstaar ved, at Mennesket betragter sin egen
Svaghed, og Angeren beror ligeledes kun paa, at Mennesket
sér sig som afmægtigt. Ingen af disse Følelser
kan derfor være Motiv til den fuldkomne Handlen,
som netop skal udspringe af Menneskets aandelige
Kraft. Vi komme snarere til den rette Vej ved Fornuft
og Kærlighed end ved Anger og Sønderknuselse.
Kun overtroiske Lærere, som bedre forstaa sig paa at
udskælde Lasten end at forkynde Dyden, som knække
Menneskenes Sind i Steden for at styrke det, og som
derfor hellere lede Menneskene ved Frygt end ved
Fornuft, kun de kunne prise disse Følelser. Den
Attraa eller Selvopholdelsesdrift, der er ét med Menneskets
fornuftige Væsen, er rettet direkte mod det gode.
% --- p. 138 ---
\setcounter{footnote}{0}%
\opage{138}Det onde undgaas indirekte. Naar vi følge vort sande
Væsen, trænge vi vel til Føde, men ikke til Medicin.
Det er kun den syge, der af Frygt for Døden tager
Medicin. Men det frie Menneske tænker ikke paa noget
mindre end paa Døden; hans Visdom har Liv, ej Død
til sit Indhold.

Dog er det kun, naar vi tale om det Menneske,
der ledes af Fornuften, om den aandelige fri, at dette
gælder. Den, som hverken bevæges til at hjælpe andre
af Fornuft eller af Medlidenhed, kaldes med Rette
umenneskelig. Og da Menneskene sjælden leve efter
Fornuftens Bud, gøre Ydmyghed og Anger, Haab og
Frygt mere Gavn end Skade, og naar man skal fejle,
maa man hellere angre og frygte for meget end for
lidet. Naar de ufornuftige, aandelig ufri Mennesker
alle vare lige overmodige, vilde de ikke undse sig for
noget og ikke frygte noget, og intet vilde kunne styre
dem. Profeterne, der sørgede, ej for nogle faa Menneskers,
men for den fælles Nytte, have derfor ivrig indskærpet
Ydmyghed, Anger og Ærbødighed. Og i
Virkeligheden kunne de, der beherskes af disse Følelser,
langt lettere end andre føres til at leve efter
Fornuftens Bud, d. v. s. være fri og føre et saligt Liv.

Spinoza angiver en Række Midler, hvorved Lidenskaben
kan føres til at blive aktiv i Steden for passiv
\textit{(remedia affectuum)}. Først og fremmest gælder det at
erkende Lidenskabens Natur og Aarsag; ti saa snart
den erkendes, hører den op at være passiv Lidenskab.
Alle passive Lidenskaber opstaa jo ved, at vi ere Dele
af Naturen og derfor kun have ufuldkomne Ideer; det
passive falder derfor bort ved den fuldkomne Erkendelse.
Det er en almindelig Erfaring, at Sorgen
% --- p. 139 ---
\setcounter{footnote}{0}%
\opage{139}mildnes, naar man indsér, at Ulykken var nødvendig.
Dog er det en og samme Lidenskab, der først fremtræder
som passiv og siden som aktiv. Den ærgerrige
og stolte bestræber sig lige saa vel som den, der ledes
af Fornuften, efter at alle skulle leve efter hans Tanker;
men hvad hin attraar af ufuldkommen Indsigt, stræber
denne efter med fuld og sand Erkendelse. --- Et andet
vigtigt Middel er at skille Sindsbevægelsen eller Lidenskaben
fra Tanken om den ydre Aarsag og forbinde
den med andre Tanker. Kærlighed og Had opstod jo
af Glæde og Sorg i Forbindelse med Tanken om de
Ting, der voldte disse Følelser; naar denne Forbindelse
ophæves igen, maa Lidenskaben altsaa ogsaa faa en
anden Karakter. Vi maa benytte de Mellemrum, hvor
Lidenskaben ikke raser, til at til Veje bringe og indøve
en saadan Forbindelse mellem Forestillingerne, som Fornuften
fordrer, altsaa navnlig føre dem alle tilbage til den
højeste Tanke, Ideen om Guddommen. Da Tanken
omfatter de fælles Love, der stedse virke i alt, ville
de Følelser, der udspringe af Tanken om det evige
stedse have Anledning til at opstaa og ville efter
Haanden faa Overhaand over de andre Lidenskaber.
Vor endelige Natur fører det desuden med sig, at vi
maa søge Næring og Styrke hos noget uden for os; derfor
kunne vi ikke befries fra Kærligheden som fra de
andre Lidenskaber. Det er i det hele en fortræffelig
Ting, at alle Lidenskaber, der ere gode, ere af en
saadan Natur, at vi uden dem ikke kunne være eller
bestaa. Jo flere Aarsager og Tilknytningspunkter en
Lidenskab har, des større Magt har den. Og dette
gælder netop om den Lidenskab, der udspringer af
% --- p. 140 ---
\setcounter{footnote}{0}%
\opage{140}Tanken om det evige og guddommelige, der ligger til
Grund for alt. En klar og tydelig Erkendelse fører
Glæde med sig, fordi vor Handlekraft naar sin Fuldendelse
i den. Og da denne Glæde, naar Erkendelsen
gaar tilbage til den første Aarsag, har Gud til Aarsag,
opstaar Kærligheden til Gud, som næres ved alt, hvad
vi erkende, og derfor maa være den mægtigste Lidenskab
af alle. Selv om Kærligheden til det uforanderlige
og evige ikke ganske kan fortrænge Lidenskaberne,
vil den dog indskrænke dem til den mindste Del af
Sjælen.

Den intellektuelle Kærlighed --- saaledes kalder
Spinoza (efter Descartes og ældre jødiske Religionsfilosofer)
den af sand Erkendelse udspringende Glæde
ved det evige og uendelige, --- som er den forklarede
Lidenskab, viser os Menneskets Væsen fra en hel anden
Side end den, vi hidtil have betragtet det fra: ikke længer
som et Enkeltvæsen, men som et Væsen, der ved
sin Enhed med det uendelige Grundvæsen faar Bestaaen
i sig selv og er evig frit. Men denne intellektuelle
Kærlighed er ikke blot et Maal, Mennesket nærmer
sig hen imod og stræber at naa. Naar Gudsbegrebet
tænkes efter Systemets strenge Konsekvents, stiller det
sig anderledes. --- Som absolut Væsen kan Gud ikke
gaa over til større eller mindre Fuldkommenhed og kan
altsaa ikke føle Kærlighed eller Had til nogen. Den,
der vilde ønske, at Gud elskede ham igen, vilde ønske,
at Gud ikke var Gud. At Gud ikke elsker Menneskene,
maa dog ej forstaas, som om de vare ham ligegyldige.
Mennesket og alt, hvad der er, er saaledes i Gud, og
Gud bestaar saaledes af alle Væsener, at han ikke kan
% --- p. 141 ---
\setcounter{footnote}{0}%
\opage{141}føle Kærlighed til noget eller nogen, da der intet gives
uden for hans eget Væsen. Menneskets Kærlighed til
Gud er en Del af den uendelige Kærlighed, hvormed
Gud omfatter sit eget Væsen. Mennesket betragter
sig i den som Led i og Form for Guds uendelige
Væsen. Da dette ikke kan træde frem i nogen bestemt
Form, „som en Genstand blandt andre flere“, faar den
intellektuelle Kærlighed ingen Genstand, bliver egentlig
ikke Kærlighed til noget. Derfor kan der heller ikke
være Tale om Genkærlighed; Kærlighedens Subjekt og
Objekt ere ikke forskellige. Det er, kan man sige, en
immanent Kærlighed; den har samme. Karakter som
Grundvæsenets Virken i det hele: Gud kan kun være
immanent (reflexiv) Aarsag, ikke virke paa noget uden
for sig. --- Denne Lære er ej blot konsekvent, men
indeholder ogsaa den store Sandhed, at enhver ideal
Stræben har Maalet, ikke uden for sig, men i sig. Saligheden,
siger Spinoza, er ikke Dydens Belønning, men
er selve Dyden. For den mytologiske Fantasi stiller
det sig saaledes, at Arbejdet er ét, Lønnen et andet.
Men efter Spinozas Anskuelse gælder det, hvad der
staar i Hejbergs Digt:

I det du søger, du ham alt har fundet,

I det du spørger, er alt Svaret vundet.\footnote[1]{„Naar dette forholder sig saaledes, kunne vi med Rette erklære det for en stor Urimelighed, hvad mange, som man ellers anser for store Teologer, sige, at naar der af Kærligheden til Gud ikke fulgte noget evigt Liv, saa skulde man helst søge sit eget Jeg, som om man deri kunde finde noget bedre end i Gud. Dette vilde være lige saa taabeligt, som hvis en Fisk, der jo ikke kan leve uden i Vandet, vilde sige: Naar der for mig efter dette Liv i Vandet ikke følger noget evigt Liv, saa vil jeg op af Vandet, op paa Land.“ \emph{Spinoza}: „Kortfattet Afhandling“ II, 26. (Schaarschmidts tyske Oversættelse).}

% --- p. 142 ---
\setcounter{footnote}{0}%
\opage{142}Den intellektuelle Kærlighed maa, lige som den
sande Erkendelse, være evig, har hverken Begyndelse
eller Ende. Mennesket er jo i den ét med Grundvæsenet
og altsaa evigt som dette, og det beror kun
paa en ufuldkommen Erkendelse, naar han opfattes
som værende i Tiden. Det gør intet Skaar i denne
Evighed, at vi ikke kunne erindre at have været til
før vor Existents i Tiden. Erindring forudsætter Tidsforhold,
men det evige har slet intet med Tiden at
gøre. Den sande Erkendelse, der ser alt under Evighedens
Synspunkt, er hævet over Sanseforestilling og
Erindring, som kun tilhøre Tiden. De fleste formaa
rigtig nok ikke at naa denne Erkendelse. Lige som de
ved Gud forestille sig en enkelt Genstand, saaledes
forstaa de ved Aandens Evighed en ligefrem Fortsættelse
og Vedvaren af deres Existents i Tiden. Men
Aandens Evighed betegner ikke andet end dens Natur
som nødvendigt Element i Grundvæsenet (evig Tankeform,
\textit{æternus cogitandi modus)}.

I denne mystiske Lære om den intellektuelle Kærlighed
bøjer Spinozas Filosofi tilbage mod sit Hovedprincip.
Efter at den experimentale Understrøm har
været den stærkeste i Løbet af de psykologiske og etiske
Undersøgelser, træder nu den religiøse Spekulation
atter i Forgrunden. Modsætningen mellem de to Strømninger
gør sig end ogsaa her skarpt gældende, i det
Spinoza afskærer Muligheden af at naa det ideale
Standpunkt gennem den endelige Udvikling. Lige som
% --- p. 143 ---
\setcounter{footnote}{0}%
\opage{143}han i en tidligere Sammenhæng nægtede, at Enkeltvæsenernes
individuelle Mangfoldighed kunde udledes
af Grundvæsenet, saaledes nægter han nu, at Enkeltvæsenet
gennem en Udvikling i Tiden kan naa Erkendelsen
af det evige. Hele den Befrielsesproces, vi have
set Mennesket gennemgaa, beror altsaa paa en Illusion\footnote[1]{Dette er et af de Punkter, som \emph{Brøchner} med stor Klarhed har belyst i sin Monografi „Benedikt Spinoza“ (Kbhvn. 1857). Kap. 6 og 8.},
den samme Illusion, der er ejendommelig for
den laveste Erkendelsesform: den nemlig, at den i
Erfaringen givne, i Tiden sig udfoldende Tilværelse
har Realitet. I Virkeligheden er Grundvæsenet det
ene existerende, og Mennesket er som Element i dets
Natur selv evigt, --- kan altsaa ej være passive Lidenskaber
underkastet og ingen Udvikling gennemgaa.

Spinoza forlader derfor ogsaa her sine tidligere
psykologiske Bestemmelser. Han opfattede jo Glæden
som Overgang til større Fuldkommenhed, og vi fremhævede
det sande i denne Bestemmelse. Men den er
ham ej længer nok. I det absolute og evige kan der
ikke være Tale om Overgang: derfor siger han, at i
Steden for Glæden \textit{(lætitia)} træder nu Saligheden
\textit{(beatitudo)}, som bestaar i, at Aanden er i Besiddelse
af selve Fuldkommenheden. --- Det er intet Under, at
Spinoza her har vanskeligt ved at gøre sig forstaaelig,
og vi ved at følge ham. En fuldkommen Afslutning
og Fuldendelse strider mod Psykologien. Den kender,
lige saa vel som Fysiken og Fysiologien, kun relative
Tilstandsformer. Det har Spinozas samtidige, Thomas
Hobbes, klart set og udtrykt saaledes (\textit{De homine}
% --- p. 144 ---
\setcounter{footnote}{0}%
\opage{144}XI, 15): „Naar der skal være et sidste Maal, er der,
naar det er naaet, intet at tragte efter og attraa; saa
er der heller intet, der staar for Mennesket som et
Gode, ja han kan ikke engang længer have nogen Følelse,
ti enhver Følelse er forbunden med Attraa eller
Afsky, og ikke at føle er det samme som ikke at leve.
Det største Gode bestaar derimod i, saa uhindret som
muligt at gaa fremad til stedse højere Formaal. Livet
er en stadig Bevægelse, og naar denne Bevægelse ej
kan gaa i lige Linje, gaar den i en Kreds.“ --- Man
vil let se, at vi her igen staa ved det flere Gange omtalte,
at al deduktiv Spekulation egentlig gør Etiken
umulig, da det evige Grundvæsen er det ene værende.
Relativitetsloven derimod viser os, at vort Liv som alt,
hvad vi erkende, er betinget, er Led i en stor Udvikling.
Arbejde og Stræben ere altsaa ikke Illusioner,
men Udtryk for en almengyldig Lov. En Tilværelse
med idel Væren uden Vorden, er os uforstaaelig.

Ogsaa Platon, med hvem Spinoza i det hele har
saa meget til fælles, stansede ved en uovervindelig
Modsætning mellem Vorden og Væren. Idealverdenen
i sin Evighed var for ham den sande Væren, og den i
Erfaringen givne Tilværelse, hvor Vorden, Udvikling,
Stræben, altsaa ogsaa det etiske Arbejde hører hjemme,
blev til sidst et uforklarligt Skin. Derfor var det
højeste for Platon, lige som for Spinoza, Teorien, den
spekulative Anskuen af det evige. Og dog udmærker
Platon, lige som Spinoza, sig ved sin dybe psykologiske
Forstaaelse af Driftens Væsen.

Spinoza minder her ogsaa om den religiøse Mystik,
der blomstrede i Slutningen af det 17de Aarhundrede,
% --- p. 145 ---
\setcounter{footnote}{0}%
\opage{145}særlig inden for den kartesianske Skole (Fénélon, Malebranche).
Han er paa sin Vis en Forkæmper for den
„uinteresserede Kærlighed“, og staar her Side om Side
med den ædle Erkebiskop af Cambray. Ogsaa de religiøse
Mystikere søge en umiddelbar Forening med det
evige og hige efter et Standpunkt, hvor Tiden og Endeligheden
forsvinder. Hvad der adskiller Spinoza fra
dem (afset fra Modsætningen mellem hans Filosofi og
deres Teologi), er, at de væsentlig bevæge sig paa Følelsens
Omraade, han derimod paa Tankens. Deres
uinteresserede Kærlighed til det evige er sentimental,\footnote[1]{Dette Ord brugt i sin egentlige, ej i dadlende Betydning.}
hans er intellektuel. Han har, som Jakobi træffende
sagde om ham, sin Himmel i Forstanden. Det hører
til hans Psykologis Ensidighed, at han ser Aandens
Væsen udtømt i Tanken. Viljen er kun Tilslutningen
til den erkendte Sandhed, Følelsen kun en forvirret
Tanke. I sin Lære om Lidenskaberne viser han,
hvorledes disse ved en Metamorfose gaa over til at
blive Tanker. Den rene intellektuelle Virksomhed er for
ham Aandens egentlige Væsen. Det var først hen
imod Slutningen af det 18de Aarhundrede, at Psykologien
ret blev klar over Følelsens Ejendommelighed og
Selvstændighed som sjæleligt Fænomen ved Siden af
Erkendelse og Vilje. Det var sikkert det 18de Aarhundredes
Fremdragen af Individualiteten, dens Betydning
og dens Ret, der her aabnede Øjnene. Følelsen
har en mere individuel Karakter end de andre aandelige
Livsytringer. Vi kunne til en vis Grad tænke og handle
som andre, men kun føle som os selv.
% --- p. 146 ---
\setcounter{footnote}{0}%
\opage{146}Allen gehört was du denkst, dein eigen ist nur,

\begin{center}was du fühlest!\end{center}


(Schiller).

Spinoza kunde vel fordybe sig i Betragtning af de
individuelle Sindsbevægelser og deres Love, men han
søgte dog stedse ud over dette individuelle til den dybe
Enhed og den store Helhed, hvori alt individuelt for
ham udstrømmer, i det det naar sin Fuldendelse. Hans
personlige Karakter synes derhos at have været afgjort
kontemplativ; at bevæges af Haab og Frygt, af de mange
skiftende Stemninger, hvorunder de fleste Menneskers
personlige Liv udvikler sig og maa udvikle sig, var ham
fremmed, fordi han én Gang for alle havde viet sig til
Erkendelsen og sammensluttet sin personlige Udvikling
i Tankens Udvikling. Følelsen fremtræder hos ham
som det milde Skær, den stille Glæde, vunden Erkendelse
udbreder; de stille Kyster, hvor han færdedes,
hjemsøgtes ej af Livets voldsommere Brændinger.

% --- p. 147 ---
\setcounter{footnote}{0}%
\opage{147}\begin{center}V.\end{center}


Saa ensom og tilbagetrukken en Stilling Spinoza
end indtager i sin Tænkning og i sit Liv, bære dog
alle hans Skrifter Præg af, at han følte sig som „Menneske
blandt Mennesker“. Den Opgave, han havde sat
sig, og hvis Løsning han ansaa for det højeste Gode,
Erkendelsen af vor Aands Sammenhæng med hele Naturen,
kunde efter hans Overbevisning ikke naas af den
enkelte for sig alene, men kun i Samfund med andre.
Opgavens Løsning forudsætter derfor dels en Erkendelsens
Reformation, dels Dannelsen af et Samfund, i
hvilket saa mange som muligt opdrages og udvikles til
det angivne Maal. I hans etiske Lære se vi, hvorledes
den Aandskraft, som for ham betegner Karakterens
højeste Udvikling, fremtræder ej blot som Livsmod
(i Individets egne Forhold), men ogsaa som Højmod
(i Forholdet til andre). --- Den etiske Karakter forudsætter
en Udvikling, der ikke er til Stede hos alle, og
naar Spørgsmaalet derfor bliver om Grundlaget for det
menneskelige Samfund, maa vi gaa længere tilbage,
tilbage til det samme Punkt, hvorfra den etiske Udvikling
begyndte: Mennesket som rent Naturvæsen, passive
Lidenskabers Herredømme underkastet. Statens Væsen
og Oprindelse kan ikke forstaas, uden man forstaar
Lidenskabens Natur. Spinoza fremhæver derfor her,
% --- p. 148 ---
\setcounter{footnote}{0}%
\opage{148}lige som i Etiken, Nødvendigheden af at studere Lidenskaben
som ethvert andet Naturfænomen. Statsmændene,
siger han, have skrevet langt bedre om Politik end Filosoferne,
fordi hine have fulgt Erfaringen, medens disse
have lastet Lidenskaberne og konstrueret en Stat for
rene Fornuftsvæsener, som ikke findes, eller som i hvert
Tilfælde ingen Love behøve. Erfaringen frembyder
saa mange forskellige Statsformer, at vi vanskelig kunne
udtænke noget aldeles nyt, som det praktiske Livs nødvendige
Krav ikke skulde have bragt for Dagen et eller
andet Sted. --- Da Menneskene væsentlig ledes af Lidenskaberne,
over for hvilke Fornuftens og Religionens
Bud kun formaa lidet, maa Staten ordnes saaledes, at
intet ganske overlades til de enkeltes Troskab og Moralitet,
men at de herskende ved deres egen Interesse
hindres i at handle uredelig, hvilket Sindelag der saa
maatte ligge til Grund for deres Handlen. Aandens
Frihed og Kraft er en privat Dyd; og i Staten, hvor
det først og fremmest gælder Sikkerhed for alle, tør
man ikke regne paa den.

Ethvert Naturvæsen er en endelig Form for det
uendelige Grundvæsen, og den Kraft, hvormed det existerer,
virker og hævder sig i sin Tilværelse, er en Del
af det guddommelige Væsens uendelige Kraft. Derfor
har ethvert Naturvæsen oprindelig --- hvad enten det
ledes af Fornuft eller Lidenskab --- lige saa megen
Ret, som det har Magt. Den naturlige Ret er ét med
den naturlige Magt eller med den Naturlov, der oprindelig
bestemmer Individets Handlen. Det er altsaa
i Følge en naturlig Ret, at de store Fisk fortære de
smaa. For enhvers Vedkommende bestemmes, hvad
% --- p. 149 ---
\setcounter{footnote}{0}%
\opage{149}han anser for godt og ondt, ved hans Interesse. Han
holder f. Ex. kun et givet Løfte, saa vidt det stemmer
med hans egen Fordel. Naar en Røver har tvunget
mig til at love at afstaa ham min Formue, betænker
jeg mig ej paa at bryde dette Løfte, naar jeg kan gøre
det uden Fare. I Naturtilstanden er intet forbudt,
uden hvad ingen falder paa at gøre. Da nu Lidenskaberne
--- især Vrede, Had og Misundelse --- bringe
Menneskene i Strid med hinanden, ere de af Naturen
Fjender. Den ene maa stedse frygte den anden. Under
disse Forhold existerer ingen Ejendom --- enhver besidder
kun, hvad han har Magt til at forsvare --- og
ingen Sikkerhed. Men derfor er den naturlige Ret
egentlig slet ikke til, saa længe den bestemmes ved
enhvers Magt. Menneskene kunne kun ved forenede
Kræfter opholde deres Liv og uddanne deres Aand.
Man kan derfor sige, at Mennesket er et selskabeligt
Væsen: hans naturlige Vilkaar lede ham til Stiftelse af
et Samfundsliv. Jo flere der forene deres Kræfter, des
større Magt, Sikkerhed og Ret besidde de.

Naar et saadant Samfund stiftes, overgive de enkelte
til Dels deres Ret til Samfundsmagten \textit{(imperium)}.
Hvad Samfundsmagten erklærer for ret og uret, maa
de enkelte gøre eller undlade. Hver enkelt har nu
kun den Ret, Samfundsmagten tilstaar ham. Men hvor
vidt strækker Samfundets Ret sig? Kun saa vidt som
dets virkelige Magt strækker sig. Samfundet hviler
paa det Princip, at en Lidenskab kun kan hæmmes
ved en anden, stærkere Lidenskab, og at Mennesket
derfor maa holdes fra at paaføre andre Skade ved
Frygten for selv at lide større Skade. Hvad Samfundet
% --- p. 150 ---
\setcounter{footnote}{0}%
\opage{150}ikke kan tvinge Mennesket til at gøre, har det ingen
Ret til at befale. Samfundsmagten maa derfor i Følge
Selvopholdelsens Lov, som er ét med Fornuftens Bud,
ikke befale noget, som vilde gøre Ende paa dens eget
Herredømme. Lige som det enkelte Individ er stærkest,
naar det følger Fornuften, saaledes kan kun det
Samfund bestaa, der bestemmes af Fornuften. At sige,
at Samfundsmagten har Lov at gøre, hvad den vil, er
det samme som at sige, at et Menneske har Lov at
gaa fra Forstanden. Den er bleven til for at hjælpe
de enkelte til deres naturlige Ret og maa derfor søge
det, der er nyttigst for alle. Hvad der strider mod
den menneskelige Natur og oprører den menneskelige
Følelse, kan Statsmagten ikke befale uden at stile mod
sin egen Undergang. Den bliver da selv et Onde for
Mennesket af samme Art som de, han led under i
Naturtilstanden; de enkelte ville søge Midler til at
sikre sig imod den, --- og Samfundet vil altsaa opløses.
Navnlig vilde Staten overskride sin Ret, hvis den vilde
gribe ind paa den fri Overbevisnings Omraade, hvor
Trusler og Løfter intet formaa. Ingen kan fraskrive
sig sin Dømmekraft eller hindre en anden i at bruge
hans Dømmekraft. Trods alle Forbud og al Tvang
ville Menneskene elske og hade, som det stemmer med
deres Natur.

Der er altsaa ingen absolut Modsætning mellem
Livet i Naturtilstanden og Livet i Samfundet. Spinoza
adskiller sig i dette Punkt fra sin samtidige Thomas
Hobbes, der lærte, at Selvopholdelsesdriften, som fører
Menneskene ud over Naturtilstanden, denne Krig af
alle mod alle, byder fuldstændig Underkastelse under
% --- p. 151 ---
\setcounter{footnote}{0}%
\opage{151}Statsmagtens Vilje som det eneste Middel til at sikre
Freden. Spinoza indrømmer, at ethvert Medlem af
Samfundet ganske vist maa opgive Retten til at leve
efter sit eget Hoved; i Steden for sin egen Afgørelse
af, hvad der er godt og ondt, maa han nu følge Statens.
Han kan heller ikke have Ret til at fortolke Lovene
paa sin egen Vis, ti saa vilde han være sin egen
Dommer, og det var jo netop for at faa en fælles
Maalestok for Ret og Uret, at Staten stiftedes. Men
den Resignation, som her maa vises, strider ikke mod
Fornuften. Selv om jeg i nogle Tilfælde maa gøre,
hvad jeg selv ikke ansér for rigtigt, vælger jeg dog det
mindste af to Onder: det vilde være et langt større
Onde, om Samfundslivet opløstes og Naturtilstanden
begyndte igen!

Statens Oprindelse er altsaa i Følge Spinoza at
søge i den menneskelige Egoisme. For at undgaa den
voldsomme Brydning af Interesserne overdrager man
Samfundet som Helhed Varetagelsen af det fælles Tarv.
Men Spinoza sér ikke Statens Opgave udtømt i den
blotte Sikkerhed. Dens Formaal er ikke blot at herske
og ved Frygt holde den ene fra at skade den anden.
Om en Stat, hvis Borgere kun ved Frygt holdes fra at
gribe til Vaaben mod hverandre, kan man snarere sige,
at den er uden Krig, end at den har Fred. Fred er
ej blot Fraværelse af Krig, men den sande Fred fremgaar
af aandelig Kraft og af Tilslutning til Statens
fælles Love. Det er netop Statens Opgave ganske at
forjage Frygten. Staten gaar ikke ud paa at gøre
Menneskene til Maskiner eller ufornuftige Dyr, men
paa at gøre et aandeligt Liv muligt. Statens Maal er
% --- p. 152 ---
\setcounter{footnote}{0}%
\opage{152}derfor Frihed. At dette Maal ikke naas, maa ikke
saa meget tilskrives Borgernes Slethed, som Forfatningens
Daarlighed. Hvor der hersker Splid, Oprør og
Foragt for Lovene, maa det komme af, at Institutionerne
ikke due. Menneskene fødes nemlig ikke som
Borgere, men maa udvikle sig til at blive det.

Spinoza ansér Demokratiet baade for den naturligste
og den fuldkomneste Statsforfatning. Det er naturligst,
at Menneskene ikke strax overdrage hele Magten
til en enkelt eller nogle enkelte, men til hele den samlede
Mængde, som bliver enig om at skabe et Samfund.
Hver enkelt faar da, i det han selv hører til Mængden,
Del i de fremtidige Raadslagninger om Statsmagtens
Brug. Derved bliver tillige Haab snarere end Frygt
Samfundets Grundlag. Den fuldkomneste Forfatning
er Demokratiet, fordi det nærmer sig mest til den Frihed,
der af Naturen tilkommer enhver.

Spinoza er den første nyere politiske Forfatter, der
har erklæret sig for en demokratisk Republik. De republikanske
Forfattere før ham i det 16de og 17de
Aarhundrede vare mere eller mindre Aristokrater. Han
er i denne Henseende Rousseau's Forgænger. Dog indrømmer
han, at hvis Patricierne lode sig lede, ikke af
deres Lidenskaber, men af Hensynet til det offentlige
Vel, vilde en aristokratisk Forfatning være allerbedst;
men en saadan Forudsætning ansér han for stridende
mod Erfaringen. Des værre døde han, just som han i
sin politiske Afhandling havde begyndt paa at fremstille
den demokratiske Forfatning, saa at vi ikke kunne se,
hvilke specielle Institutioner han har ansét for nødven
% --- p. 153 ---
\setcounter{footnote}{0}%
\opage{153}dige under den.\footnote[1]{Den sidste Paragraf i den ufuldendte Afhandling, altsaa vistnok det sidste, Spinoza har skrevet før sin Død, handler om, hvor vidt Kvinderne i en demokratisk Republik bør have lige Ret med Mændene. Hvis Kvinderne af Naturen vare Mændene ligestillede, vilde --- mener Spinoza --- Historien have frembudt Exempler paa en saadan Ligeberettigelse. Desuden kunne Kvinderne ikke uden store Ulemper faa Del i den politiske Magt: Elskovsdriften og de Lidenskaber, den fører med sig (S kinsyge o. s. v.), ville da gribe forstyrrende ind. --- Den første af disse Indvendinger har Stuart Mill grundig gendrevet i sit bekendte Skrift den anden derimod omtaler han slet ikke.} Derimod fremstiller han udførlig
Aristokratiets og Monarkiets Institutioner. Han søger
at paavise, hvorledes baade de herskende og de beherskede
--- med eller mod deres Vilje --- kunne føres
til at handle efter Fornuftens Bud. Et absolut Monarki
er umuligt, da Kongen behøver Raadgivere, af
hvilke han stedse vil være mere eller mindre afhængig,
og Forfatningen vil altsaa i Virkeligheden blive aristokratisk.
Det Kongedømme, Spinoza beskriver, svarer
til, hvad vi kalde et konstitutionelt Monarki \textit{(imperium
monarchicum, qvod a libera multitudine instituitur)}.
Kongen er endogsaa bunden til at vælge mellem de
Forslag, der have faaet flest Stemmer i det store Raad,
hvori hele Folket er repræsenteret. Den dømmende
Magt udøves af et særskilt Kollegium. Et absolut
Eneherredømme vilde, dersom det var muligt, maaske
føre til større Sikkerhed, og Statshemmeligheder vilde
bedre kunne bevares under en saadan Forfatning; men
disse Fordele vilde købes med Opgivelsen af langt
væsentligere Goder. Naar Eneherskeren ikke har
andet at tage Hensyn til end sin egen Lyst, har man
% --- p. 154 ---
\setcounter{footnote}{0}%
\opage{154}ingen Garanti for, at han stræber efter det almindelige
Vel, hvorimod han naturligen vil være en Fjende af
alle fremragende Mænd, som kunde synes at gøre ham
Rangen stridig. Efter at Spinoza har omtalt de Former
og Institutioner, hvorved Kongemagten bør indskrænkes,
bemærker han, at hans Fremstilling maaske
vil vække Latter hos dem, der ere vante til at tænke
sig alle Laster paa Mængdens Side og alle Dyder hos
de fornemme. Men den menneskelige Natur er en og
samme hos alle, og netop de herskende ere ved deres
Stilling udsatte for mange Fejl.

Udgangspunktet for Spinozas Statslære\footnote[1]{Spinozas Statslære er --- ej blot for sin Tid --- ene staaende i Klarhed, Konsekvents og Frisindethed. Men vistnok paa Grund af hans isolerede Stilling maatte han paa dette Omraade, saa vel som i Psykologien, træde i Skygge for Hobbes, skønt dennes parodoxe og kortsynede Pessimisme, i det mindste ved første Øjekast, syntes at gøre ham til Fremskridtets Modstander.}% PRINTED AS IS: „parodoxe“
 er Individets
naturlige Ret, der ikke skal opgives, men netop
befæstes og udvikles i Samfundslivet. Han er her en
Forgænger for det 18de Aarhundrede og dets Kamp
for Individets Rettigheder. Kant's Sætning om Frihed

--- i Betydning af Uafhængighed af enhver andens
tvingende Vilkaarlighed --- som den ene, oprindelige
Menneskerettighed, og den Gennemførelse, som Wilhelm
von Humboldt og Stuart Mill have givet denne Sætning,
er i Principet udtalt hos Spinoza. Og hvad der overrasker
hos en spekulativ og abstrakt Filosof: Gennemførelsen
af hans Lære har overalt en praktisk og real
Karaktér og bygges paa skarp psykologisk Opfattelse.

Men uagtet Spinoza ikke anerkender nogen fuld
% --- p. 155 ---
\setcounter{footnote}{0}%
\opage{155}stændig Modsætning mellem Naturtilstanden og Samfundslivet,
betragter han dem dog som forskellige og
lader Overgangen fra den ene til den anden ske ved
en Overenskomst. Det var en Forestilling, som i det
17de og 18de Aarhundrede ofte lagdes til Grund, og
som væsentlig udsprang af mangelfuld Indsigt i den
menneskelige (og al) Udviklings Love. Lige som man
i Geologien tidligere tænkte sig Jordens nuværende
Tilstand som Resultat af en Række store Revolutioner
og ganske forskellig fra tidligere Tilstande, saaledes
tænkte man sig Overgangen mellem Menneskelivets
forskellige Tilstande som en vilkaarlig Akt. Saadanne
bratte Overgange indrømmer en mere udviklet psykologisk
og historisk Indsigt ikke. Kontinuitetens Lov
hersker ogsaa her. Trods alle indgribende Begivenheder
og Katastrofer har Samfundslivet udviklet sig
væsentlig under de samme Love, som det nu beror paa,
ved en stille, ofte umærkelig Væxt.

--- Der er en nøje Forbindelse mellem Spinozas Statslære
og hans Lære om Religionen. I sin „teologiskpolitiske
Afhandling“ søgte han paa Grundlag af sin
Lære om Forholdet mellem Tro og Viden, Teologi og
Filosofi at vise Religionsfrihedens Mulighed og Nødvendighed.
Og i Statslæren have vi set, at der er et
Omraade, hvor Samfundsmagten ikke kan trænge ind,
nemlig den enkeltes frie Overbevisning. Hans ydre
Handlinger kan Staten bestemme, men Religionen tilhører
væsentlig det indre og bestaar i Erkendelsen af
det guddommelige og i Kærlighed til Næsten. De
ydre Kultusformer ere ligegyldige, og den religiøse vil
% --- p. 156 ---
\setcounter{footnote}{0}%
\opage{156}netop vise sin Kærlighed ved at bevare Freden i Samfundet
og ikke vække Splid ved sekterisk Optræden.

Den Religiøsitet, som svarer til Spinozas Standpunkt,
bestemmer han selv saaledes: „Alt hvad vi
attraa og udføre, for saa vidt vi ledes af Ideen om
Gud, henfører jeg til Religionen. Den Lyst til at gøre
vel, der opstaar af, at vi leve efter Fornuftens Bud,
kalder jeg Fromhed.“ Men da dette filosofiske Begreb
om Religionen strider mod den teologiske Opfattelse
af den, søger han nærmere at bestemme Forholdet
mellem Teologi og Filosofi. Det teologiske Religionsbegreb
hviler paa Ideen om en Aabenbaring. Her
tales naturligvis ikke om Aabenbaring i Ordets videre
Betydning, hvor al Erkendelse, vi vinde af Naturen,
kan kaldes Aabenbaring, selv om vi kun bruge „det
naturlige Lys“. Aabenbaring, i den strenge Betydning
af Ordet, er forskellig fra Fornufterkendelsen: den
kommer ej som denne gennem klare og tydelige Begreber,
men sker ved Ord og Tegn. Ikke alle kunne
modtage saadanne Ord og Tegn og meddele dem til
andre; dertil udkræves Individer med en særegen Begavelse,
som kaldes Profeter. Men denne særegne Begavelse
kan ikke bestaa i fremragende Forstand og Kundskab,
da det jo ej er rene Tanker, men Billeder og Symboler,
der skulle opfattes og meddeles. Derimod gælder det
navnlig om en stærk og levende Indbildningskraft. Der
er en naturlig Modsætning mellem saadanne Individer,
hos hvem Indbildningskraften, og saadanne, hos hvem
Forstanden er mest udviklet; den ene af disse Evner
udvikles gerne paa den andens Bekostning. Derfor er
det en stor Misforstaaelse at søge Kundskab om
% --- p. 157 ---
\setcounter{footnote}{0}%
\opage{157}Naturens og Aandens Væsen i Profeternes Skrifter.
En saadan Kundskab maa vindes paa en hel anden
Maade. Hvor skulde ogsaa en Kriger som Josva forstaa
sig paa Astronomi eller Noah paa Geografi? Ja.
Adam vidste jo ej engang, at Gud var alvidende, siden
han skjulte sig mellem Havens Træer; ej heller vidste
Abraham det, siden han bad Gud regne efter, hvor
mange retfærdige der vari Sodoma. Profeterne maatte
desuden lempe sig efter dem, de talte til. Kun om
Kristus gælder det, at han ikke beherskedes af Indbildningskraften,
men modtog Aabenbaringen ved umiddelbar
og sand Erkendelse. Det er, som vi have hørt
Spinoza sige i et af sine Breve (oven for P. 49), ej
nødvendigt at kende den historiske Kristus („Kristus
efter Kødet“), men „Guds Søn“ i Betydning af den
evige Visdom, der aabenbarer sig i alle Ting, mest i
den menneskelige Aand og allermest i Kristus, der
lærte Menneskene den universelle Fornuftlov at kende.
I Bjergprædikenen ser Spinoza, lige som den nyere
Rationalisme, Indbegrebet af den sande Kristendom.
Mose Lov var en national og politisk Institution, der
bortfaldt ved det jødiske Riges Undergang. Jøderne
have intet at bryste sig af over for andre Folk. Ogsaa
andre Nationer end Jøderne have haft Profeter; det er
ikke noget særeget for Jødefolket.

Aabenbaringen faar en forskellig Karakter hos de
forskellige Profeter, i det deres Personlighed, Legemsbeskaffenhed
og tidligere Meninger nødvendigvis faa
Indflydelse paa deres Forkyndelse. Hofmanden Esajas
taler anderledes end Hyrden Amos. Man kan derfor
ikke uden videre benytte den enes Skrifter til at for
% --- p. 158 ---
\setcounter{footnote}{0}%
\opage{158}tolke den andens. Man maa paa Grundlag af sikker
Sprogkundskab og Kendskab til Forfatterens Liv, de
historiske Forhold og Textens Ægthed finde den virkelige
Mening af hvert enkelt Sted. Om denne Mening
i sig selv er sand, er et ganske andet Spørgsmaal. Vi
kunne ikke uden videre gaa ud fra, at Skriften stemmer
med Fornuften; det vilde føre til den Konsekvents, at
den sidste kirkelige Avtoritet maatte søges hos Filosoferne,
men denne Avtoritet vilde næppe møde stor
Respekt. Hvad Skriften lærer, og hvor vidt det, den
lærer, er sandt, er altsaa to helt forskellige Spørgsmaal.


Spinoza, som er den første, der har taget Ordet
for en rent historisk Fortolkning af Skriften, er ogsaa
den første, der har anlagt en historisk Kritik paa de
bibelske Bøger. Han hævder, at hvert enkelt Skrifts
Ægthed maa bevises for sig, og kommer til det Resultat,
at de gammeltestamentlige Skrifter i deres nuværende
Form ikke kunne have de Forfattere, hvis
Navne de bære.

De Vanskeligheder, hvorved den videnskabelige
Fortolkning af Skriften stanser, behøve dog i Følge
Spinoza ikke at besvære dem, der søge Opbyggelse i
Skriften. Det opbyggelige \textit{(pietatis documenta)} i Skriften
vil altid let kunne forstaas. Og Profeternes egentlige
Øjemed var at forkynde en Lære, der kunde føre
til Fromhed og til Lydighed mod Gud. Derfor fremstille
de Gud som en Lovgiver og Fyrste, hvem Menneskene
maa lyde. Teologien adskiller sig netop fra Filosofien
ved at lære, at Menneskene frelses ved Lydighed, medens
Filosofien hævder, at den sande Dyd udspringer
% --- p. 159 ---
\setcounter{footnote}{0}%
\opage{159}af Indsigt i Naturens Sammenhæng. Teologien søger
Fromhed, Filosofien søger Sandhed. Sandheden er i
sig selv hverken from eller ufrom. Kun for saa vidt
kan man sige, at der er Fromhed eller Ugudelighed i
et Menneskes Meninger, som han af dem tager Anledning
til Lydighed eller Ulydighed. Den, der ved at
tro det sande bliver hovmodig, har en ugudelig, den,
der ved at tro det falske bliver lydig, en from Tro.
Den Tro, der skal være fælles for alle Mennesker, kan
derfor ikke indbefatte noget, hvorom hæderlige Mennesker
kunne være uenige. Naar der forlanges, at man
skal tro paa Gud, saa maa der ved Gud forstaas et
idealt Forbillede for det menneskelige Liv \textit{(veræ vitæ
exemplar)}, men det er ligegyldigt, hvorledes enhver
tænker sig Gud --- som Ild eller Lys, Aand, Tanke,
paa panteistisk eller teologisk Vis o. s. v. Her kommer
den forskellige intellektuelle Begavelse og Udvikling til
at spille en Rolle. Alle kunne lyde og føre et Liv i
Retfærdighed og Kærlighed, men ikke alle kunne erhverve
sig en fornuftbegrundet Erkendelse.

Ved at henvise den positive Religion til det praktiske
Omraade søger Spinoza at sikre Videnskabens
Frihed. Han overlader Teologien Fromheden, men
forbeholder Filosofien Sandheden, en Fordeling, Teologerne
ikke kunde være ganske tilfredse med. Man
fristes til at anse en saadan Fordeling for ironisk ment,
og dog spores ingen Ironi i Spinozas Fremstilling.
Klart er det dog, at det ikke har været hans virkelige
Mening dermed at stille Tro og Viden ved Siden af
hinanden som lige stillede. Sandheden staar for ham
over Fromheden og regulerer denne, ikke omvendt.
% --- p. 160 ---
\setcounter{footnote}{0}%
\opage{160}Ikke blot bemærker han, at man for at bøje sig i
Troens Lydighed ikke maa vide, at det er usandt, hvad
man tror. Men han siger endog, at det er Filosofien
umuligt at indsé, at Menneskene frelses ved Lydighed,
fordi Lydigheden kun tager Hensyn til den bydendes
Vilje, ej til Sagens Nødvendighed og Sandhed. Fremdeles
kritiserer han i sin „teologisk-politiske Afhandling“
Miraklet fra et rent filosofisk Synspunkt: da Naturen
er ét med Gud, kan man ikke som i Mirakeltroen
appellere fra Naturens Kraft til Guds Kraft, og det
nytter ej at sige, at Miraklet er over, ej mod Naturen,
efter som Miraklet jo sker midt i Naturen og altsaa
afbryder dens Orden. Og endelig udtaler han ligefrem,
at Profeterne kun paa Grund af deres mangelfulde
Erkendelse have fremstillet Gud som Lovgiver og
Fyrste i Steden for som det evige og uendelige Grundvæsen.
De have gjort det for at lempe sig efter
Mængdens Fatteevne. De fleste gøre ikke det gode for
dets egen Skyld og af Kærlighed til det, men fordi
de frygte en almægtig Lovgiver og bevæges ved den
Løn og Straf, han har fastsat. Mængden trænger lige
ledes til Ceremonier og historiske Fortællinger og
Exempler, som kunne gøre Indtryk paa dens Fantasi.
Fornuftslutninger gøre intet Indtryk paa den. Den,
der bruger sin Fornuft, behøver derimod ikke den
historiske Religion. Loven er kun given for de ufri,
ikke for dem, der have „Aandens Frugter“ og besidde
den sande Erkendelse.

Der bliver altsaa en bestemt Modsætning tilbage
mellem de vidende og de troende. Det er to forskellige
Livsanskuelser, der staa over for hinanden- Og
% --- p. 161 ---
\setcounter{footnote}{0}%
\opage{161}Modsætningen er ret træffende angivet ved Spinozas
Antitese: Lydighed --- Sandhed. Hvad der adskiller
de to Retninger, er nemlig væsentlig den forskellige
Opfattelse af Avtoritetens Betydning: den positive Tro
fastholder, at der gives absolute Avtoriteter, den rationelle
Tro kender kun relative Avtoriteter. Spinoza
har dog ikke fremdraget det historiske Forhold mellem
de to Retninger, som abstrakt set ere aldeles uforsonlige
Modsætninger. Den historiske Virkelighed viser
os nemlig Avtoriterne som opdragende Magter i den
menneskelige Verden, og der finder en stadig Given
og Modtagen Sted mellem de to Poler; de bestemme
i Historiens Løb stedse hinanden. Sammenligner man
Religionernes og i det hele Avtoriteternes Række i
Historien, vil man finde, at de højere staaende stedse
have en mere rationel Karakter. Paa den anden Side
er Avtoriteten --- hvad Spinoza selv har vist i sin Statslære
--- en nødvendig Støtte for den sig udviklende
Frihed. Lige som Universet efter nogle Astronomers
Mening viser os Stjernegrupper paa meget forskellige
Trin af Udvikling fra Kaos til System, saaledes se vi
til alle Tider i den menneskelige Verden mange forskellige
Stadier af Udvikling fra Avtoritet til Frihed.
Og Spinoza har ved sit eget Exempel godtgjort, at
Overgang fra det ene Standpunkt til det andet ikke
medfører Opgivelse af den sande Religiøsitet. Den
religiøse Trang kan finde sin Tilfredsstillelse under
flere Former, end den positive Religions Dogmatisme
indrømmer. Spinozas inderste Tankegang var af religiøs
Karakter, naar man ved Religiøsitet forstaar Hen
% --- p. 162 ---
\setcounter{footnote}{0}%
\opage{162}givelse til den højere Sammenhæng, hvori den enkelte
er et Led. En Følelse af denne Sammenhæng er Livsaaren
i hans System. ---

Saa systematisk en Aand Spinoza var, maatte
Menneskene dog efter hans Overbevisning søge Tilfredsstillelse
for deres religiøse Trang hver paa sin
Maade. Menneskenes aandelige Natur \textit{(ingenium)} er
saa forskellig, at de ikke kunne finde Hvile ved de
samme Meninger. Samfundsmagten vil derfor kun gøre
Skade og arbejde sig selv imod, dersom den vil krænke
Tænke- og Talefriheden. Enhver er Herre over sin
Overbevisning og maa have Ret til at udtale den, naar
han ikke gør det paa en Maade, der kan stifte Had og
Oprør i Staten. Fromhed og Ugudelighed angaa jo,
som vi have hørt, ikke Meningerne i sig selv, men de
Handlinger, der udspringe af Meningerne. Dersom
Staten vilde gribe ind paa den frie Overbevisnings
Omraade, vilde den netop møde Modstand hos de ædleste
og bedste Mennesker. De, der kun tænke paa at
samle Penge og pleje deres Bug, vilde naturligvis ej
bryde sig om et saadant Indgreb. Men hvilken større
Ulykke kan der tænkes for en Stat end at hæderlige
Mænd sendes i Landflygtighed som Forbrydere, fordi
de nære ejendommelige Anskuelser og ikke forstaa at
hykle? Hvad er skrækkeligere, end at Mennesker føres
til Døden, ikke paa Grund af nogen Forbrydelse eller
Misgjerning, men paa Grund af deres frie Aandsretning?
Friheden kan ganske vist medføre Ulemper;
men der er saa mange Ulemper, man maa finde sig i,
for at opnaa des større Goder. --- Hvem skulde i hvert
% --- p. 163 ---
\setcounter{footnote}{0}%
\opage{163}Tilfælde have Ret til at gribe ind i den enkeltes Overbevisning?
Man kan ikke her beraabe sig paa nogen
guddommelig Myndighed. Gud kan kun paa usand
Maade fremstilles som Fyrste og Lovgiver; den guddommelige
Lov er ét med den evige Naturs Nødvendighed.
Historien, særlig Jødernes Historie, viser
noksom, at det fører til betænkelige Konsekventser,
naar Gejstligheden faar politisk Magt. Staten maa
derfor have Ret til at ordne den ydre Gudsdyrkelse,
som det stemmer med dens Tarv. Men det er dog
den bedste Ordning, at de forskellige religiøse Partier
faa Frihed til at bestaa ved Siden af hverandre. Spinoza
anfører Amsterdam som et Exempel paa, at en
saadan Ordning kan gennemføres uden uheldige Følger.
Intet er sikrere for Staten, end at Fromhed og Religiøsitet
kun sættes i Kærlighed og Retfærdighed, og at
Øvrighedens Ret, baade hvad helligt og verdsligt angaar,
indskrænkes til Handlingerne alene, medens enhver
i øvrigt faar Lov til at mene, hvad han vil, og
sige, hvad han mener.

Spinoza kunde uden Tvivl være gaaet endnu
videre og have begrænset Samfundsmagtens Ret til
Indgriben ogsaa paa Handlingernes Omraade endnu
mere. I denne Retning have Kant og Stuart Mill
fortsat hans Tankegang. Men det er allerede stort,
at han midt under det 17de Aarhundredes Dogmatisme
og Avtoritetsherredømme har udtalt Tænke- og
Talefrihedens Princip. Han godtgør her, med hvor liden
Ret man saa ofte har frakendt ham Blik for Personlighedens
Betydning. Som hans religiøse og filosofiske
% --- p. 164 ---
\setcounter{footnote}{0}%
\opage{164}Overbevisning var i Harmoni med hans klare, stille
og dog energiske Personlighed, saaledes viser han
ogsaa en dyb Forstaaelse af den personlige Overbevisnings
Betydning for Menneskelivets fulde og rige
Udvikling.

% --- p. 165 ---
\setcounter{footnote}{0}%
\opage{165}\begin{center}\emph{Liste}\end{center}


over de vigtigste filosofiske Udtryk og de Steder i Bogen,

\begin{center}hvor deres Betydning kan ses.\end{center}


Absolut (smlgn. relativ). 69---71.97 f.
Antropomorfisme (Opfattelse af det guddommelige under mennelige Former). 118. % PRINTED AS IS: „mennelige“
At isme. 54---58.
Attribut (Grundform). 70 f. 73 ff. 80 f. 89 f.
Deduktion, deduktiv (smlgn. Induktion). 6. 14 f. 44---46. 68. 95.
Deisme, Deist. 54---59.
Determinisme 120 ff. (Nødvendighedslære), deterministisk. 26. 49. 51---54.
Diskursiv (smlgn. intuitiv). 110 ff.
Dogmatisme, dogmatisk. 4---6. 118.
Dualisme (Tohedslære), dualistisk (smlgn. Monisme). 7---8. 59.

94. 107.
Empirisme, empirisk. Se Induktion.
Immanents (Iboen), immanent. 48. 83. 141.
Induktion (Erfaringsmetode), induktiv. 6. 12---15. 44---46. 76---79.

106 f. 110.
Intuition (anskuende Tænken), intuitiv. 110---115.
Kritisk Filosofi (smlgn. Dogmatisme). 4---6.
Modus (Enkeltvæsen). 73. 84 ff.
Monade. 67. 71. 93 f.
Monisme (Enhedslære), monistisk (smlgn. Dualisme). 8. 59.
Optimisme. 31. 134 ff.
Panteisme, Panteist. 8---10. 53. 57---59. 82. 140---141
Relation (Forhold), Relativitet (Betingethed), relativ. 59. 70 f.

81. 88. 118 f. 130. 143 f.
Spekulation, spekulativ. 2. 13. 16. 86 f. 97 f.
Substants (Grundvæsen). 6---9. 17. 26. 67. 70 f. 72 ff. 76 ff.
Teisme. 53 kfr. 57.

% --- p. 166 ---
\setcounter{footnote}{0}%
\opage{166}Oversættelse af det franske Vers Side 56.

Da kom en lille, gusten og langnæset Jøde,

Fattig, men tilfreds, tankefuld indadvendt,

En fin, spidsfindig Aand, mindre læst end anerkendt.

Hyllet i Descartes’, sin Læremesters, Klæder,

Langsomt han hen for den højeste træder

Og hvisker sagte til ham: „Om Forladelse, jeg vil

Betro Dem unter uns, jeg tror ikke, De er til.“ H. T.

Anmærkning til Side 127.

Der burde her være tilføjet en Note om, at Spinoza sikkert
ikke har Ret i umiddelbart at anvende Ideassociationens Love
paa Følelserne. Det fortjener i hvert Tilfælde nøjere psykologisk
Undersøgelse, om den ene Følelse umiddelbart kan fremkalde
den anden, eller om ikke snarere Vejen fra den ene Følelse til
den anden gaar gennem de til Følelsen knyttede Tanker og
Forestillinger og de Ideassociationer, som disse fremkalde.
Enhver Ting kan altsaa indirekte blive Aarsag til Glæde, Sorg
og Attraa, ikke fordi disse Følelser umiddelbart gaa over i
hinanden, men fordi der til Tingen er knyttet forskellige Forestillinger,
af hvilke den ene vækker Glæde, den anden Sorg
o. s. v. En Erindring vil f. Ex. vække Glæde, naar Forestillingen
om den tidligere Lykke er fremherskende, men Sorg,
naar hin Forestilling afløses af Forestillingen om, at denne
Lykke nu er forbi.

\begin{center}\emph{Rettelser.}\end{center}


\begin{center}S. 8, L. 5 f. n. for l. fra.\end{center}


S. 10, L. 3 f. n. hallucition l. hallucination.
S. 132, L. 11 f. n. gneerositas l. generositas.

\end{document}
